\documentclass[11pt,a4paper]{article}
\usepackage[a4paper,margin=1in]{geometry}
\usepackage{lmodern}
\usepackage[T1]{fontenc}
\usepackage{amsmath,amssymb,amsthm,mathtools}
\usepackage{booktabs,tabularx}
\usepackage{xcolor}
\usepackage{microtype}
\usepackage{hyperref}

\definecolor{linkblue}{RGB}{18,52,110}
\definecolor{citegreen}{RGB}{20,108,52}

\hypersetup{
  colorlinks=true,
  linkcolor=linkblue,
  citecolor=citegreen,
  urlcolor=linkblue,
  pdftitle={Finite-Temperature Bounds on Operator Growth and Quantum Chaos},
  pdfauthor={Prince Anand}
}

\pagestyle{plain}
\numberwithin{equation}{section}
\emergencystretch=2em
\setlength{\parskip}{3pt plus 1pt minus 1pt}

\theoremstyle{plain}
\newtheorem{theorem}{Theorem}[section]
\newtheorem{lemma}[theorem]{Lemma}
\newtheorem{proposition}[theorem]{Proposition}
\newtheorem{corollary}[theorem]{Corollary}

\theoremstyle{definition}
\newtheorem{definition}[theorem]{Definition}

\newcommand{\Hop}{\mathcal{H}_{\mathrm{op}}}
\newcommand{\HS}{\mathrm{HS}}
\newcommand{\MSS}{\mathrm{MSS}}
\newcommand{\Tr}{\operatorname{Tr}}
\newcommand{\Var}{\operatorname{Var}}
\newcommand{\ad}{\operatorname{ad}}
\newcommand{\supp}{\operatorname{supp}}
\newcommand{\dist}{\operatorname{dist}}
\newcommand{\diam}{\operatorname{diam}}
\newcommand{\Range}{\operatorname{Range}}
\newcommand{\Ker}{\operatorname{Ker}}
\newcommand{\sech}{\operatorname{sech}}
\newcommand{\arctanh}{\operatorname{arctanh}}
\newcommand{\sgn}{\operatorname{sgn}}
\newcommand{\E}{\mathbb{E}}

\begin{document}

\begin{center}
  {\LARGE Finite-Temperature Bounds on Operator Growth and Quantum Chaos}\\[12pt]
  {\large Prince Anand}\\[4pt]
  {\normalsize November 2026}
\end{center}

\vspace{4pt}
\begin{quote}\small
\noindent\textbf{Abstract.}
We prove microscopic bounds relating the quantum Lyapunov exponent $\lambda_{L,T}$ of thermally regulated out-of-time-order correlators to the Wightman Lanczos growth rate $\alpha_T^{(W)}$ for $r$-local many-body Hamiltonians. Using a method-of-characteristics inequality on the Krylov chain and a Wightman-orthogonalized operator-size filtration, we derive non-asymptotic exponential tail bounds and reduce $\lambda_{L,T}\le 2\alpha_T^{(W)}$ at arbitrary $\beta\ge 0$ to a static equilibrium estimate on the Gibbs state. By Dyck-path comparison with the Meixner chain, we obtain $2\liminf_{n\to\infty} b_{n,T}^{(W)}/n \le 2\pi/(\beta+2\kappa) < 2\pi/\beta$ without smoothness hypotheses, together with holomorphy of the regulated four-point correlator on $|\operatorname{Im}z|<\frac{1}{2}\max(r_W,r_V)$. For arbitrary non-commuting $r$-local qubit Hamiltonians in any spatial dimension $d\ge 1$ at $\beta\in[0,\beta_{\mathrm{wc}}]$, a depolarized projective $2$-design cluster expansion yields a global quantum reverse-Poincar\'e inequality and, under uniform linear domination of the finite-volume Wightman Lanczos coefficients, proves $\lambda_{L,T}\le 2\alpha_T^{(W)}\le 2\pi/(\beta+2\kappa)$ for both Wightman-commutator and Maldacena--Shenker--Stanford regulated correlators.
\end{quote}

\vspace{6pt}

\section{Introduction}\label{sec:intro}

Understanding how local perturbations spread and delocalize in interacting quantum many-body systems at thermal equilibrium is a foundational problem in quantum statistical mechanics and mathematical physics. At zero temperature or in operator norms, the propagation of local observables is constrained by Lieb--Robinson bounds~\cite{Lieb1972,Hastings2006,Nachtergaele2006}, which establish an effective ballistic light cone for commutators $[W(t),V]$ of spatially separated operators. At finite inverse temperature $\beta = 1/T$, however, thermal fluctuations suppress high-energy transitions and impose a fundamentally different speed limit on quantum information scrambling. In their analysis of thermally regulated out-of-time-order correlators (OTOCs), Maldacena, Shenker, and Stanford~\cite{MSS2016} showed that under thermal clustering and complex-time analyticity assumptions, the quantum Lyapunov exponent $\lambda_{L,T}$ governing the exponential growth of four-point correlators satisfies
\begin{equation}\label{eq:intro-mss}
  \lambda_{L,T} \le \frac{2\pi}{\beta}
\end{equation}
in natural units $\hbar = k_B = 1$. The bound~\eqref{eq:intro-mss} is saturated by black holes in Einstein gravity and by the low-temperature Sachdev--Ye--Kitaev (SYK) model~\cite{MS2016}, whereas generic lattice spin systems and weakly coupled models exhibit $\lambda_{L,T} < 2\pi/\beta$.

A microscopic approach to operator spreading was subsequently proposed by Parker, Cao, Avdoshkin, Scaffidi, and Altman~\cite{Parker2019}. Given a Hamiltonian $H$, an initial Hermitian local operator $\mathcal{O}$, and a thermal inner product on the operator space $\Hop$, the Heisenberg evolution $\mathcal{O}(t) = e^{i\mathcal{L}t}\mathcal{O}$ generated by the Liouvillian $\mathcal{L} = [H,\cdot]$ is mapped by the Lanczos algorithm onto a tight-binding problem on a semi-infinite chain with hopping amplitudes $\{b_n\}_{n\ge 1}$. The Universal Operator Growth Hypothesis~\cite{Parker2019} posits that in non-integrable many-body systems, the Lanczos coefficients grow asymptotically linearly, $b_n \sim \alpha n$, at a maximal rate consistent with the geometry of interactions, and that the resulting Krylov complexity $K(t) = \sum_{n\ge 0} n|\varphi_n(t)|^2$ grows exponentially at rate $2\alpha$. At infinite temperature ($\beta=0$), Parker et al.~\cite{Parker2019} proved that Krylov complexity upper-bounds a broad family of ``$q$-complexities,'' including summed squared commutators for $r$-local Hamiltonians, thereby establishing $\lambda_{L,\infty} \le 2\alpha_\infty$.

At finite inverse temperature $\beta > 0$, the natural inner product compatible with the symmetric thermal regulation of~\cite{MSS2016} is the Wightman inner product
\begin{equation}\label{eq:intro-wightman}
  (A|B)_\beta^W := \Tr\!\bigl(\rho_\beta^{1/2} A^\dagger \rho_\beta^{1/2} B\bigr),
\end{equation}
where $\rho_\beta := e^{-\beta H}/Z(\beta)$ is the Gibbs state. In this inner product, the Kubo--Martin--Schwinger (KMS) relation enforces exponential decay of the Wightman spectral measure at rate at least $e^{-\beta|\omega|/2}$, which under regularity hypotheses on $\{b_{n,T}^{(W)}\}$ implies $\alpha_T^{(W)} \le \pi/\beta$~\cite{Parker2019,Avdoshkin2020}. Parker et al.~\cite{Parker2019} conjectured the finite-temperature chaos--operator growth bound
\begin{equation}\label{eq:intro-uogh}
  \lambda_{L,T} \le 2\alpha_T^{(W)} \le \frac{2\pi}{\beta}.
\end{equation}
Whenever $\lambda_{L,T} < 2\pi/\beta$, the left inequality of~\eqref{eq:intro-uogh} provides a microscopic tightening of~\eqref{eq:intro-mss} governed solely by the thermal two-point autocorrelation function; in the large-$q$ SYK model, it is saturated across the entire temperature range~\cite{Parker2019}. The status of~\eqref{eq:intro-uogh} at generic finite temperature has been investigated from multiple perspectives, including eigenstate thermalization~\cite{Murthy2019}, Euclidean operator growth~\cite{Avdoshkin2020}, operator-size distributions~\cite{Roberts2018,Qi2019}, all-to-all scrambling bounds~\cite{Yin2020}, deformed SYK models~\cite{Garcia2025}, hydrodynamic effective field theories~\cite{Demulder2026}, Krylov winding and holographic teleportation~\cite{Perugu2026}, and recent surveys~\cite{Baiguera2025}.

\subsection{Structural Obstructions at Finite Temperature}\label{subsec:obstructions}

Extending the infinite-temperature proof of~\cite{Parker2019} to $\beta > 0$ encounters three distinct mathematical obstructions, each of which must be addressed separately.

The first obstruction is operator-algebraic. At $\beta = 0$, the Hilbert--Schmidt inner product factorizes across lattice sites, so the Pauli operator-size superoperator $\mathcal{S}$ is simultaneously diagonal in the Pauli string basis and $(r-1)$-banded with respect to $\mathcal{L}$. At $\beta > 0$, the Gibbs weight $\rho_\beta^{1/2}$ couples distant sites. Consequently, the Wightman adjoint of a local commutator $\mathcal{L}_{V_x}(A) = [V_x, A]$ is given by $\mathcal{L}_{V_x}^{\dagger W}(A) = V_x(i\beta/2)A - A V_x(-i\beta/2)$, which is nonlocal in space and unbounded in norm as $\beta\to\infty$. In particular, the Hilbert--Schmidt size superoperator $\mathcal{S}$ fails to be banded in $(\cdot|\cdot)_\beta^W$, and the infinite-temperature $q$-complexity theorem of~\cite[App.~E]{Parker2019} does not apply directly.

The second obstruction is thermodynamic. As observed by Qi and Streicher~\cite{Qi2019}, the summed MSS-regulated four-point correlator equals the thermal-subtracted Hilbert--Schmidt operator size $\mathcal{S}_\beta(\mathcal{O}(t)) = \langle\mathcal{J}_\beta(\mathcal{O}(t)), \mathcal{S}\mathcal{J}_\beta(\mathcal{O}(t))\rangle_\HS - \langle\rho_\beta^{1/2},\mathcal{S}\rho_\beta^{1/2}\rangle_\HS$, where $\mathcal{J}_\beta(A) := \rho_\beta^{1/4} A \rho_\beta^{1/4}$. Because $\rho_\beta^{1/2}$ itself carries an extensive operator size $\langle\rho_\beta^{1/2},\mathcal{S}\rho_\beta^{1/2}\rangle_\HS = \Theta(N)$, any naive bound on $\langle\mathcal{J}_\beta(\mathcal{O}(t)),\mathcal{S}\mathcal{J}_\beta(\mathcal{O}(t))\rangle_\HS$ risks mixing the dynamical growth of $\mathcal{O}(t)$ with $O(N)$ background fluctuations of the thermal state. Controlling the cancellation of this extensive background uniformly over all real times $t\in\mathbb{R}$ is essential for extracting the Lyapunov exponent in the scrambling regime $t \sim \lambda_{L,T}^{-1}\ln N$.

The third obstruction is spectral and analytic. Two technical links in the UOGH framework have remained reliant on smoothness or continuum assumptions: first, passing from a linear upper bound $b_n \le \alpha n + \gamma$ to an exponential bound $K(t) \le C e^{2\alpha t}$ without assuming smooth asymptotic behavior of $\{b_n\}$ (an open technical step noted in~\cite[Sec.~V.1]{Parker2019}); and second, deriving the right-hand inequality $2\alpha_T^{(W)} \le 2\pi/\beta$ from exponential decay of the Wightman spectral measure without assuming that $\{b_{n,T}^{(W)}\}$ varies smoothly with $n$~\cite{Parker2019,Avdoshkin2020}.

\subsection{Setting and Notation}\label{subsec:intro-notation}

Before stating our main theorems, we fix the basic notation used throughout the paper. Let $\Lambda = \{1,\dots,N\}$ be a finite set of $N$ qubits with Hilbert space $\mathcal{H} = \bigotimes_{x\in\Lambda}\mathbb{C}^2$ ($\dim\mathcal{H}=2^N$), or $N$ Majorana fermions $\{\gamma_i\}_{i=1}^N$ satisfying $\{\gamma_i,\gamma_j\}=\delta_{ij}$ ($\dim\mathcal{H}=2^{N/2}$ for even $N$). Let $\Hop := \mathcal{B}(\mathcal{H})$ denote the space of linear operators on $\mathcal{H}$, equipped with the Hilbert--Schmidt inner product $\langle A,B\rangle_\HS := \Tr(A^\dagger B)$ and norm $\|A\|_\HS := \langle A,A\rangle_\HS^{1/2}$.

Let $H = \sum_{X\in\mathcal{E}} h_X$ be a Hermitian $r$-local Hamiltonian ($r\ge 1$), meaning that $\mathcal{E}$ is a finite collection of nonempty subsets $X\subseteq\Lambda$ of cardinality $|X|\le r$ and each $h_X = h_X^\dagger$ acts nontrivially only on $X$. In the Majorana setting, we assume throughout that every $h_X$ has even fermion parity so that $H$ and $\rho_\beta$ are bosonic. We measure the local interaction strength by the per-site norm sum
\begin{equation}\label{eq:intro-Jsite}
  J_{\mathrm{site}} := \sup_{x\in\Lambda} \sum_{X\in\mathcal{E}:\, X\ni x} \|h_X\|_\infty \in (0,\infty),
\end{equation}
and for $r\ge 2$ we define the inverse locality scale
\begin{equation}\label{eq:intro-kappa}
  \kappa := \frac{1}{2(r-1)J_{\mathrm{site}}} > 0.
\end{equation}
For $0\le s\le N$, let $\mathcal{V}_s \subset \Hop$ denote the linear span of all Pauli (or Majorana) strings of weight at most $s$, and let $\Pi_s^{(\beta)}:\Hop\to\mathcal{V}_s$ be the $(\cdot|\cdot)_\beta^W$-orthogonal projector onto $\mathcal{V}_s$. We define the \emph{Wightman-orthogonalized operator-size superoperator} by
\begin{equation}\label{eq:intro-Qbeta}
  \mathcal{Q}^{(\beta)} := \sum_{s=0}^{N-1}\bigl(\mathbf{1}-\Pi_s^{(\beta)}\bigr).
\end{equation}
For each qubit site $x\in\Lambda$, let $\mathcal{V}_{(x^c)}\subset\Hop$ denote the subalgebra of operators acting trivially on site $x$, let $P_{x,0}^{(\beta)}:\Hop\to\mathcal{V}_{(x^c)}$ be the $(\cdot|\cdot)_\beta^W$-orthogonal projector onto $\mathcal{V}_{(x^c)}$, and define the \emph{quantum Glauber generator} in Wightman space by
\begin{equation}\label{eq:intro-Qtilde}
  \widetilde{\mathcal{Q}}^{(\beta)} := \sum_{x=1}^N \bigl(\mathbf{1} - P_{x,0}^{(\beta)}\bigr).
\end{equation}
We also denote by $\mathcal{E}_x(A) := \frac{1}{2}\Tr_x(A)\otimes I_x$ the normalized single-site partial trace and by $\mathcal{S}_x := \mathbf{1} - \mathcal{E}_x$ the Hilbert--Schmidt orthogonal projector onto operators traceless at site $x$.

Two closely related finite-temperature OTOCs will be studied side by side. The first is the summed Wightman squared-commutator of~\cite{Parker2019},
\begin{equation}\label{eq:intro-CW}
  C_W(A) := \frac{1}{8}\sum_{x=1}^N\sum_{\alpha=1}^3 \bigl\|[\sigma_x^\alpha, A]\bigr\|_{\beta,W}^2.
\end{equation}
The second is the thermal-subtracted operator size $\mathcal{S}_\beta(A) := \langle\mathcal{J}_\beta(A),\mathcal{S}\mathcal{J}_\beta(A)\rangle_\HS - \langle\rho_\beta^{1/2},\mathcal{S}\rho_\beta^{1/2}\rangle_\HS$, which for every Hermitian $A$ with $\|A\|_{\beta,W}=1$ on $N$ qubits equals the summed MSS-regulated four-point correlator (Proposition~\ref{prop:mss-size}):
\begin{equation}\label{eq:intro-Sbeta}
  \mathcal{S}_\beta(A) = \frac{1}{4}\sum_{x=1}^N\sum_{\alpha=1}^3 \bigl(C_{\beta,\sigma_x^\alpha}^W(0) - F_\MSS(A,\sigma_x^\alpha)\bigr),
\end{equation}
where $F_\MSS(A,V) := \Tr(\rho_\beta^{1/4}A\rho_\beta^{1/4}V\rho_\beta^{1/4}A\rho_\beta^{1/4}V)$ and $C_{\beta,V}^W(t) := \Tr(\rho_\beta^{1/2}V(t)\rho_\beta^{1/2}V)$. As we show in Section~\ref{sec:entropic}, the dynamical leading term of $\mathcal{S}_\beta(\mathcal{O}(t))$ is governed by the imaginary-time shifted commutator sum
\begin{equation}\label{eq:intro-I}
  \mathrm{I}(A) := \frac{1}{8}\sum_{x=1}^N\sum_{\alpha=1}^3 \bigl\|[\sigma_x^\alpha(-i\beta/4), A]\bigr\|_{\beta,W}^2,
\end{equation}
where $V(z) := e^{izH} V e^{-izH}$. Finally, we define the static thermal size-filtration function
\begin{equation}\label{eq:intro-Lambda}
  \Lambda_\beta^{\MSS}(s) := \sup_{\substack{A\in\mathcal{V}_s \\ \|A\|_{\beta,W}=1}} \mathcal{S}_\beta(A) \qquad (s\in\{0,1,\dots,N\}),
\end{equation}
noting that since $\mathcal{J}_\beta(A^\dagger) = \mathcal{J}_\beta(A)^\dagger$ and $\mathcal{S}(M^\dagger) = \mathcal{S}(M)^\dagger$, the quadratic form $\mathcal{S}_\beta(A)$ preserves the real subspace of Hermitian operators, so the supremum over all $A\in\mathcal{V}_s$ coincides with the supremum over Hermitian $A\in\mathcal{V}_s$.

\subsection{Main Results}\label{subsec:main-theorems}

Our first main theorem supplies non-asymptotic exponential tail bounds for general linearly bounded Krylov chains and reduces the left inequality of~\eqref{eq:intro-uogh} at arbitrary inverse temperature $\beta\ge 0$ to a static sub-linear estimate on $\Lambda_\beta^{\MSS}(s)$.

\begin{theorem}\label{thm:main-tails-reduction}
\begin{enumerate}
  \setlength{\itemsep}{3pt}\setlength{\parskip}{0pt}
  \item[\emph{(a)}] Let $\{b_n\}_{n\ge 1}$ be any sequence of nonnegative real numbers satisfying $b_n \le \alpha n + \gamma$ for all $n\ge 1$ with $\alpha > 0$ and $\gamma \ge 0$, and let $\{\varphi_n(t)\}_{n\ge 0}$ be the $\ell^2$-solution to the Krylov evolution equation $\partial_t\varphi_n = b_n\varphi_{n-1} - b_{n+1}\varphi_{n+1}$ ($b_0=0$, $\varphi_{-1}=0$) with initial data $\varphi_n(0)=\delta_{n0}$. Then for every $n\ge 0$, every $t\ge 0$, and every $\theta_0 > 0$,
  \begin{equation}\label{eq:main-tail}
    \sum_{m=n}^\infty |\varphi_m(t)|^2 \le \exp\!\left(\left(\frac{1}{2}+\frac{\gamma}{\alpha}\right)\theta_0\right)\exp\!\left(-4\tanh\!\left(\frac{\theta_0}{4}\right)ne^{-2\alpha t}\right).
  \end{equation}
  In particular, the Krylov complexity $K(t) = \sum_{n=0}^\infty n|\varphi_n(t)|^2$ obeys
  \begin{equation}\label{eq:main-Kbound}
    K(t) \le \frac{\exp\!\bigl((\frac{1}{2}+\frac{\gamma}{\alpha})\theta_0\bigr)}{4e\tanh(\theta_0/4)}\,e^{2\alpha t} \qquad \forall t\ge 0,
  \end{equation}
  and hence $\limsup_{t\to\infty}\frac{1}{t}\ln K(t) \le 2\alpha$.
  \item[\emph{(b)}] Consider a family of $r$-local Hamiltonians $\{H_N\}$ on $N$ degrees of freedom at fixed $\beta\ge 0$, together with $s_0$-local normalized Hermitian operators $\mathcal{O}_N$ (with $r, s_0$ independent of $N$) whose Wightman Lanczos coefficients satisfy $b_{n,T}^{(W)} \le \alpha n + \gamma$ uniformly in $N$ ($\alpha = \alpha_T^{(W)} > 0$, $\gamma \ge 0$). Suppose that:
  \begin{itemize}
    \setlength{\itemsep}{2pt}\setlength{\parskip}{0pt}
    \item[\emph{(H1)}] There exist $\lambda_L > 0$, $c > 0$, and arbitrarily small $\epsilon > 0$ such that at $t_N(\epsilon) := \frac{1-\epsilon}{\lambda_L}\ln N$, $\mathcal{S}_\beta(\mathcal{O}_N(t_N(\epsilon))) \ge c N^{1-\epsilon}$ for all sufficiently large $N$.
    \item[\emph{(H2)}] There exist $\delta\in[0,1)$ and $N$-independent constants $C_1(\beta), C_2(\beta) > 0$ such that
    \begin{equation}\label{eq:main-H2}
      \Lambda_\beta^{\MSS}(s) \le C_1(\beta)\,s + C_2(\beta)\,N^\delta s^{1-\delta} \qquad \forall s\in\{1,\dots,N\}.
    \end{equation}
  \end{itemize}
  Then $\lambda_L \le 2\alpha_T^{(W)}$.
\end{enumerate}
\end{theorem}

The proof of Theorem~\ref{thm:main-tails-reduction}(a) uses a method-of-characteristics inequality for the moment generating function $G(\theta,t)=\sum_{n\ge 0} e^{\theta n}|\varphi_n(t)|^2$ on finite truncations, followed by Carleman's essential self-adjointness criterion and Fatou's lemma as the truncation length tends to infinity. Because $4\tanh(\theta_0/4)\to 4$ as $\theta_0\to\infty$, the bound~\eqref{eq:main-tail} matches the Meixner exponential decay rate $4ne^{-2\alpha t}$ without any smoothness assumption on $\{b_n\}$. Part~(b) then combines~\eqref{eq:main-tail} with a filtration truncation argument (Lemma~\ref{lem:truncation}) built from $\mathcal{Q}^{(\beta)}$, eliminating all real-time dynamics in favor of the static equilibrium bound~\eqref{eq:main-H2}. In Section~\ref{sec:filtration} we verify~\eqref{eq:main-H2} with $\delta=1/2$ for every full-rank product thermal state (Proposition~\ref{prop:product}).

Our second main theorem addresses the right-hand inequality of~\eqref{eq:intro-uogh} and the complex-time analyticity domain of the MSS-regulated four-point correlator. Let $r_W := \sup\{r > 0 : C_{\beta,W}^W(z)\text{ extends holomorphically to }|\operatorname{Im}z|<r\}$, and define $r_V$ analogously.

\begin{theorem}\label{thm:main-spectral}
Let $H = \sum_{X\in\mathcal{E}} h_X$ be an $r$-local Hamiltonian ($r\ge 2$) with per-site interaction strength $J_{\mathrm{site}} < \infty$ and $\kappa = (2(r-1)J_{\mathrm{site}})^{-1}$, and let $W,V$ be local Hermitian operators with $\|W\|_\infty,\|V\|_\infty < \infty$.
\begin{enumerate}
  \setlength{\itemsep}{3pt}\setlength{\parskip}{0pt}
  \item[\emph{(a)}] The nonnegative Wightman spectral measure $d\mu_{\beta,W}^W(\omega)$ satisfies
  \begin{equation}\label{eq:main-measure}
    \int_{\mathbb{R}} e^{(\frac{\beta}{2}+\kappa')|\omega|}\,d\mu_{\beta,W}^W(\omega) < \infty \qquad \forall\kappa'\in(0,\kappa).
  \end{equation}
  In particular, $r_W \ge \frac{\beta}{2}+\kappa$.
  \item[\emph{(b)}] For any thermodynamic-limit Wightman Lanczos sequence $\{b_{n,T}^{(W)}\}_{n\ge 1}$ with $b_{n,T}^{(W)} > 0$ for all $n\ge 1$ associated with a spectral measure obeying~\eqref{eq:main-measure},
  \begin{equation}\label{eq:main-right}
    2\liminf_{n\to\infty}\frac{b_{n,T}^{(W)}}{n} \le \frac{\pi}{r_W} \le \frac{2\pi}{\beta+2\kappa} < \frac{2\pi}{\beta}.
  \end{equation}
  In particular, whenever $\alpha_T^{(W)} := \lim_{n\to\infty} b_{n,T}^{(W)}/n$ exists, $2\alpha_T^{(W)} \le 2\pi/(\beta+2\kappa) < 2\pi/\beta$.
  \item[\emph{(c)}] The regulated four-point correlator $F_N(z) := \Tr(\rho_\beta^{1/4}W(z)\rho_\beta^{1/4}V\rho_\beta^{1/4}W(z)\rho_\beta^{1/4}V)$ satisfies, in every finite volume $N$ and for all $z = t + i\tau\in\mathbb{C}$,
  \begin{equation}\label{eq:main-otoc-bound}
    |F_N(t+i\tau)| \le \min\!\bigl(\|V\|_\infty^2 C_{\beta,W,N}^W(2i\tau),\; \|W\|_\infty^2 C_{\beta,V,N}^W(2i\tau),\; T_{\tau,N}(t)\bigr),
  \end{equation}
  where $T_{\tau,N}(t) := \Tr(\rho_\beta^{1/4}W(t-i\tau)\rho_\beta^{1/2}W(t+i\tau)\rho_\beta^{1/4}V^2)$. Every thermodynamic limit $F(t)$ of $F_N(t)$ along $\mathbb{R}$ extends holomorphically to the strip $\{z\in\mathbb{C} : |\operatorname{Im}z| < \frac{\beta}{4}+\frac{\kappa}{2}\}$ and obeys~\eqref{eq:main-otoc-bound} there (with the right-hand side replaced by its $\liminf_{N\to\infty}$, or by the limiting correlators whenever $C_{\beta,W,N}^W(2i\tau)$, $C_{\beta,V,N}^W(2i\tau)$, and $T_{\tau,N}(t)$ converge along the same subsequence). Moreover, if $C_{\beta,W,N}^W(2i\tau)$ (respectively, $C_{\beta,V,N}^W(2i\tau)$) is uniformly bounded in $N$ on compact subsets of $|\tau| < r_W/2$ (respectively, $|\tau| < r_V/2$), then $F(z)$ extends holomorphically to
  \begin{equation}\label{eq:main-strip}
    \mathcal{S}_{\mathrm{OTOC}} := \left\{z\in\mathbb{C} : |\operatorname{Im}z| < \frac{1}{2}\max(r_W,r_V)\right\} \supset \left\{z\in\mathbb{C} : |\operatorname{Im}z| < \frac{\beta}{4}+\frac{\kappa}{2}\right\}
  \end{equation}
  and satisfies~\eqref{eq:main-otoc-bound} throughout $\mathcal{S}_{\mathrm{OTOC}}$.
\end{enumerate}
\end{theorem}

Part~(b) of Theorem~\ref{thm:main-spectral} removes the smoothness assumption previously required to link $r_W$ to $\alpha_T^{(W)}$: expanding the $2n$-th moment as a sum over Dyck paths, positivity of all path weights allows a direct monotone comparison with the $\eta=1$ Meixner chain, whose moments are the Euler secant numbers.

Our third main theorem resolves~\eqref{eq:intro-uogh} for arbitrary non-commuting $r$-local qubit Hamiltonians in any spatial dimension $d\ge 1$ (or on any bounded-degree hypergraph) at high temperature, for both the Wightman-commutator OTOC $C_W$ and the MSS-regulated OTOC $\mathcal{S}_\beta$. Define the inverse-temperature threshold
\begin{equation}\label{eq:intro-betawc}
  \beta_{\mathrm{wc}} := \frac{1}{4e^2 r 6^r J_{\mathrm{site}}} < 2\kappa,
\end{equation}
which lies strictly inside the convergence radius $(8e^2 r J_{\mathrm{site}})^{-1}$ of our trace cluster expansion (Lemma~\ref{lem:kp-cumulant}).

\begin{theorem}\label{thm:main-highT}
Let $H = \sum_{X\in\mathcal{E}} h_X$ be any (not necessarily commuting) $r$-local Hamiltonian ($r\ge 2$) on $N$ qubits with per-site interaction strength $J_{\mathrm{site}} < \infty$. For every $\beta\in[0,\beta_{\mathrm{wc}}]$:
\begin{enumerate}
  \setlength{\itemsep}{3pt}\setlength{\parskip}{0pt}
  \item[\emph{(a)}] The global quantum reverse-Poincar\'e inequality holds on the entire operator Hilbert space $\Hop$:
  \begin{equation}\label{eq:main-rp}
    \bigl(A\big|\widetilde{\mathcal{Q}}^{(\beta)}\big|A\bigr)_\beta^W \le \sum_{x=1}^N \|\mathcal{S}_x(A)\|_{\beta,W}^2 \le 9\bigl(A\big|\mathcal{Q}^{(\beta)}\big|A\bigr)_\beta^W \qquad \forall A\in\Hop.
  \end{equation}
  \item[\emph{(b)}] Both the Wightman squared-commutator and the imaginary-time shifted MSS commutator are globally dominated by $\mathcal{Q}^{(\beta)}$ on $\Hop$:
  \begin{equation}\label{eq:main-CW-I}
    C_W(A) \le 54\bigl(A\big|\mathcal{Q}^{(\beta)}\big|A\bigr)_\beta^W, \qquad \mathrm{I}(A) \le 216\bigl(A\big|\mathcal{Q}^{(\beta)}\big|A\bigr)_\beta^W \qquad \forall A\in\Hop.
  \end{equation}
  Consequently, for any $s_0$-local Hermitian $\mathcal{O}$ with $\|\mathcal{O}\|_{\beta,W}=1$ and all $t\in\mathbb{R}$,
  \begin{equation}\label{eq:main-dyn-bounds}
    C_W(\mathcal{O}(t)) \le 54\bigl((r-1)K_\beta^W(t) + s_0\bigr), \qquad \mathrm{I}(\mathcal{O}(t)) \le 216\bigl((r-1)K_\beta^W(t) + s_0\bigr).
  \end{equation}
  \item[\emph{(c)}] Consider a family $\{H_N\}$ on $N$ qubits with uniform $(r,J_{\mathrm{site}})$ ($r, s_0$ independent of $N$), $s_0$-local Hermitian operators $\mathcal{O}_N$ obeying $\|\mathcal{O}_N\|_{\beta,W}=1$ and $\|\mathcal{O}_N\|_\infty = O(1)$, and Wightman Lanczos coefficients obeying $b_{n,T,N}^{(W)} \le \alpha_T^{(W)} n + \gamma$ uniformly in $N$, where $\alpha_T^{(W)} := \lim_{n\to\infty} b_{n,T,\infty}^{(W)}/n$ is the growth rate of the thermodynamic-limit sequence $b_{n,T,\infty}^{(W)} > 0$. Define $\lambda_{L,T}^{(W)}$ (respectively, $\lambda_{L,T}^{(\MSS)}$) as the supremum of all $\lambda_L > 0$ for which $C_W(\mathcal{O}_N(t_N(\epsilon))) \ge c N^{1-\epsilon}$ (respectively, $\mathcal{S}_\beta(\mathcal{O}_N(t_N(\epsilon))) \ge c N^{1-\epsilon}$) along $t_N(\epsilon) = \frac{1-\epsilon}{\lambda_L}\ln N$ for some $c>0$ and arbitrarily small $\epsilon > 0$. Then for every $\beta\in[0,\beta_{\mathrm{wc}}]$:
  \begin{itemize}
    \setlength{\itemsep}{2pt}\setlength{\parskip}{0pt}
    \item the Wightman-commutator Lyapunov exponent obeys $\lambda_{L,T}^{(W)} \le 2\alpha_T^{(W)} \le 2\pi/(\beta+2\kappa) < 2\pi/\beta$;
    \item the MSS-regulated thermal operator size satisfies
    \begin{equation}\label{eq:main-Sbeta-dyn}
      \mathcal{S}_\beta(\mathcal{O}_N(t)) \le 216\bigl((r-1)K_\beta^W(t)+s_0\bigr) + C_1(\beta)\sqrt{N\bigl((r-1)K_\beta^W(t)+s_0\bigr)} + C_2(\beta)\sqrt{N}
    \end{equation}
    for all sufficiently large $N$ and all $t\in\mathbb{R}$, and hence the MSS Lyapunov exponent obeys
    \begin{equation}\label{eq:main-mss-resolved}
      \lambda_{L,T}^{(\MSS)} \le 2\alpha_T^{(W)} \le \frac{2\pi}{\beta+2\kappa} < \frac{2\pi}{\beta}.
    \end{equation}
  \end{itemize}
\end{enumerate}
\end{theorem}

The key technical mechanism behind Theorem~\ref{thm:main-highT} is a depolarized on-site projective $2$-design decomposition of $\rho_{\beta/2}$. Recent works on high-temperature quantum Gibbs states~\cite{Bakshi2024,Kim2026} express $\rho_\beta$ as a convex mixture of tensor products of rank-$1$ projectors via a cluster expansion~\cite{Kotecky1986,Ueltschi2004}. Rank-$1$ projectors, however, are singular and do not yield a non-degenerate product inner product on $\Hop$. By depolarizing the six single-qubit Pauli eigenstate projectors to $\widetilde{P}_{s_x} := \frac{1}{2}P_{s_x} + \frac{1}{4}I_x$ (which satisfy $\frac{1}{4}I_x \preceq \widetilde{P}_{s_x} \preceq \frac{3}{4}I_x$, uniform condition number $3$), we obtain a strictly positive convex decomposition $\rho_{\beta/2} = \sum_{\mathbf{s}}\widetilde{\pi}(\mathbf{s})\widetilde{P}_{\mathbf{s}}$ for $\beta\le\beta_{\mathrm{wc}}$ whose induced Wightman norm is a convex average of well-conditioned product inner products. On each product component, the single-site projections and filtration projections are simultaneously diagonalized, and averaging the Hilbert-space projection theorem over the mixture produces~\eqref{eq:main-rp}. Expanding $\sigma_i^\alpha(-i\beta/4)$ into causal commutator sequences and reversing the root via Cayley's spanning-tree formula then yields the MSS-commutator bound $\mathrm{I}(A)\le 216(A|\mathcal{Q}^{(\beta)}|A)_\beta^W$ at the same temperature threshold $\beta_{\mathrm{wc}}$, which pairs with entropic background decoupling to prove~\eqref{eq:main-Sbeta-dyn} and~\eqref{eq:main-mss-resolved}.

\subsection{Relation to Prior Literature and Organization of the Paper}\label{subsec:org}

Our results connect two active lines of research: the operator-growth and Krylov-complexity program initiated in~\cite{Parker2019,Avdoshkin2020,Qi2019,Roberts2018,Murthy2019} and the mathematical statistical mechanics of quantum lattice Gibbs states~\cite{Araki1969,Lieb1972,Kotecky1986,Ueltschi2004,Kliesch2014,Frohlich2015,Abanin2015,Arad2016,Kato2019,Kuwahara2020a,Kuwahara2020b,Anshu2021,Bluhm2025,Alhambra2023,Bakshi2024,Rouze2026,Kim2026}. Relative to~\cite{Murthy2019}, which analyzed~\eqref{eq:intro-uogh} within the eigenstate thermalization hypothesis and isolated unknown constraints on four-point matrix-element correlations, Theorem~\ref{thm:main-tails-reduction}(b) bypasses eigenstate thermalization altogether and isolates the static criterion~\eqref{eq:main-H2}. Relative to the hydrodynamic effective field theory analysis of~\cite{Demulder2026}, where shift-symmetric non-microscopic correlators exhibited $\lambda_K = \lambda_L/2$, Proposition~\ref{prop:nogo}(1) shows that microscopic $r$-locality of both the operator and the Hamiltonian is structurally necessary for~\eqref{eq:intro-uogh}.

The remainder of the paper is organized as follows. Section~\ref{sec:prelim} develops the operator-algebraic framework, establishes the nested-commutator cluster bound (Lemma~\ref{lem:nested}), and proves the Wightman--Hilbert--Schmidt isometry and thermal-size identity (Propositions~\ref{prop:isometry} and~\ref{prop:mss-size}). Section~\ref{sec:tails} proves Theorem~\ref{thm:main-tails-reduction}(a) via the method of characteristics and Carleman's criterion. Section~\ref{sec:filtration} constructs the Wightman filtration, proves the finite-temperature $q$-complexity bound (Proposition~\ref{prop:q-complexity}) and size-truncation reduction (Lemma~\ref{lem:truncation}), proves Theorem~\ref{thm:main-tails-reduction}(b), and verifies~\eqref{eq:main-H2} for product thermal states (Proposition~\ref{prop:product}). Section~\ref{sec:spectral} proves Theorem~\ref{thm:main-spectral}. Section~\ref{sec:entropic} proves the high-temperature quantum Glauber bound (Proposition~\ref{prop:glauber}), the self-contained trace cumulant and entropic background decoupling estimates (Lemma~\ref{lem:kp-cumulant} and Proposition~\ref{prop:entropic}), and the resulting Krylov--Lieb--Robinson bound (Corollary~\ref{cor:klr}). Section~\ref{sec:highT} develops the depolarized projective $2$-design cluster expansion in modular lemmas and proves Theorem~\ref{thm:main-highT}. Section~\ref{sec:ising} establishes all-temperature bounds for one-dimensional quantum spin chains (Proposition~\ref{prop:1D}) and arbitrary $2$-local classical Ising models (Proposition~\ref{prop:ising}). Section~\ref{sec:nogo} proves three structural necessity results on locality, thermal clustering, and initial-state spread (Proposition~\ref{prop:nogo}). Section~\ref{sec:numerics} presents finite-size numerical diagonalization benchmarks. Appendices~\ref{app:trees} and~\ref{app:doublewell} supply self-contained proofs of the rooted-tree polymer estimates and the double-well Jacobi perturbation bounds.

\section{Preliminaries and Operator-Algebraic Framework}\label{sec:prelim}

In this section we formalize the lattice and fermionic operator algebras, prove the nested-commutator cluster estimate that underpins our complex-time analyticity bounds, and record the basic identities linking Wightman norms, Hilbert--Schmidt norms, Krylov wavefunctions, and thermal operator size.

\subsection{Local Hamiltonians and Imaginary-Time Evolution}\label{subsec:local-ham}

Let $\Lambda = \{1,\dots,N\}$ be a finite set of sites. We work in one of two finite-dimensional settings:
\begin{enumerate}
  \setlength{\itemsep}{2pt}\setlength{\parskip}{0pt}
  \item[\emph{(i)}] \emph{Quantum spin-$1/2$ systems:} Each site $x\in\Lambda$ carries a local Hilbert space $\mathcal{H}_x = \mathbb{C}^2$, the global Hilbert space is $\mathcal{H} = \bigotimes_{x\in\Lambda}\mathcal{H}_x$ ($\dim\mathcal{H} = 2^N$), and $\Hop = \mathcal{B}(\mathcal{H}) \cong \bigotimes_{x\in\Lambda} M_2(\mathbb{C})$. The Pauli operators at site $x$ are denoted $\sigma_x^1 = X_x$, $\sigma_x^2 = Y_x$, $\sigma_x^3 = Z_x$, together with the identity $I_x$.
  \item[\emph{(ii)}] \emph{Majorana fermion systems:} For even $N$, $\mathcal{H}$ ($\dim\mathcal{H} = 2^{N/2}$) carries $N$ Hermitian Majorana generators $\{\gamma_i\}_{i=1}^N$ obeying the Clifford relations $\{\gamma_i,\gamma_j\} = \delta_{ij}\mathbf{1}$, so that the normalized probes $V_i := \sqrt{2}\gamma_i$ satisfy $V_i^\dagger = V_i$ and $V_i^2 = \mathbf{1}$.
\end{enumerate}
The operator space $\Hop$ is equipped with the Hilbert--Schmidt inner product $\langle A,B\rangle_\HS := \Tr(A^\dagger B)$ and the Schatten norms $\|A\|_\HS := (\Tr(A^\dagger A))^{1/2}$, $\|A\|_1 := \Tr\sqrt{A^\dagger A}$, and $\|A\|_\infty := \sup_{\|\psi\|=1}\|A\psi\|$.

A Hermitian Hamiltonian $H = \sum_{X\in\mathcal{E}} h_X$ is $r$-local ($r\ge 1$) if each $h_X = h_X^\dagger$ has support $\supp(h_X) \subseteq X$ with $|X| \le r$. In the Majorana setting, we assume that every $h_X$ is a polynomial of even degree in the $\{\gamma_i\}$, so that $H$ commutes with the global fermion parity operator and $[h_X, P] = 0$ whenever $X \cap \supp(P) = \emptyset$. The per-site interaction strength $J_{\mathrm{site}}$ and inverse locality scale $\kappa = (2(r-1)J_{\mathrm{site}})^{-1}$ ($r\ge 2$) are defined in~\eqref{eq:intro-Jsite}--\eqref{eq:intro-kappa}.

For a fixed inverse temperature $\beta\in[0,\infty)$, the Gibbs state $\rho_\beta = e^{-\beta H}/Z(\beta)$ is strictly positive definite, so~\eqref{eq:intro-wightman} defines a positive-definite inner product $(\cdot|\cdot)_\beta^W$ on $\Hop$ with norm $\|A\|_{\beta,W} := (A|A)_\beta^{W\,1/2}$. The Liouvillian $\mathcal{L}(A) := [H,A] = \ad_H(A)$ generates the complex-time Heisenberg evolution
\[
  A(z) := e^{izH} A e^{-izH} = e^{iz\mathcal{L}} A, \qquad z\in\mathbb{C}.
\]
In particular, along the imaginary time axis $z = \pm i\tau$ ($\tau\in\mathbb{R}$), we have the relations
\begin{equation}\label{eq:prelim-imag}
  A(i\tau) = e^{-\tau H} A e^{\tau H} = \rho_\beta^{\tau/\beta} A \rho_\beta^{-\tau/\beta}, \qquad
  A(-i\tau) = e^{\tau H} A e^{-\tau H} = \rho_\beta^{-\tau/\beta} A \rho_\beta^{\tau/\beta},
\end{equation}
and for any Hermitian $A = A^\dagger$ and $z = t + i\tau$,
\begin{equation}\label{eq:prelim-adj}
  \bigl(A(t+i\tau)\bigr)^\dagger = \bigl(e^{(it-\tau)H} A e^{-(it-\tau)H}\bigr)^\dagger = e^{(it+\tau)H} A e^{-(it+\tau)H} = A(t-i\tau).
\end{equation}
For a qubit system, we recall the single-site normalized partial trace $\mathcal{E}_x(A) := \frac{1}{2}\Tr_x(A)\otimes I_x$ and its complementary projector $\mathcal{S}_x := \mathbf{1} - \mathcal{E}_x$. The maps $\{\mathcal{E}_x\}_{x\in\Lambda}$ are pairwise commuting orthogonal projectors in $(\Hop,\langle\cdot,\cdot\rangle_\HS)$ onto the subalgebras $\mathcal{V}_{(x^c)}$ of operators acting trivially on site $x$. For any subset $U\subseteq\Lambda$, $\mathcal{E}_U := \prod_{x\in U}\mathcal{E}_x$ is the Hilbert--Schmidt orthogonal projector onto $\mathcal{V}_{(U^c)}$, and $\mathbf{1}-\mathcal{E}_U$ is the orthogonal projector onto operators with nontrivial support in $U$.

The following lemma controls the growth of nested commutators $\ad_H^m(O_U)$ in operator norm, yielding an elementary convergence radius for imaginary-time evolution in terms of $r$ and $J_{\mathrm{site}}$ alone (compare~\cite{Abanin2015,Arad2016}).

\begin{lemma}[Nested-Commutator Cluster Bound]\label{lem:nested}
Let $H = \sum_{X\in\mathcal{E}} h_X$ be an $r$-local Hamiltonian ($r\ge 2$) with per-site interaction strength $J_{\mathrm{site}} < \infty$ and $\kappa = (2(r-1)J_{\mathrm{site}})^{-1}$. Let $O_U\in\Hop$ be supported on $U\subseteq\Lambda$ with $|U| = s \ge 1$. Then for every integer $m\ge 0$:
\begin{enumerate}
  \setlength{\itemsep}{2pt}\setlength{\parskip}{0pt}
  \item[\emph{(1)}] Expanding $\ad_H^m(O_U) = \sum_{\Gamma:\,|\Gamma|=m} T_\Gamma$ over ordered causal sequences $\Gamma = (X_1,\dots,X_m)\in\mathcal{E}^m$ with $X_1\cap U\ne\emptyset$ and $X_k\cap(U\cup X_1\cup\dots\cup X_{k-1})\ne\emptyset$ ($2\le k\le m$), where $T_\Gamma := [h_{X_m},[\dots,[h_{X_1},O_U]\dots]]$, we have
  \begin{equation}\label{eq:causal-norm-sum}
    \sum_{\Gamma:\,|\Gamma|=m}\|T_\Gamma\|_\infty \le \|O_U\|_\infty\,\kappa^{-m}\prod_{j=0}^{m-1}\left(\frac{s}{r-1}+j\right),
  \end{equation}
  and in particular $\|\ad_H^m(O_U)\|_\infty \le \|O_U\|_\infty\,\kappa^{-m}\prod_{j=0}^{m-1}(\frac{s}{r-1}+j)$.
  \item[\emph{(2)}] For every $|\tau| < 4\kappa$ and every single-site operator $V_i$ with $\|V_i\|_\infty \le 1$,
  \begin{equation}\label{eq:nested-Mbeta}
    \|V_i(\pm i\tau/4)\|_\infty \le \left(1 - \frac{|\tau|}{4\kappa}\right)^{-1}.
  \end{equation}
  \item[\emph{(3)}] For every $|\tau| < 4\kappa$ and every $r$-local interaction term $h_X$ ($|X|\le r$),
  \begin{equation}\label{eq:nested-hX}
    \|h_X(\pm i\tau/4)\|_\infty \le \|h_X\|_\infty\left(1 - \frac{|\tau|}{4\kappa}\right)^{-2}.
  \end{equation}
\end{enumerate}
\end{lemma}

\begin{proof}
We prove~\eqref{eq:causal-norm-sum} by induction on $m\ge 0$. The base case $m=0$ is immediate, with the empty product equal to $1$. Assume~\eqref{eq:causal-norm-sum} holds for a given $m\ge 0$. For each causal sequence $\Gamma = (X_1,\dots,X_m)$ of length $m$, the operator $T_\Gamma$ is supported on $U_\Gamma := U\cup X_1\cup\dots\cup X_m$. Because each $X_k$ intersects $U\cup X_1\cup\dots\cup X_{k-1}$ in at least one site and has cardinality $|X_k|\le r$, adding $X_k$ introduces at most $r-1$ new sites, whence $|U_\Gamma| \le s + m(r-1)$. Applying $\ad_H$ yields
\[
  \ad_H^{m+1}(O_U) = \sum_{\Gamma:\,|\Gamma|=m} \sum_{\substack{X\in\mathcal{E} \\ X\cap U_\Gamma\ne\emptyset}} [h_X, T_\Gamma].
\]
Using $\|[h_X,T_\Gamma]\|_\infty \le 2\|h_X\|_\infty\|T_\Gamma\|_\infty$ and summing over all sites $x\in U_\Gamma$, we obtain
\begin{align*}
  \sum_{\Gamma':\,|\Gamma'|=m+1} \|T_{\Gamma'}\|_\infty
  &\le \sum_{\Gamma:\,|\Gamma|=m} 2\|T_\Gamma\|_\infty \sum_{x\in U_\Gamma} \sum_{X\ni x}\|h_X\|_\infty \\
  &\le 2\bigl(s + m(r-1)\bigr)J_{\mathrm{site}} \sum_{\Gamma:\,|\Gamma|=m}\|T_\Gamma\|_\infty.
\end{align*}
Since $2(s + m(r-1))J_{\mathrm{site}} = 2(r-1)J_{\mathrm{site}}(\frac{s}{r-1}+m) = \kappa^{-1}(\frac{s}{r-1}+m)$, the inductive step follows and establishes~\eqref{eq:causal-norm-sum}.

To deduce~\eqref{eq:nested-Mbeta}, note that for $s=1$ and $r\ge 2$ we have $\frac{1}{r-1}\le 1$, so $\prod_{j=0}^{m-1}(\frac{1}{r-1}+j) \le \prod_{j=0}^{m-1}(1+j) = m!$. Expanding $V_i(\pm i\tau/4) = \sum_{m=0}^\infty \frac{(\mp\tau/4)^m}{m!}\ad_H^m(V_i)$ in operator norm gives
\[
  \|V_i(\pm i\tau/4)\|_\infty \le \sum_{m=0}^\infty \frac{(|\tau|/4)^m}{m!}\,\kappa^{-m} m! = \sum_{m=0}^\infty \left(\frac{|\tau|}{4\kappa}\right)^m = \left(1-\frac{|\tau|}{4\kappa}\right)^{-1}
\]
for all $|\tau|<4\kappa$. Similarly, for $s\le r$ and $r\ge 2$ we have $\frac{s}{r-1}\le 2$, so $\prod_{j=0}^{m-1}(\frac{s}{r-1}+j) \le (m+1)!$, whence
\[
  \|h_X(\pm i\tau/4)\|_\infty \le \|h_X\|_\infty \sum_{m=0}^\infty (m+1)\left(\frac{|\tau|}{4\kappa}\right)^m = \|h_X\|_\infty\left(1-\frac{|\tau|}{4\kappa}\right)^{-2}. \qedhere
\]
\end{proof}

\subsection{Wightman--Hilbert--Schmidt Isometry and Krylov Dynamics}\label{subsec:isometry}

The Wightman inner product $(\cdot|\cdot)_\beta^W$ is unitarily equivalent to the Hilbert--Schmidt inner product via symmetric multiplication by $\rho_\beta^{1/4}$. This elementary observation allows us to transfer Hilbert--Schmidt operator-size identities directly into Wightman space while preserving the Liouvillian $\mathcal{L}$.

\begin{proposition}[Wightman--Hilbert--Schmidt Isometry]\label{prop:isometry}
Define the linear map $\mathcal{J}_\beta:\Hop\to\Hop$ by $\mathcal{J}_\beta(A) := \rho_\beta^{1/4} A \rho_\beta^{1/4}$. Then:
\begin{enumerate}
  \setlength{\itemsep}{2pt}\setlength{\parskip}{0pt}
  \item[\emph{(1)}] $\mathcal{J}_\beta$ is a unitary isometry from $(\Hop,(\cdot|\cdot)_\beta^W)$ onto $(\Hop,\langle\cdot,\cdot\rangle_\HS)$.
  \item[\emph{(2)}] $\mathcal{J}_\beta$ commutes with the Liouvillian $\mathcal{L}=[H,\cdot]$: $\mathcal{J}_\beta(\mathcal{L}A) = \mathcal{L}(\mathcal{J}_\beta A)$ for all $A\in\Hop$.
  \item[\emph{(3)}] For any Hermitian initial operator $\mathcal{O}$ normalized to $\|\mathcal{O}\|_{\beta,W}=1$, the Wightman Lanczos coefficients $\{b_{n,T}^{(W)}\}$ and Krylov complexity $K_\beta^W(t)$ of $\mathcal{O}$ coincide with the Hilbert--Schmidt Lanczos coefficients and Krylov complexity of $\widetilde{\mathcal{O}}_0 := \mathcal{J}_\beta(\mathcal{O})$ evolving under $\mathcal{L}$.
\end{enumerate}
\end{proposition}

\begin{proof}
For all $A,B\in\Hop$, cyclicity of the trace and Hermiticity of $\rho_\beta^{1/4}$ yield
\[
  \langle\mathcal{J}_\beta(A),\mathcal{J}_\beta(B)\rangle_\HS
  = \Tr\!\bigl(\rho_\beta^{1/4}A^\dagger\rho_\beta^{1/4}\rho_\beta^{1/4}B\rho_\beta^{1/4}\bigr)
  = \Tr\!\bigl(\rho_\beta^{1/2}A^\dagger\rho_\beta^{1/2}B\bigr)
  = (A|B)_\beta^W.
\]
Since $\rho_\beta^{1/4}>0$ is invertible, $\mathcal{J}_\beta$ is bijective and hence unitary. Because $[H,\rho_\beta^{1/4}]=0$, we have $\rho_\beta^{1/4}[H,A]\rho_\beta^{1/4} = [H,\rho_\beta^{1/4}A\rho_\beta^{1/4}]$ for all $A\in\Hop$. Assertion~(3) follows immediately from unitary equivalence of the pairs $(\mathcal{L},\mathcal{O})$ in $(\Hop,(\cdot|\cdot)_\beta^W)$ and $(\mathcal{L},\widetilde{\mathcal{O}}_0)$ in $(\Hop,\langle\cdot,\cdot\rangle_\HS)$.
\end{proof}

We now recall the Lanczos recursion in $(\Hop,(\cdot|\cdot)_\beta^W)$. Starting from a Hermitian operator $\mathcal{O}$ with $\|\mathcal{O}\|_{\beta,W}=1$, set $|\mathcal{O}_0)_\beta^W := |\mathcal{O})_\beta^W$, $b_{0,T}^{(W)} := 0$, and $|\mathcal{O}_{-1})_\beta^W := 0$. For $n\ge 1$, define inductively
\begin{equation}\label{eq:lanczos-rec}
  |A_n) := \mathcal{L}|\mathcal{O}_{n-1})_\beta^W - b_{n-1,T}^{(W)}|\mathcal{O}_{n-2})_\beta^W, \qquad
  b_{n,T}^{(W)} := \|A_n\|_{\beta,W}, \qquad
  |\mathcal{O}_n)_\beta^W := \bigl(b_{n,T}^{(W)}\bigr)^{-1}|A_n),
\end{equation}
continuing until $b_{K_{\max},T}^{(W)}=0$ (which occurs at a finite Krylov dimension $K_{\max}\le (\dim\mathcal{H})^2$ in finite volume). Let us verify that the diagonal matrix elements vanish: first, $\mathcal{L}$ is self-adjoint in $(\cdot|\cdot)_\beta^W$, since for any $A,B\in\Hop$, cyclicity of the trace and $[H,\rho_\beta^{1/2}]=0$ give
\[
  (A|\mathcal{L}B)_\beta^W = \Tr\!\bigl(\rho_\beta^{1/2} A^\dagger\rho_\beta^{1/2}(HB-BH)\bigr)
  = \Tr\!\bigl(\rho_\beta^{1/2}(A^\dagger H - H A^\dagger)\rho_\beta^{1/2}B\bigr)
  = (\mathcal{L}A|B)_\beta^W.
\]
Second, because $\mathcal{O}_0$ is Hermitian and $\mathcal{L}$ maps Hermitian operators to anti-Hermitian operators and vice versa, $i^n\mathcal{O}_n^{(\beta)}$ is Hermitian for every $n\ge 0$, i.e., $(\mathcal{O}_n^{(\beta)})^\dagger = (-1)^n \mathcal{O}_n^{(\beta)}$. Hence
\[
  (\mathcal{O}_n|\mathcal{L}|\mathcal{O}_n)_\beta^W
  = (-1)^n \Tr\!\bigl(\rho_\beta^{1/2}\mathcal{O}_n^{(\beta)}\rho_\beta^{1/2}(H\mathcal{O}_n^{(\beta)} - \mathcal{O}_n^{(\beta)}H)\bigr) = 0,
\]
as the two terms coincide by cyclicity of the trace and $[H,\rho_\beta^{1/2}]=0$. Expanding $\mathcal{O}(t) = e^{i\mathcal{L}t}\mathcal{O} = \sum_{n\ge 0} i^n\varphi_n^{(\beta)}(t)\mathcal{O}_n^{(\beta)}$, the amplitudes $\varphi_n^{(\beta)}(t)$ are real-valued, satisfy $\varphi_n^{(\beta)}(0)=\delta_{n0}$ and $\sum_{n\ge 0}|\varphi_n^{(\beta)}(t)|^2=1$, and obey the discrete Schr\"odinger equation
\begin{equation}\label{eq:krylov-ode}
  \partial_t\varphi_n^{(\beta)}(t) = b_{n,T}^{(W)}\varphi_{n-1}^{(\beta)}(t) - b_{n+1,T}^{(W)}\varphi_{n+1}^{(\beta)}(t) \qquad (n\ge 0).
\end{equation}
The Wightman Krylov complexity is $K_\beta^W(t) := \sum_{n\ge 0} n|\varphi_n^{(\beta)}(t)|^2$.

\subsection{Operator Size and the MSS-Regulated Four-Point Correlator}\label{subsec:size-mss}

In the $N$-qubit setting, let $\{P\}$ denote the orthogonal basis of $4^N$ Pauli strings, and let $|P|\in\{0,1,\dots,N\}$ be the number of non-identity tensor factors of $P$. The Hilbert--Schmidt operator-size superoperator $\mathcal{S}$ acts diagonally by $\mathcal{S}(P) = |P|P$ and admits the commutator representation
\begin{equation}\label{eq:size-comm-rep}
  \langle M,\mathcal{S}M\rangle_\HS = \frac{1}{8}\sum_{x=1}^N\sum_{\alpha=1}^3 \bigl\|[M,\sigma_x^\alpha]\bigr\|_\HS^2 \qquad \forall M\in\Hop.
\end{equation}
To verify~\eqref{eq:size-comm-rep}, write $M = \sum_P c_P P$. For a fixed site $x$ and Pauli $\sigma_x^\alpha$, the map $P\mapsto [P,\sigma_x^\alpha]$ sends each Pauli string either to $0$ (if $P$ commutes with $\sigma_x^\alpha$) or to $2P\sigma_x^\alpha$ (if $P$ anticommutes with $\sigma_x^\alpha$). Since $P\mapsto P\sigma_x^\alpha$ is bijective on Pauli strings up to phases, distinct Pauli strings remain orthogonal after taking the commutator with $\sigma_x^\alpha$. At each site $x\in\supp(P)$, $P$ commutes with one Pauli and anticommutes with two, contributing $\frac{1}{8}(2\cdot 4)\|P\|_\HS^2 = \|P\|_\HS^2$; at each site $x\notin\supp(P)$, all three commutators vanish. Summing over $x\in\Lambda$ gives~\eqref{eq:size-comm-rep}.

In the $N$-Majorana setting with normalized probes $V_i = \sqrt{2}\gamma_i$ ($V_i^2 = \mathbf{1}$), a Majorana string $Q = \gamma_{i_1}\cdots\gamma_{i_k}$ of odd degree $k$ commutes with $V_i$ when $i\in\{i_1,\dots,i_k\}$ and anticommutes with $V_i$ when $i\notin\{i_1,\dots,i_k\}$, whereas an even-degree string anticommutes with $V_i$ for $i\in\{i_1,\dots,i_k\}$ and commutes otherwise. Thus
\[
  \langle M,\mathcal{S}M\rangle_\HS =
  \begin{cases}
    \dfrac{1}{4}\displaystyle\sum_{i=1}^N \|\{M,V_i\}\|_\HS^2 & \text{if $M$ has odd fermion parity},\\[10pt]
    \dfrac{1}{4}\displaystyle\sum_{i=1}^N \|[M,V_i]\|_\HS^2 & \text{if $M$ has even fermion parity}.
  \end{cases}
\]
Applying these identities to $M = \mathcal{J}_\beta(A)$ and subtracting the thermal background $M = \rho_\beta^{1/2}$ yields the relation of Qi and Streicher~\cite{Qi2019} between thermal operator size and the MSS-regulated OTOC.

\begin{proposition}[MSS-OTOC and Thermal-Size Identity]\label{prop:mss-size}
Let $A\in\Hop$ be Hermitian with $\|A\|_{\beta,W}=1$, and define the thermal-subtracted operator size
\[
  \mathcal{S}_\beta(A) := \langle\mathcal{J}_\beta(A),\mathcal{S}\mathcal{J}_\beta(A)\rangle_\HS - \langle\rho_\beta^{1/2},\mathcal{S}\rho_\beta^{1/2}\rangle_\HS.
\]
\begin{enumerate}
  \setlength{\itemsep}{2pt}\setlength{\parskip}{0pt}
  \item[\emph{(1)}] In an $N$-qubit system,
  \begin{equation}\label{eq:mss-size-spin}
    \mathcal{S}_\beta(A) = \frac{1}{4}\sum_{x=1}^N\sum_{\alpha=1}^3 \bigl(C_{\beta,\sigma_x^\alpha}^W(0) - F_\MSS(A,\sigma_x^\alpha)\bigr).
  \end{equation}
  \item[\emph{(2)}] In an $N$-Majorana system with even-parity $H$, odd-parity Hermitian $A$, and $V_i=\sqrt{2}\gamma_i$,
  \begin{equation}\label{eq:mss-size-maj}
    \mathcal{S}_\beta(A) = \frac{1}{2}\sum_{i=1}^N \bigl(C_{\beta,V_i}^W(0) + F_\MSS(A,V_i)\bigr).
  \end{equation}
\end{enumerate}
\end{proposition}

\begin{proof}
For any Hermitian operator $M$ and Hermitian unitary $P$ ($P^2=\mathbf{1}$), expanding the trace gives
\[
  \|[M,P]\|_\HS^2 = \Tr\bigl((PM-MP)(MP-PM)\bigr) = 2\Tr(M^2) - 2\Tr(MPMP).
\]
Setting $M = \mathcal{J}_\beta(A)$ yields $\Tr(M^2) = \|A\|_{\beta,W}^2 = 1$ and $\Tr(MPMP) = F_\MSS(A,P)$. Setting $M = \rho_\beta^{1/2}$ yields $\Tr(M^2) = \Tr(\rho_\beta) = 1$ and $\Tr(MPMP) = C_{\beta,P}^W(0)$. Subtracting the two expressions cancels the constant $2\Tr(M^2)=2$, and summing over $x\in\Lambda$ and $\alpha\in\{1,2,3\}$ with the factor $1/8$ gives~\eqref{eq:mss-size-spin}.

In the Majorana case, even parity of $H$ ensures that $\rho_\beta^{1/4}$ and $\rho_\beta^{1/2}$ have even parity, so $\mathcal{J}_\beta(A)$ has odd parity. Expanding $\|\{M,V_i\}\|_\HS^2 = 2\Tr(M^2)+2\Tr(MV_iMV_i)$ at $M=\mathcal{J}_\beta(A)$ gives $2+2F_\MSS(A,V_i)$, while $\|[M,V_i]\|_\HS^2 = 2\Tr(M^2)-2\Tr(MV_iMV_i)$ at $M=\rho_\beta^{1/2}$ gives $2-2C_{\beta,V_i}^W(0)$. Subtracting and summing over $i=1,\dots,N$ with the factor $1/4$ yields~\eqref{eq:mss-size-maj}.
\end{proof}

\section{Non-Asymptotic Krylov Tail Estimates}\label{sec:tails}

Before addressing the finite-temperature operator-algebraic structure, we resolve the dynamical problem on the Krylov chain itself. In~\cite[Sec.~V.1]{Parker2019}, exponential growth $K(t) \sim e^{2\alpha t}$ was computed for the Meixner family $b_n = \alpha\sqrt{n(n-1+\eta)}$ and conjectured to hold for general linearly bounded sequences $b_n \le \alpha n + \gamma$ without continuum heuristics or smoothness assumptions. In this section we prove Theorem~\ref{thm:main-tails-reduction}(a).

To keep the argument transparent, we separate the finite-chain differential inequality (Lemma~\ref{lem:char-finite}) from the infinite-volume limit via Carleman's essential self-adjointness criterion and Fatou's lemma.

\begin{lemma}[Finite-Chain Moment Generating Inequality]\label{lem:char-finite}
Fix $N_{\max}\ge 1$, $\alpha > 0$, and $\gamma \ge 0$. Let $\{b_n\}_{n=1}^{N_{\max}}$ be nonnegative real numbers obeying $b_n \le \alpha n + \gamma$ for $1\le n\le N_{\max}$, with $b_0 = b_{N_{\max}+1} = 0$, and let $\{\varphi_n(t)\}_{n=0}^{N_{\max}}$ be the real solution to~\eqref{eq:krylov-ode} with $\varphi_n(0)=\delta_{n0}$. Define
\[
  G_{N_{\max}}(\theta,t) := \sum_{n=0}^{N_{\max}} e^{\theta n}\varphi_n(t)^2 \qquad (\theta > 0,\; t\ge 0).
\]
Then for every $\theta_0 > 0$ and $t\ge 0$, defining $\theta(t)\in(0,\theta_0]$ implicitly by
\begin{equation}\label{eq:char-curve}
  \tanh\!\left(\frac{\theta(t)}{4}\right) = \tanh\!\left(\frac{\theta_0}{4}\right) e^{-2\alpha t},
\end{equation}
we have
\begin{equation}\label{eq:mgf-finite}
  G_{N_{\max}}(\theta(t), t) \le \exp\!\left(\left(\frac{1}{2}+\frac{\gamma}{\alpha}\right)\bigl(\theta_0-\theta(t)\bigr)\right) \le \exp\!\left(\left(\frac{1}{2}+\frac{\gamma}{\alpha}\right)\theta_0\right).
\end{equation}
\end{lemma}

\begin{proof}
Because $G_{N_{\max}}(\theta,t)$ is a finite sum of smooth functions of $(\theta,t)$, we may differentiate term by term. Inserting~\eqref{eq:krylov-ode} and re-indexing the second sum (using $b_0 = b_{N_{\max}+1} = 0$) gives
\begin{align*}
  \frac{\partial G_{N_{\max}}}{\partial t}(\theta,t)
  &= 2\sum_{n=0}^{N_{\max}} e^{\theta n}\varphi_n\bigl(b_n\varphi_{n-1} - b_{n+1}\varphi_{n+1}\bigr) \\
  &= 2\sum_{n=1}^{N_{\max}} b_n \varphi_n\varphi_{n-1}\bigl(e^{\theta n} - e^{\theta(n-1)}\bigr).
\end{align*}
We factor the exponential difference symmetrically as
\[
  e^{\theta n} - e^{\theta(n-1)} = e^{\theta(n-\frac{1}{2})}\bigl(e^{\theta/2}-e^{-\theta/2}\bigr) = 2\sinh\!\left(\frac{\theta}{2}\right) e^{\theta n/2} e^{\theta(n-1)/2}.
\]
For $\theta > 0$, we have $\sinh(\theta/2) > 0$ and $b_n \ge 0$. Applying the elementary inequality $2uv \le u^2 + v^2$ with $u = e^{\theta n/2}\varphi_n$ and $v = e^{\theta(n-1)/2}\varphi_{n-1}$ yields
\[
  \frac{\partial G_{N_{\max}}}{\partial t}(\theta,t)
  \le 2\sinh\!\left(\frac{\theta}{2}\right)\sum_{n=1}^{N_{\max}} b_n\bigl(e^{\theta n}\varphi_n^2 + e^{\theta(n-1)}\varphi_{n-1}^2\bigr).
\]
Substituting $b_n \le \alpha n + \gamma$ and splitting into two sums, the first sum satisfies
\[
  \sum_{n=1}^{N_{\max}} (\alpha n + \gamma) e^{\theta n}\varphi_n^2
  = \alpha\sum_{n=0}^{N_{\max}} n e^{\theta n}\varphi_n^2 + \gamma\sum_{n=1}^{N_{\max}} e^{\theta n}\varphi_n^2
  \le \alpha\frac{\partial G_{N_{\max}}}{\partial\theta} + \gamma G_{N_{\max}}.
\]
For the second sum, re-indexing $m = n - 1$ gives
\[
  \sum_{n=1}^{N_{\max}} (\alpha n + \gamma) e^{\theta(n-1)}\varphi_{n-1}^2
  = \sum_{m=0}^{N_{\max}-1}\bigl(\alpha(m+1)+\gamma\bigr)e^{\theta m}\varphi_m^2
  \le \alpha\frac{\partial G_{N_{\max}}}{\partial\theta} + (\alpha+\gamma)G_{N_{\max}}.
\]
Adding the two estimates yields the first-order differential inequality
\begin{equation}\label{eq:pde-ineq}
  \frac{\partial G_{N_{\max}}}{\partial t}(\theta,t)
  \le 4\alpha\sinh\!\left(\frac{\theta}{2}\right)\frac{\partial G_{N_{\max}}}{\partial\theta}(\theta,t)
  + 2(\alpha+2\gamma)\sinh\!\left(\frac{\theta}{2}\right)G_{N_{\max}}(\theta,t).
\end{equation}

We now integrate~\eqref{eq:pde-ineq} along the characteristic curve $\theta(t)$ solving the ordinary differential equation
\begin{equation}\label{eq:char-ode}
  \frac{d\theta}{dt} = -4\alpha\sinh\!\left(\frac{\theta}{2}\right), \qquad \theta(0) = \theta_0 > 0.
\end{equation}
Since
\[
  \frac{d}{d\theta}\ln\tanh\!\left(\frac{\theta}{4}\right)
  = \frac{\frac{1}{4}\sech^2(\theta/4)}{\tanh(\theta/4)}
  = \frac{1}{4\sinh(\theta/4)\cosh(\theta/4)}
  = \frac{1}{2\sinh(\theta/2)},
\]
separation of variables in~\eqref{eq:char-ode} gives $\frac{d}{dt}\ln\tanh(\theta(t)/4) = -2\alpha$, which integrates to~\eqref{eq:char-curve}. Along this curve, the chain rule and~\eqref{eq:pde-ineq} yield
\begin{align*}
  \frac{d}{dt} G_{N_{\max}}(\theta(t),t)
  &= \frac{\partial G_{N_{\max}}}{\partial t}(\theta(t),t) + \frac{d\theta}{dt}\frac{\partial G_{N_{\max}}}{\partial\theta}(\theta(t),t) \\
  &\le 2(\alpha+2\gamma)\sinh\!\left(\frac{\theta(t)}{2}\right)G_{N_{\max}}(\theta(t),t)
  = -\left(\frac{1}{2}+\frac{\gamma}{\alpha}\right)\frac{d\theta}{dt}\,G_{N_{\max}}(\theta(t),t).
\end{align*}
By unitarity of the finite-chain evolution, $G_{N_{\max}}(\theta(t),t) \ge \sum_{n=0}^{N_{\max}}\varphi_n(t)^2 = 1 > 0$. Dividing by $G_{N_{\max}}(\theta(t),t)$ and integrating from $0$ to $t$ with initial value $G_{N_{\max}}(\theta_0,0) = 1$ gives
\[
  \ln G_{N_{\max}}(\theta(t),t) \le \left(\frac{1}{2}+\frac{\gamma}{\alpha}\right)\bigl(\theta_0 - \theta(t)\bigr),
\]
which exponentiates to~\eqref{eq:mgf-finite}.
\end{proof}

\begin{proof}[Proof of Theorem~\ref{thm:main-tails-reduction}(a)]
If $b_1 = 0$, then $\varphi_n(t) = \delta_{n0}$ for all $t\ge 0$, so~\eqref{eq:main-tail}--\eqref{eq:main-Kbound} hold trivially. If $b_k = 0$ for some $k\ge 2$, then the Krylov dynamics decouples into a finite chain of length $N_{\max} = k-1 \ge 1$, and Lemma~\ref{lem:char-finite} applies directly. Assume now that $b_n > 0$ for all $n\ge 1$. Under the unitary gauge transformation $\psi_n(t) := i^n\varphi_n(t)$, the system~\eqref{eq:krylov-ode} becomes $\partial_t\psi = iT\psi$, where $T$ is the symmetric tridiagonal Jacobi operator on $\ell^2(\mathbb{Z}_{\ge 0})$ defined on finitely supported vectors by $(T\psi)_n = b_n\psi_{n-1} + b_{n+1}\psi_{n+1}$, and $|\varphi_m(t)|^2 = |\psi_m(t)|^2$. Let $T_{N_{\max}}$ be the finite truncation of $T$ with $b_n = 0$ for $n > N_{\max}$. Because $b_n \le \alpha n + \gamma$, the series $\sum_{n=1}^\infty b_n^{-1}$ diverges, so $T$ satisfies Carleman's condition and is therefore essentially self-adjoint on the domain of finitely supported vectors (see, e.g.,~\cite{Bratteli1997,Chihara1978}).

Consequently, $T_{N_{\max}}$ converges to the closure $\overline{T}$ in the strong resolvent sense as $N_{\max}\to\infty$, which implies strong convergence of the unitary groups $e^{i T_{N_{\max}} t} \to e^{i\overline{T} t}$ for every $t\in\mathbb{R}$. In particular, $|\varphi_m^{(N_{\max})}(t)|^2 \to |\varphi_m(t)|^2$ pointwise for each $m\ge 0$. Applying Fatou's lemma to~\eqref{eq:mgf-finite} (whose right-hand side is independent of $N_{\max}$) yields
\begin{equation}\label{eq:mgf-inf}
  G(\theta(t),t) := \sum_{m=0}^\infty e^{\theta(t)m}|\varphi_m(t)|^2
  \le \liminf_{N_{\max}\to\infty} G_{N_{\max}}(\theta(t),t)
  \le \exp\!\left(\left(\frac{1}{2}+\frac{\gamma}{\alpha}\right)\theta_0\right).
\end{equation}
Because $\tanh(x) \le x$ for all $x\ge 0$, equation~\eqref{eq:char-curve} implies
\[
  \theta(t) \ge 4\tanh\!\left(\frac{\theta(t)}{4}\right) = 4\tanh\!\left(\frac{\theta_0}{4}\right)e^{-2\alpha t}.
\]
Applying Markov's inequality $\sum_{m=n}^\infty |\varphi_m(t)|^2 \le e^{-\theta(t)n}G(\theta(t),t)$ together with~\eqref{eq:mgf-inf} yields~\eqref{eq:main-tail}. Finally, for all $x\ge 0$ and $\theta > 0$ we have $x \le (e\theta)^{-1}e^{\theta x}$, whence
\[
  K(t) = \sum_{n=0}^\infty n|\varphi_n(t)|^2 \le \frac{1}{e\theta(t)} G(\theta(t),t)
  \le \frac{\exp\!\bigl((\frac{1}{2}+\frac{\gamma}{\alpha})\theta_0\bigr)}{4e\tanh(\theta_0/4)}\,e^{2\alpha t},
\]
establishing~\eqref{eq:main-Kbound}.
\end{proof}

\noindent\emph{Remark.}
For the Meixner chain $b_n = \alpha\sqrt{n(n-1+\eta)}$ ($\eta\ge 1$), the Krylov probability distribution is $|\varphi_n(t)|^2 = \frac{(\eta)_n}{n!}\sech(\alpha t)^{2\eta}\tanh(\alpha t)^{2n}$~\cite{Parker2019}. Since $\tanh(\alpha t)^{2n} = (1 - 2e^{-2\alpha t} + O(e^{-4\alpha t}))^{2n} = \exp(-4ne^{-2\alpha t}(1+o(1)))$, the tail decays with exponential factor $\exp(-4ne^{-2\alpha t})$. Because $4\tanh(\theta_0/4) \to 4$ as $\theta_0\to\infty$, the estimate~\eqref{eq:main-tail} matches the Meixner exponential constant to arbitrary precision without assuming smoothness or asymptotic equality of $\{b_n\}$.

\section{Wightman Filtration Dominance and Static Reduction}\label{sec:filtration}

Having established the dynamical tail estimate on the Krylov chain, we now address the first structural obstruction described in Subsection~\ref{subsec:obstructions}: constructing a finite-temperature analogue of the infinite-temperature $q$-complexity superoperator~\cite[App.~E]{Parker2019} that remains compatible with the Liouvillian $\mathcal{L}$ in the Wightman inner product $(\cdot|\cdot)_\beta^W$. We then combine this construction with Theorem~\ref{thm:main-tails-reduction}(a) to prove Theorem~\ref{thm:main-tails-reduction}(b), and verify the resulting static criterion~\eqref{eq:main-H2} for arbitrary full-rank product thermal states.

\subsection{The Wightman-Orthogonalized Size Filtration}\label{subsec:wightman-filtration}

For each $s\in\{0,1,\dots,N\}$, recall that $\mathcal{V}_s \subset \Hop$ denotes the linear span of all Pauli strings (in the $N$-qubit setting) or Majorana strings (in the $N$-Majorana setting) of weight at most $s$, giving an ascending flag of subspaces
\[
  \mathbb{C}\mathbf{1} = \mathcal{V}_0 \subset \mathcal{V}_1 \subset \dots \subset \mathcal{V}_N = \Hop.
\]
Let $\Pi_s^{(\beta)}:\Hop\to\mathcal{V}_s$ denote the orthogonal projector onto $\mathcal{V}_s$ with respect to $(\cdot|\cdot)_\beta^W$, with the convention $\Pi_{-1}^{(\beta)} := 0$. Because $\mathcal{V}_0 \subset \mathcal{V}_1 \subset \dots \subset \mathcal{V}_N$, the projectors $\{\Pi_s^{(\beta)}\}_{s=0}^N$ form a commuting family of $(\cdot|\cdot)_\beta^W$-self-adjoint projections obeying $\Pi_{s-1}^{(\beta)} \preceq \Pi_s^{(\beta)}$ in the Loewner order on $(\Hop,(\cdot|\cdot)_\beta^W)$. Setting $\Delta\Pi_s^{(\beta)} := \Pi_s^{(\beta)} - \Pi_{s-1}^{(\beta)}$ for $0\le s\le N$, the Wightman-orthogonalized operator-size superoperator~\eqref{eq:intro-Qbeta} admits the spectral resolution
\begin{equation}\label{eq:Qbeta-spectral}
  \mathcal{Q}^{(\beta)} = \sum_{s=0}^{N-1}\bigl(\mathbf{1}-\Pi_s^{(\beta)}\bigr)
  = \sum_{s=0}^N s\,\Delta\Pi_s^{(\beta)}.
\end{equation}
In particular, $\mathcal{Q}^{(\beta)}$ is positive semidefinite and self-adjoint in $(\cdot|\cdot)_\beta^W$, satisfies $\mathcal{Q}^{(\beta)} = \mathcal{S}$ at $\beta = 0$, and measures the mean filtration degree of an operator in the Wightman inner product:
\[
  \bigl(A\big|\mathcal{Q}^{(\beta)}\big|A\bigr)_\beta^W
  = \sum_{s=1}^N s\bigl\|\Delta\Pi_s^{(\beta)}A\bigr\|_{\beta,W}^2
  = \sum_{s=0}^{N-1}\bigl\|(\mathbf{1}-\Pi_s^{(\beta)})A\bigr\|_{\beta,W}^2.
\]
Crucially, although $\Pi_s^{(\beta)}$ depends nontrivially on $\rho_\beta$, the underlying subspaces $\mathcal{V}_s$ are purely algebraic and satisfy the locality inclusion $\mathcal{L}(\mathcal{V}_s)\subseteq\mathcal{V}_{\min(N,s+r-1)}$. This is all that is needed to bound $\mathcal{Q}^{(\beta)}$ by Krylov complexity at arbitrary $\beta\ge 0$.

\begin{proposition}[Finite-Temperature $q$-Complexity Bound]\label{prop:q-complexity}
Let $H = \sum_{X\in\mathcal{E}} h_X$ be an $r$-local Hamiltonian ($r\ge 1$) on $N$ qubits or $N$ Majorana fermions (with even-parity $h_X$ in the fermionic case), and let $\mathcal{O}\in\mathcal{V}_{s_0}$ be an $s_0$-local Hermitian operator normalized to $\|\mathcal{O}\|_{\beta,W}=1$. Then for every $\beta\ge 0$:
\begin{enumerate}
  \setlength{\itemsep}{2pt}\setlength{\parskip}{0pt}
  \item[\emph{(1)}] Every Wightman Krylov basis element satisfies $\mathcal{O}_n^{(\beta)}\in\mathcal{V}_{\min(N,\,s_0+(r-1)n)}$ for all $0\le n < K_{\max}$.
  \item[\emph{(2)}] For every $t\in\mathbb{R}$,
  \begin{equation}\label{eq:q-complexity-bound}
    \bigl(\mathcal{O}(t)\big|\mathcal{Q}^{(\beta)}\big|\mathcal{O}(t)\bigr)_\beta^W
    \le (r-1)\,K_\beta^W(t) + s_0.
  \end{equation}
\end{enumerate}
\end{proposition}

\begin{proof}
We first verify that $\mathcal{L}(\mathcal{V}_s) \subseteq \mathcal{V}_{\min(N,s+r-1)}$ for every $0\le s\le N$. Let $P$ be a Pauli string of weight $|P|\le s$ and $h_X$ an interaction term supported on $X$ with $|X|\le r$. If $X\cap\supp(P)=\emptyset$, then $[h_X,P]=0$. If $X\cap\supp(P)\ne\emptyset$, then $|X\cup\supp(P)|\le |P|+|X|-1\le s+r-1$, so $[h_X,P]$ is supported on $X\cup\supp(P)$ and therefore lies in $\mathcal{V}_{\min(N,s+r-1)}$. In the Majorana setting, write $h_X$ as a linear combination of even-degree Majorana strings $Q$ supported on $X$ ($|Q|\le r$). If $\supp(Q)\cap\supp(P)=\emptyset$, then even degree of $Q$ implies $[Q,P]=0$; if $\supp(Q)\cap\supp(P)\ne\emptyset$, then $QP$ and $PQ$ have support inside $\supp(Q)\cup\supp(P)$ of cardinality at most $s+r-1$, so again $[h_X,P]\in\mathcal{V}_{\min(N,s+r-1)}$. Summing over $X\in\mathcal{E}$ yields $\mathcal{L}(\mathcal{V}_s)\subseteq\mathcal{V}_{\min(N,s+r-1)}$.

Part~(1) follows immediately by induction on $n$ from the Lanczos recursion~\eqref{eq:lanczos-rec}. For $n=0$, $\mathcal{O}_0^{(\beta)} = \mathcal{O}\in\mathcal{V}_{s_0}$. Assuming $\mathcal{O}_m^{(\beta)}\in\mathcal{V}_{\min(N,s_0+(r-1)m)}$ for all $m\le n-1$, we have $\mathcal{L}\mathcal{O}_{n-1}^{(\beta)}\in\mathcal{V}_{\min(N,s_0+(r-1)n)}$ and $\mathcal{O}_{n-2}^{(\beta)}\in\mathcal{V}_{\min(N,s_0+(r-1)(n-2))}\subseteq\mathcal{V}_{\min(N,s_0+(r-1)n)}$, whence $\mathcal{O}_n^{(\beta)}\in\mathcal{V}_{\min(N,s_0+(r-1)n)}$.

To prove part~(2), let $\mathcal{K} := \operatorname{span}\{\mathcal{O}_0^{(\beta)},\dots,\mathcal{O}_{K_{\max}-1}^{(\beta)}\}$ denote the Krylov subspace. If $r=1$, then part~(1) gives $\mathcal{K}\subseteq\mathcal{V}_{s_0}$, so $(\mathbf{1}-\Pi_s^{(\beta)})\mathcal{O}(t)=0$ for all $s\ge s_0$. Using $0\preceq\mathbf{1}-\Pi_s^{(\beta)}\preceq\mathbf{1}$ for $0\le s < s_0$ gives $(\mathcal{O}(t)|\mathcal{Q}^{(\beta)}|\mathcal{O}(t))_\beta^W \le s_0 = (r-1)K_\beta^W(t)+s_0$.

Assume now that $r\ge 2$. For each $k\ge 0$, define the nested Krylov subspace $\mathcal{K}_k := \operatorname{span}\{\mathcal{O}_0^{(\beta)},\dots,\mathcal{O}_{\min(k,K_{\max}-1)}^{(\beta)}\}$ and let $P_{\mathcal{K}_k}$ be the $(\cdot|\cdot)_\beta^W$-orthogonal projector onto $\mathcal{K}_k$. By part~(1), $\mathcal{K}_k\subseteq\mathcal{V}_{\min(N,s_0+(r-1)k)}$, which implies the Loewner operator inequality $P_{\mathcal{K}_k}\preceq \Pi_{\min(N,s_0+(r-1)k)}^{(\beta)}$ on $(\Hop,(\cdot|\cdot)_\beta^W)$. By monotonicity of $\Pi_s^{(\beta)}$ in $s$, for every integer $s$ in the block $s_0+(r-1)k \le s < \min(N, s_0+(r-1)(k+1))$ we have
\[
  \mathbf{1}-\Pi_s^{(\beta)} \preceq \mathbf{1}-\Pi_{s_0+(r-1)k}^{(\beta)} \preceq \mathbf{1}-P_{\mathcal{K}_k}.
\]
If $s_0 = N$, then $\mathcal{Q}^{(\beta)}\preceq N\mathbf{1}$ gives~\eqref{eq:q-complexity-bound} trivially. Assume $s_0 < N$. Splitting the sum~\eqref{eq:Qbeta-spectral} into $0\le s < s_0$ (where $\mathbf{1}-\Pi_s^{(\beta)}\preceq\mathbf{1}$) and blocks of length at most $r-1$ for $s_0\le s\le N-1$, and setting $k_{\max} := \lceil (N-s_0)/(r-1)\rceil - 1 \ge 0$, yields the operator inequality on $\Hop$:
\[
  \mathcal{Q}^{(\beta)}
  = \sum_{s=0}^{s_0-1}\bigl(\mathbf{1}-\Pi_s^{(\beta)}\bigr)
    + \sum_{k=0}^{k_{\max}}\sum_{s=s_0+(r-1)k}^{\min(N,\,s_0+(r-1)(k+1))-1}\bigl(\mathbf{1}-\Pi_s^{(\beta)}\bigr)
  \preceq s_0\mathbf{1} + (r-1)\sum_{k=0}^{k_{\max}}\bigl(\mathbf{1}-P_{\mathcal{K}_k}\bigr).
\]
Because $|\mathcal{O}(t))_\beta^W = \sum_{n=0}^{K_{\max}-1} i^n\varphi_n^{(\beta)}(t)|\mathcal{O}_n^{(\beta)})_\beta^W \in \mathcal{K} = \mathcal{K}_{K_{\max}-1}$, we have $(\mathcal{O}(t)|\mathbf{1}-P_{\mathcal{K}_k}|\mathcal{O}(t))_\beta^W = 0$ for every $k\ge K_{\max}-1$. Evaluating both sides on $|\mathcal{O}(t))_\beta^W$ and using nonnegativity of each summand therefore gives
\[
  \sum_{k=0}^{k_{\max}}\bigl(\mathcal{O}(t)\big|\mathbf{1}-P_{\mathcal{K}_k}\big|\mathcal{O}(t)\bigr)_\beta^W
  \le \sum_{k=0}^{K_{\max}-2}\sum_{n=k+1}^{K_{\max}-1}\bigl|\varphi_n^{(\beta)}(t)\bigr|^2
  = \sum_{n=0}^{K_{\max}-1} n\bigl|\varphi_n^{(\beta)}(t)\bigr|^2
  = K_\beta^W(t),
\]
which yields~\eqref{eq:q-complexity-bound}.
\end{proof}

\subsection{Size-Truncation Reduction and Proof of Theorem~\ref{thm:main-tails-reduction}(b)}\label{subsec:static-reduction}

Proposition~\ref{prop:q-complexity}(1) shows that truncating the Krylov expansion of $\mathcal{O}(t)$ at depth $M$ produces an operator lying strictly inside the low-weight subspace $\mathcal{V}_{s_0+(r-1)M}$. The following lemma quantifies the error incurred in $\mathcal{S}_\beta(\mathcal{O}(t))$ by this truncation in terms of the Krylov tail probability.

\begin{lemma}[Thermal-Size Truncation Inequality]\label{lem:truncation}
Let $H$ be an $r$-local Hamiltonian ($r\ge 1$) on $N$ qubits or $N$ Majorana fermions, let $\mathcal{O}\in\mathcal{V}_{s_0}$ be Hermitian with $\|\mathcal{O}\|_{\beta,W}=1$, and fix an integer $M\ge 0$ and a time $t\in\mathbb{R}$. Set $s_M := \min(N, s_0+(r-1)M)$ and let
\[
  \varepsilon_M(t) := \left(\sum_{n=M+1}^{K_{\max}-1}\bigl|\varphi_n^{(\beta)}(t)\bigr|^2\right)^{1/2}
\]
denote the $\ell^2$-tail norm beyond Krylov index $M$. Then
\begin{equation}\label{eq:truncation}
  \mathcal{S}_\beta\bigl(\mathcal{O}(t)\bigr) \le \Lambda_\beta^{\MSS}(s_M) + 2N\,\varepsilon_M(t) + N\,\varepsilon_M(t)^2.
\end{equation}
\end{lemma}

\begin{proof}
Decompose $\mathcal{O}(t) = A_M + R_M$, where
\[
  A_M := \sum_{n=0}^{\min(M,K_{\max}-1)} i^n\varphi_n^{(\beta)}(t)\mathcal{O}_n^{(\beta)}, \qquad
  R_M := \sum_{n=M+1}^{K_{\max}-1} i^n\varphi_n^{(\beta)}(t)\mathcal{O}_n^{(\beta)}.
\]
Since $i^n\mathcal{O}_n^{(\beta)}$ is Hermitian and $\varphi_n^{(\beta)}(t)\in\mathbb{R}$, both $A_M$ and $R_M$ are Hermitian. By Proposition~\ref{prop:q-complexity}(1), $A_M\in\mathcal{V}_{s_M}$. Moreover, by orthonormality of the Krylov basis and Proposition~\ref{prop:isometry}(1),
\[
  \|\mathcal{J}_\beta(A_M)\|_\HS = \|A_M\|_{\beta,W} = \sqrt{1-\varepsilon_M(t)^2} \le 1, \qquad
  \|\mathcal{J}_\beta(R_M)\|_\HS = \|R_M\|_{\beta,W} = \varepsilon_M(t).
\]
Expanding $\langle\mathcal{J}_\beta(\mathcal{O}(t)),\mathcal{S}\mathcal{J}_\beta(\mathcal{O}(t))\rangle_\HS$ using $\mathcal{J}_\beta(\mathcal{O}(t)) = \mathcal{J}_\beta(A_M) + \mathcal{J}_\beta(R_M)$ gives
\begin{align*}
  \mathcal{S}_\beta(\mathcal{O}(t))
  &= \langle\mathcal{J}_\beta(A_M),\mathcal{S}\mathcal{J}_\beta(A_M)\rangle_\HS - \langle\rho_\beta^{1/2},\mathcal{S}\rho_\beta^{1/2}\rangle_\HS \\
  &\quad + 2\operatorname{Re}\langle\mathcal{J}_\beta(A_M),\mathcal{S}\mathcal{J}_\beta(R_M)\rangle_\HS
    + \langle\mathcal{J}_\beta(R_M),\mathcal{S}\mathcal{J}_\beta(R_M)\rangle_\HS.
\end{align*}
Since $0\preceq\mathcal{S}\preceq N\mathbf{1}$ in $(\Hop,\langle\cdot,\cdot\rangle_\HS)$, the Cauchy--Schwarz inequality bounds the cross term by $2N\|\mathcal{J}_\beta(A_M)\|_\HS\|\mathcal{J}_\beta(R_M)\|_\HS \le 2N\varepsilon_M(t)$ and the remainder term by $N\|\mathcal{J}_\beta(R_M)\|_\HS^2 = N\varepsilon_M(t)^2$.

For the leading difference, observe first that the identity operator $\mathbf{1}\in\mathcal{V}_0\subseteq\mathcal{V}_{s_M}$ is Hermitian with $\|\mathbf{1}\|_{\beta,W} = (\Tr\rho_\beta)^{1/2} = 1$ and $\mathcal{S}_\beta(\mathbf{1}) = 0$, whence $\Lambda_\beta^{\MSS}(s_M)\ge 0$. If $A_M = 0$, then the leading difference equals $-\langle\rho_\beta^{1/2},\mathcal{S}\rho_\beta^{1/2}\rangle_\HS \le 0 \le \Lambda_\beta^{\MSS}(s_M)$. If $A_M\ne 0$, set $\widehat{A}_M := A_M/\|A_M\|_{\beta,W} \in \mathcal{V}_{s_M}$; then $\|\widehat{A}_M\|_{\beta,W}=1$, and since $\langle\rho_\beta^{1/2},\mathcal{S}\rho_\beta^{1/2}\rangle_\HS \ge 0$,
\begin{align*}
  \langle\mathcal{J}_\beta(A_M),\mathcal{S}\mathcal{J}_\beta(A_M)\rangle_\HS - \langle\rho_\beta^{1/2},\mathcal{S}\rho_\beta^{1/2}\rangle_\HS
  &= \|A_M\|_{\beta,W}^2\,\mathcal{S}_\beta(\widehat{A}_M) - \varepsilon_M(t)^2\langle\rho_\beta^{1/2},\mathcal{S}\rho_\beta^{1/2}\rangle_\HS \\
  &\le \|A_M\|_{\beta,W}^2\,\Lambda_\beta^{\MSS}(s_M) \le \Lambda_\beta^{\MSS}(s_M).
\end{align*}
Combining these bounds establishes~\eqref{eq:truncation}.
\end{proof}

We now combine Theorem~\ref{thm:main-tails-reduction}(a) with Lemma~\ref{lem:truncation} to prove Theorem~\ref{thm:main-tails-reduction}(b).

\begin{proof}[Proof of Theorem~\ref{thm:main-tails-reduction}(b)]
Suppose for contradiction that $\lambda_L > 2\alpha$, where $\alpha := \alpha_T^{(W)} > 0$. By~(H1), we may choose $\epsilon\in(0,1)$ sufficiently small that
\begin{equation}\label{eq:eps-choice}
  \eta := \delta + (1-\delta)\frac{2\alpha(1-\epsilon)}{\lambda_L} < 1-\epsilon,
\end{equation}
which is possible because as $\epsilon\downarrow 0$ the left-hand side tends to $\delta + (1-\delta)(2\alpha/\lambda_L) < 1$ (since $\delta\in[0,1)$ and $2\alpha/\lambda_L < 1$). Let $t_N := \frac{1-\epsilon}{\lambda_L}\ln N$, so that
\[
  e^{2\alpha t_N} = N^{2\alpha(1-\epsilon)/\lambda_L}.
\]
Fix $\theta_0 = 4$ in Theorem~\ref{thm:main-tails-reduction}(a), and define the truncation depth
\[
  M_N := \left\lceil \frac{1}{\tanh(1)}\,e^{2\alpha t_N}\ln N \right\rceil.
\]
By~\eqref{eq:main-tail}, the $\ell^2$-tail at $t = t_N$ satisfies
\[
  \varepsilon_{M_N}(t_N)^2 = \sum_{m=M_N+1}^\infty\bigl|\varphi_m^{(\beta)}(t_N)\bigr|^2
  \le \exp\!\left(4\left(\frac{1}{2}+\frac{\gamma}{\alpha}\right)\right)\exp\!\bigl(-4\tanh(1)(M_N+1)e^{-2\alpha t_N}\bigr)
  \le C_0 N^{-4},
\]
where $C_0 := \exp(2 + 4\gamma/\alpha)$ is independent of $N$. Hence $\varepsilon_{M_N}(t_N) \le C_0^{1/2} N^{-2}$, so both error terms in Lemma~\ref{lem:truncation} decay at least as $O(N^{-1})$:
\[
  2N\,\varepsilon_{M_N}(t_N) + N\,\varepsilon_{M_N}(t_N)^2 \le 2C_0^{1/2} N^{-1} + C_0 N^{-3} = o(1) \qquad (N\to\infty).
\]
On the other hand, setting $s_{M_N} := \min(N, s_0+(r-1)M_N)$, we have $s_{M_N} = O(N^{2\alpha(1-\epsilon)/\lambda_L}\ln N)$. Substituting this into Hypothesis~(H2) gives
\begin{align*}
  \Lambda_\beta^{\MSS}(s_{M_N})
  &\le C_1(\beta)\,s_{M_N} + C_2(\beta)\,N^\delta s_{M_N}^{1-\delta} \\
  &= O\!\left(N^{2\alpha(1-\epsilon)/\lambda_L}\ln N + N^{\delta + (1-\delta)\frac{2\alpha(1-\epsilon)}{\lambda_L}}(\ln N)^{1-\delta}\right)
  = O\!\left(N^\eta \ln N\right).
\end{align*}
Applying Lemma~\ref{lem:truncation} at $t = t_N$ therefore yields
\[
  \mathcal{S}_\beta\bigl(\mathcal{O}_N(t_N)\bigr) \le O\!\left(N^\eta \ln N\right) + o(1) = o\!\left(N^{1-\epsilon}\right)
\]
because $\eta < 1-\epsilon$ by~\eqref{eq:eps-choice}. This contradicts Hypothesis~(H1), which asserts $\mathcal{S}_\beta(\mathcal{O}_N(t_N)) \ge c N^{1-\epsilon}$ for all sufficiently large $N$. Hence $\lambda_L \le 2\alpha_T^{(W)}$.
\end{proof}

\subsection{Verification of the Static Criterion for Product Thermal States}\label{subsec:product-states}

To illustrate the mechanism behind Hypothesis~(H2) in a setting where the Wightman inner product factorizes across sites, we verify~\eqref{eq:main-H2} with exponent $\delta = 1/2$ for every full-rank product thermal state on $N$ qubits.

\begin{proposition}[Static Filtration Bound for Product States]\label{prop:product}
Let $\rho_\beta = \bigotimes_{x=1}^N \rho_x$ be a full-rank product state on $N$ qubits, where up to single-site unitaries $\rho_x = \frac{1}{2}(I_x + m_x Z_x)$ with $|m_x| \le m_{\max} < 1$ for all $x\in\Lambda$. Then for every $s\in\{1,\dots,N\}$,
\begin{equation}\label{eq:product-bound}
  \Lambda_\beta^{\MSS}(s) \le s + m_{\max}\sqrt{Ns}.
\end{equation}
In particular, Hypothesis~\emph{(H2)} of Theorem~\ref{thm:main-tails-reduction}(b) holds with $\delta = 1/2$.
\end{proposition}

\begin{proof}
Because $\mathcal{S}_\beta$ is invariant under local single-site unitary rotations, we may assume $\rho_x = \frac{1}{2}(I_x + m_x Z_x)$ with $m_x\in(-1,1)$. Set $\lambda_{x,\pm} := \frac{1\pm m_x}{2} > 0$, $\sigma_x := \sqrt{1-m_x^2} = 2\sqrt{\lambda_{x,+}\lambda_{x,-}} > 0$, and $p_x := \sqrt{\lambda_{x,+}} + \sqrt{\lambda_{x,-}}$, $q_x := \sqrt{\lambda_{x,+}} - \sqrt{\lambda_{x,-}}$. Observe that $p_x^2 + q_x^2 = 2$, $p_x q_x = m_x$, $p_x^2 - q_x^2 = 2\sigma_x$, and
\[
  \rho_x^{1/2} = \frac{1}{2}\bigl(p_x I_x + q_x Z_x\bigr).
\]
In the single-site Wightman inner product $(a|b)_x := \Tr(\rho_x^{1/2}a^\dagger\rho_x^{1/2}b)$, let us compute the Gram matrix of $\{I_x, X_x, Y_x, Z_x\}$:
\[
  (I_x|I_x)_x = 1, \qquad
  (X_x|X_x)_x = (Y_x|Y_x)_x = \sigma_x, \qquad
  (Z_x|Z_x)_x = 1, \qquad
  (I_x|Z_x)_x = \Tr(\rho_x Z_x) = m_x,
\]
while all other pairs are orthogonal. Applying Gram--Schmidt orthogonalization to $\{I_x, Z_x\}$ produces the single-site $(\cdot|\cdot)_x$-orthonormal basis
\[
  E_x^{(0)} := I_x, \qquad
  E_x^{(1)} := \sigma_x^{-1/2} X_x, \qquad
  E_x^{(2)} := \sigma_x^{-1/2} Y_x, \qquad
  E_x^{(3)} := \sigma_x^{-1}(Z_x - m_x I_x).
\]
Notice that $\operatorname{span}\{E_x^{(0)}\} = \mathbb{C}I_x$, so the tensor products $E_{\mathbf{a}} := \bigotimes_{x=1}^N E_x^{(a_x)}$ ($\mathbf{a}\in\{0,1,2,3\}^N$) form a $(\cdot|\cdot)_\beta^W$-orthonormal basis of $\Hop$ such that $\mathcal{V}_s = \operatorname{span}\{E_{\mathbf{a}} : |\mathbf{a}|\le s\}$, where $|\mathbf{a}| := \#\{x : a_x \ne 0\}$.

We now compute the single-site matrix elements $M_x(a,b) := \langle\mathcal{J}_x(E_x^{(a)}),\,\mathcal{S}_x\mathcal{J}_x(E_x^{(b)})\rangle_\HS - \delta_{ab}\langle\rho_x^{1/2},\mathcal{S}_x\rho_x^{1/2}\rangle_\HS$, where $\mathcal{J}_x(a) = \rho_x^{1/4}a\rho_x^{1/4}$ and $\mathcal{S}_x$ projects onto traceless single-site operators. For any $a\in\{0,1,2,3\}$, since $\|\mathcal{J}_x(E_x^{(a)})\|_\HS = 1$, we have $\langle\mathcal{J}_x(E_x^{(a)}),\mathcal{S}_x\mathcal{J}_x(E_x^{(b)})\rangle_\HS = \delta_{ab} - \frac{1}{2}\Tr(\rho_x^{1/2}E_x^{(a)})^*\Tr(\rho_x^{1/2}E_x^{(b)})$. Using $\Tr(\rho_x^{1/2}I_x) = p_x$, $\Tr(\rho_x^{1/2}X_x) = \Tr(\rho_x^{1/2}Y_x) = 0$, and, using $p_x^2 = 1 + \sigma_x$ and $m_x = p_x q_x$,
\[
  \Tr\bigl(\rho_x^{1/2}E_x^{(3)}\bigr) = \sigma_x^{-1}\bigl(q_x - m_x p_x\bigr)
  = \sigma_x^{-1} q_x(1 - p_x^2) = -q_x.
\]
Therefore $\langle\rho_x^{1/2},\mathcal{S}_x\rho_x^{1/2}\rangle_\HS = 1 - \frac{1}{2}p_x^2 = \frac{1}{2}q_x^2$, and the subtracted single-site operator $\Delta\mathcal{S}_x := \mathcal{J}_x^\dagger\mathcal{S}_x\mathcal{J}_x - \frac{1}{2}q_x^2\mathbf{1}$ has matrix elements in $\{E_x^{(0)}, E_x^{(1)}, E_x^{(2)}, E_x^{(3)}\}$:
\begin{align*}
  M_x(0,0) &= \left(1 - \frac{1}{2}p_x^2\right) - \frac{1}{2}q_x^2 = 0, \\
  M_x(1,1) &= M_x(2,2) = 1 - \frac{1}{2}q_x^2 \le 1, \\
  M_x(3,3) &= \left(1 - \frac{1}{2}q_x^2\right) - \frac{1}{2}q_x^2 = 1 - q_x^2 = \sigma_x \le 1, \\
  M_x(0,3) &= M_x(3,0) = -\frac{1}{2}p_x(-q_x) = \frac{1}{2}m_x,
\end{align*}
with all other entries equal to zero. Because $\Delta\mathcal{S}_x$ subtracts $\frac{1}{2}q_x^2$ only from the overall expectation $\mathcal{S}_\beta(A) = \sum_{x=1}^N (\langle\mathcal{J}_\beta(A),\mathcal{S}_x\mathcal{J}_\beta(A)\rangle_\HS - \frac{1}{2}q_x^2\|A\|_{\beta,W}^2)$, we have for every $A = \sum_{|\mathbf{a}|\le s} c_{\mathbf{a}} E_{\mathbf{a}}$ with $\sum |c_{\mathbf{a}}|^2 = 1$:
\[
  \mathcal{S}_\beta(A)
  = \sum_{x=1}^N \sum_{\substack{\mathbf{a},\mathbf{b} \\ \mathbf{a}_{x^c}=\mathbf{b}_{x^c}}} c_{\mathbf{a}}^* c_{\mathbf{b}}\, M_x(a_x, b_x).
\]
Splitting into diagonal ($a_x = b_x \ne 0$) and off-diagonal ($\{a_x,b_x\}=\{0,3\}$) terms gives
\begin{align*}
  \mathcal{S}_\beta(A)
  &\le \sum_{|\mathbf{a}|\le s} |c_{\mathbf{a}}|^2 |\mathbf{a}|
    + \sum_{x=1}^N |m_x| \sum_{\substack{\mathbf{a}: a_x=0 \\ |\mathbf{a}|\le s-1}} |c_{\mathbf{a}}|\,|c_{\mathbf{a}+3e_x}|.
\end{align*}
The diagonal sum is at most $s\sum |c_{\mathbf{a}}|^2 = s$. For the off-diagonal sum, applying Cauchy--Schwarz over $(x,\mathbf{a})$ yields
\begin{align*}
  \sum_{x=1}^N |m_x|\sum_{\substack{\mathbf{a}: a_x=0 \\ |\mathbf{a}|\le s-1}} |c_{\mathbf{a}}|\,|c_{\mathbf{a}+3e_x}|
  &\le m_{\max}\left(\sum_{x=1}^N\sum_{\mathbf{a}:a_x=0}|c_{\mathbf{a}}|^2\right)^{1/2}
    \left(\sum_{x=1}^N\sum_{\mathbf{b}:b_x=3}|c_{\mathbf{b}}|^2\right)^{1/2} \\
  &\le m_{\max}\sqrt{N}\left(\sum_{|\mathbf{b}|\le s}|\mathbf{b}|\,|c_{\mathbf{b}}|^2\right)^{1/2}
  \le m_{\max}\sqrt{Ns},
\end{align*}
establishing~\eqref{eq:product-bound}.
\end{proof}

The $\sqrt{Ns}$ term in~\eqref{eq:product-bound} reflects the coherent overlap between an operator of weight $s$ and the $N$-site thermal background $\rho_\beta^{1/2}$ via the off-diagonal matrix element $M_x(0,3) = \frac{1}{2}m_x$. In Section~\ref{sec:entropic} we will recover an analogous $\sqrt{N K_\beta^W(t)}$ scaling for interacting high-temperature Hamiltonians via quantum relative entropy and Gaussian concentration.

\section{Spectral Decay Bounds and Complex-Time OTOC Analyticity}\label{sec:spectral}

We now turn to the spectral and analytic side of the conjecture~\eqref{eq:intro-uogh}, proving Theorem~\ref{thm:main-spectral}. In~\cite{Parker2019,Avdoshkin2020}, the right-hand inequality $2\alpha_T^{(W)} \le 2\pi/\beta$ was deduced from exponential decay of the Wightman spectral measure under the assumption that the Lanczos coefficients $\{b_{n,T}^{(W)}\}$ vary smoothly with $n$ (or converge to a smooth continuum profile). In this section we remove all smoothness hypotheses by combining the nested-commutator cluster bound (Lemma~\ref{lem:nested}) with a monotone Dyck-path comparison against the $\eta=1$ Meixner chain, and we establish the holomorphic strip and time-ordered dominance bound for the MSS-regulated four-point correlator.

\subsection{Exponential Decay of the Wightman Spectral Measure}\label{subsec:spectral-measure}

Let $W$ be a local Hermitian operator supported on $U\subseteq\Lambda$ with $|U|=s_W$ and $\|W\|_\infty < \infty$. In finite volume, letting $\{|a\rangle\}$ be an orthonormal eigenbasis of $H$ with eigenvalues $\{E_a\}$, the Wightman autocorrelation function is
\begin{equation}\label{eq:CbetaW-eigen}
  C_{\beta,W}^W(t) = \Tr\!\bigl(\rho_\beta^{1/2}W(t)\rho_\beta^{1/2}W\bigr)
  = \frac{1}{Z(\beta)}\sum_{a,b} e^{-\frac{\beta}{2}(E_a+E_b)}\,|W_{ab}|^2\,e^{i(E_a-E_b)t}
  = \int_{\mathbb{R}} e^{i\omega t}\,d\mu_{\beta,W}^W(\omega),
\end{equation}
where the Wightman spectral measure
\[
  d\mu_{\beta,W}^W(\omega) := \frac{1}{Z(\beta)}\sum_{a,b} e^{-\frac{\beta}{2}(E_a+E_b)}\,|W_{ab}|^2\,\delta\bigl(\omega-(E_a-E_b)\bigr)\,d\omega
\]
is nonnegative, even ($d\mu_{\beta,W}^W(-\omega) = d\mu_{\beta,W}^W(\omega)$ since $|W_{ab}|^2 = |W_{ba}|^2$), and has total mass $\int_{\mathbb{R}} d\mu_{\beta,W}^W(\omega) = \|W\|_{\beta,W}^2$.

\begin{proof}[Proof of Theorem~\ref{thm:main-spectral}(a)]
Fix any $\kappa'\in(0,\kappa)$. Using the eigenbasis expansion~\eqref{eq:CbetaW-eigen}, we compute the bilateral Laplace transform of $d\mu_{\beta,W}^W$ at $\tau = \frac{\beta}{2}+\kappa'$:
\begin{align*}
  \int_{\mathbb{R}} e^{(\frac{\beta}{2}+\kappa')\omega}\,d\mu_{\beta,W}^W(\omega)
  &= \frac{1}{Z(\beta)}\sum_{a,b} e^{-\frac{\beta}{2}(E_a+E_b) + (\frac{\beta}{2}+\kappa')(E_a-E_b)}\,|W_{ab}|^2 \\
  &= \frac{1}{Z(\beta)}\sum_{a,b} e^{-\beta E_b}\,e^{\kappa'(E_a-E_b)}\,|W_{ab}|^2.
\end{align*}
Relabeling the dummy indices $a\leftrightarrow b$ and recognizing $(W(-i\kappa'))_{ba} = \langle b|e^{\kappa' H} W e^{-\kappa' H}|a\rangle = e^{\kappa'(E_b-E_a)}W_{ba}$, we obtain the identity
\begin{equation}\label{eq:laplace-W}
  \int_{\mathbb{R}} e^{(\frac{\beta}{2}+\kappa')\omega}\,d\mu_{\beta,W}^W(\omega)
  = \Tr\!\bigl(\rho_\beta\,W\,W(-i\kappa')\bigr).
\end{equation}
Because $d\mu_{\beta,W}^W$ is an even nonnegative measure, $e^{(\frac{\beta}{2}+\kappa')|\omega|} \le e^{(\frac{\beta}{2}+\kappa')\omega} + e^{-(\frac{\beta}{2}+\kappa')\omega}$, so~\eqref{eq:laplace-W} implies
\begin{equation}\label{eq:exp-moment-bound}
  \int_{\mathbb{R}} e^{(\frac{\beta}{2}+\kappa')|\omega|}\,d\mu_{\beta,W}^W(\omega)
  \le 2\Tr\!\bigl(\rho_\beta\,W\,W(-i\kappa')\bigr)
  \le 2\|W\|_\infty\,\|W(-i\kappa')\|_\infty.
\end{equation}
Expanding $W(-i\kappa') = \sum_{m=0}^\infty \frac{(\kappa')^m}{m!}\ad_H^m(W)$ and applying Lemma~\ref{lem:nested}(1) with $s = s_W$ yields
\begin{equation}\label{eq:W-imag-norm}
  \|W(-i\kappa')\|_\infty
  \le \|W\|_\infty \sum_{m=0}^\infty \frac{(\kappa'/\kappa)^m}{m!}\prod_{j=0}^{m-1}\left(\frac{s_W}{r-1}+j\right)
  = \|W\|_\infty \left(1 - \frac{\kappa'}{\kappa}\right)^{-s_W/(r-1)} < \infty,
\end{equation}
where we used the generalized binomial series $(1-u)^{-\alpha} = \sum_{m=0}^\infty \frac{u^m}{m!}\prod_{j=0}^{m-1}(\alpha+j)$ for $|u|<1$. Notice that the right-hand side of~\eqref{eq:exp-moment-bound}--\eqref{eq:W-imag-norm} depends only on $\|W\|_\infty$, $s_W$, $r$, and $\kappa'/\kappa$, and is independent of the system size $N$. Consequently,~\eqref{eq:main-measure} holds both in finite volume and for any weak thermodynamic limit of $d\mu_{\beta,W}^W$ (by Fatou's lemma).

Finally, by Morera's theorem and dominated convergence, the Fourier--Laplace integral $C_{\beta,W}^W(z) = \int_{\mathbb{R}} e^{i\omega z}\,d\mu_{\beta,W}^W(\omega)$ defines a holomorphic function on the open strip $|\operatorname{Im}z| < \frac{\beta}{2}+\kappa'$ for every $\kappa'\in(0,\kappa)$, whence $r_W \ge \frac{\beta}{2}+\kappa$.
\end{proof}

\subsection{Dyck-Path Comparison and Proof of Theorem~\ref{thm:main-spectral}(b)}\label{subsec:dyck}

We now remove the smoothness hypothesis from the upper bound on the asymptotic Lanczos growth rate. Let $\nu := \|W\|_{\beta,W}^{-2}\,d\mu_{\beta,W}^W$ be the normalized Wightman spectral measure, and let $\{b_n\}_{n\ge 1} = \{b_{n,T}^{(W)}\}_{n\ge 1}$ be its associated Lanczos sequence (with $b_n > 0$ for all $n\ge 1$). Because $\nu$ has finite exponential moments~\eqref{eq:main-measure}, all polynomial moments $\mu_{2n} := \int_{\mathbb{R}}\omega^{2n}\,d\nu(\omega)$ are finite, and the symmetric tridiagonal Jacobi operator $T$ associated with $\{b_n\}$ satisfies $\mu_{2n} = \langle e_0, T^{2n} e_0\rangle$ for every $n\ge 0$~\cite{Chihara1978}.

\begin{lemma}[Dyck-Path Monotonicity and Shifted Meixner Lower Bound]\label{lem:dyck-meixner}
Let $\{b_n\}_{n\ge 1}$ be a sequence of positive real numbers with moments $\mu_{2n} = \langle e_0, T^{2n} e_0\rangle$. Suppose there exist $\alpha_1 > 0$ and an integer $N_0 \ge 1$ such that $b_n \ge \alpha_1 n$ for all $n\ge N_0$. Then for every integer $n > N_0$,
\begin{equation}\label{eq:dyck-lower}
  \mu_{2n} \ge \left(\prod_{h=1}^{N_0} b_h^2\right)\frac{2}{\pi}\,(2n-2N_0)!\left(\frac{2\alpha_1}{\pi}\right)^{2n-2N_0}.
\end{equation}
\end{lemma}

\begin{proof}
Expanding $\mu_{2n} = \langle e_0, T^{2n} e_0\rangle$ in the standard basis $\{e_h\}_{h\ge 0}$ of $\ell^2(\mathbb{Z}_{\ge 0})$ gives the Dyck-path representation
\[
  \mu_{2n} = \sum_{\gamma\in\mathcal{D}_{2n}} \prod_{k=1}^{2n} b_{\max(\gamma_{k-1},\gamma_k)},
\]
where $\mathcal{D}_{2n}$ is the set of lattice paths $\gamma = (\gamma_0,\gamma_1,\dots,\gamma_{2n})$ with $\gamma_0 = \gamma_{2n} = 0$, $\gamma_k \ge 0$, and $|\gamma_k - \gamma_{k-1}| = 1$. Since every $b_h > 0$, all path weights are strictly positive, so dropping any subset of paths can only decrease the sum.

For $n > N_0$, restrict the sum to those paths $\gamma\in\mathcal{D}_{2n}$ that climb monotonically from $0$ to $N_0$ in the first $N_0$ steps ($\gamma_k = k$ for $0\le k\le N_0$), follow a shifted Dyck path $\gamma_{N_0+j} = N_0 + \tilde{\gamma}_j$ ($0\le j\le 2(n-N_0)$) for some $\tilde{\gamma}\in\mathcal{D}_{2(n-N_0)}$, and descend monotonically from $N_0$ to $0$ in the final $N_0$ steps. The initial climb and final descent contribute the fixed factor $W_{N_0} := \prod_{h=1}^{N_0} b_h^2 > 0$. Along the middle segment, every step occurs between heights $N_0 + h - 1$ and $N_0 + h$ for some $h\ge 1$, and by hypothesis carries weight
\[
  b_{N_0+h} \ge \alpha_1(N_0 + h) \ge \alpha_1 h =: b_h^{\mathrm{M}}(\alpha_1).
\]
Consequently,
\[
  \mu_{2n} \ge W_{N_0}\sum_{\tilde{\gamma}\in\mathcal{D}_{2(n-N_0)}} \prod_{j=1}^{2(n-N_0)} b_{\max(\tilde{\gamma}_{j-1},\tilde{\gamma}_j)}^{\mathrm{M}}(\alpha_1)
  = W_{N_0}\,\mu_{2(n-N_0)}^{\mathrm{M}}(\alpha_1),
\]
where $\mu_{2m}^{\mathrm{M}}(\alpha_1)$ denotes the $2m$-th moment of the $\eta=1$ Meixner chain $b_h^{\mathrm{M}}(\alpha_1) = \alpha_1 h$.

Recall~\cite{Chihara1978,Parker2019} that the normalized spectral measure of the $\eta=1$ Meixner chain $b_h^{\mathrm{M}}(\alpha_1) = \alpha_1 h$ is $d\nu^{\mathrm{M}}(\omega) = \frac{1}{2\alpha_1}\sech(\frac{\pi\omega}{2\alpha_1})\,d\omega$ (whose Fourier transform is $\sech(\alpha_1 t)$). Using the elementary inequality $\sech(u) = \frac{2}{e^u+e^{-u}} \ge e^{-u}$ for all $u\ge 0$, we obtain for every $m\ge 0$:
\[
  \mu_{2m}^{\mathrm{M}}(\alpha_1)
  = \frac{1}{\alpha_1}\int_0^\infty \omega^{2m}\sech\!\left(\frac{\pi\omega}{2\alpha_1}\right)d\omega
  \ge \frac{1}{\alpha_1}\int_0^\infty \omega^{2m} e^{-\frac{\pi\omega}{2\alpha_1}}\,d\omega
  = \frac{2}{\pi}\,(2m)!\left(\frac{2\alpha_1}{\pi}\right)^{2m}.
\]
Setting $m = n - N_0$ yields~\eqref{eq:dyck-lower}.
\end{proof}

\begin{proof}[Proof of Theorem~\ref{thm:main-spectral}(b)]
Let $\alpha_* := \liminf_{n\to\infty} b_{n,T}^{(W)}/n$, and suppose for contradiction that $2\alpha_* > \pi/r_W$. Choose $\alpha_1\in(\pi/(2r_W), \alpha_*)$; then there exists $N_0\ge 1$ such that $b_{n,T}^{(W)} \ge \alpha_1 n$ for all $n\ge N_0$. By Lemma~\ref{lem:dyck-meixner}, for every $n > N_0$,
\[
  \left(\frac{\mu_{2n}}{(2n)!}\right)^{1/(2n)}
  \ge \left(\frac{2W_{N_0}}{\pi}\right)^{1/(2n)} \left(\frac{(2n-2N_0)!}{(2n)!}\right)^{1/(2n)} \left(\frac{2\alpha_1}{\pi}\right)^{1 - N_0/n}.
\]
As $n\to\infty$, the first factor tends to $1$, the third tends to $2\alpha_1/\pi$, and since $\frac{(2n)!}{(2n-2N_0)!} \le (2n)^{2N_0}$ we have $((2n-2N_0)!/(2n)!)^{1/(2n)} \ge (2n)^{-N_0/n} \to 1$. Hence
\begin{equation}\label{eq:moment-liminf}
  \liminf_{n\to\infty}\left(\frac{\mu_{2n}}{(2n)!}\right)^{1/(2n)} \ge \frac{2\alpha_1}{\pi}.
\end{equation}
On the other hand, because $C_{\beta,W}^W(z)/\|W\|_{\beta,W}^2 = \sum_{n=0}^\infty \frac{(-1)^n \mu_{2n}}{(2n)!} z^{2n}$ extends holomorphically to the strip $|\operatorname{Im}z| < r_W$ (and in particular to the open disk $|z| < r_W$), the Cauchy--Hadamard formula yields
\[
  \limsup_{n\to\infty}\left(\frac{\mu_{2n}}{(2n)!}\right)^{1/(2n)} \le \frac{1}{r_W}.
\]
Combining this with~\eqref{eq:moment-liminf} gives $2\alpha_1/\pi \le 1/r_W$, contradicting $\alpha_1 > \pi/(2r_W)$. Therefore $2\alpha_* \le \pi/r_W$, and part~(a) gives $r_W \ge \frac{\beta}{2}+\kappa$, whence~\eqref{eq:main-right} holds.
\end{proof}

\noindent\emph{Remark.}
In the large-$q$ SYK model, the interaction is all-to-all rather than $r$-local with finite $J_{\mathrm{site}}$ ($\kappa = 0$), and $C_{\beta,W}^W(z)$ has a branch point at $z = \pm i\frac{\pi}{2v}\cdot\frac{\beta}{2}$ with $v\in(0,1)$ solving $\frac{\pi v}{\cos(\pi v/2)} = \beta\mathcal{J}$~\cite{MS2016,Parker2019}. Thus $r_W = \frac{\beta}{2v} > \frac{\beta}{2}$, and Theorem~\ref{thm:main-spectral}(b) gives $2\alpha_T^{(W)} \le \frac{\pi}{r_W} = \frac{2\pi v}{\beta}$, which matches the large-$q$ SYK value $2\alpha_T^{(W)} = \lambda_{L,T} = \frac{2\pi v}{\beta}$ across all temperatures.

\subsection{Complex-Time Analyticity of the Regulated Four-Point Correlator}\label{subsec:otoc-strip}

We finish the proof of Theorem~\ref{thm:main-spectral} by deriving the holomorphic strip $\mathcal{S}_{\mathrm{OTOC}}$ and the time-ordered dominance inequality~\eqref{eq:main-otoc-bound} for the MSS-regulated four-point correlator.

\begin{proof}[Proof of Theorem~\ref{thm:main-spectral}(c)]
Fix $z = t + i\tau\in\mathbb{C}$ with $|\tau| < \frac{1}{2}\max(r_W,r_V)$. Define the operator
\[
  A(t,\tau) := \rho_\beta^{1/4} W(t+i\tau) \rho_\beta^{1/4} V,
\]
so that $F(t+i\tau) = \Tr(A(t,\tau)^2) = \langle A(t,\tau)^\dagger, A(t,\tau)\rangle_\HS$. By the Cauchy--Schwarz inequality in $(\Hop,\langle\cdot,\cdot\rangle_\HS)$ and the identity $\|A^\dagger\|_\HS = \|A\|_\HS$,
\begin{equation}\label{eq:F-Ttau}
  |F(t+i\tau)| \le \|A(t,\tau)^\dagger\|_\HS\,\|A(t,\tau)\|_\HS = \|A(t,\tau)\|_\HS^2 = \Tr\bigl(A(t,\tau)^\dagger A(t,\tau)\bigr).
\end{equation}
Using~\eqref{eq:prelim-adj}, $A(t,\tau)^\dagger = V\rho_\beta^{1/4} W(t-i\tau)\rho_\beta^{1/4}$, so cyclicity of the trace yields
\[
  \Tr\bigl(A(t,\tau)^\dagger A(t,\tau)\bigr)
  = \Tr\!\bigl(\rho_\beta^{1/4} W(t-i\tau)\rho_\beta^{1/2} W(t+i\tau)\rho_\beta^{1/4} V^2\bigr)
  = T_\tau(t),
\]
establishing the third bound in~\eqref{eq:main-otoc-bound}.

Next, set $B(t,\tau) := \rho_\beta^{1/4} W(t+i\tau)\rho_\beta^{1/4}$, so that $B(t,\tau)^\dagger B(t,\tau) = \rho_\beta^{1/4} W(t-i\tau)\rho_\beta^{1/2} W(t+i\tau)\rho_\beta^{1/4} \succeq 0$. Since $V^2 \preceq \|V\|_\infty^2\mathbf{1}$, we obtain
\begin{align*}
  T_\tau(t)
  &= \Tr\bigl(B(t,\tau)^\dagger B(t,\tau) V^2\bigr)
   \le \|V\|_\infty^2 \Tr\bigl(B(t,\tau)^\dagger B(t,\tau)\bigr) \\
  &= \|V\|_\infty^2 \Tr\!\bigl(\rho_\beta^{1/2} W(t-i\tau)\rho_\beta^{1/2} W(t+i\tau)\bigr).
\end{align*}
By time-translation invariance of $\rho_\beta^{1/2}$, shifting by $e^{-itH}(\cdot)e^{itH}$ eliminates $t$ and leaves
\[
  \Tr\!\bigl(\rho_\beta^{1/2} W(-i\tau)\rho_\beta^{1/2} W(i\tau)\bigr)
  = \Tr\!\bigl(\rho_\beta^{1/2} W(-2i\tau)\rho_\beta^{1/2} W\bigr)
  = C_{\beta,W}^W(2i\tau),
\]
since $C_{\beta,W}^W(2i\tau) = \int_{\mathbb{R}} e^{-2\tau\omega}\,d\mu_{\beta,W}^W(\omega)$ is even in $\tau$. Thus $|F(t+i\tau)| \le T_\tau(t) \le \|V\|_\infty^2 C_{\beta,W}^W(2i\tau)$.

Finally, by time-translation invariance and cyclicity of the trace,
\[
  F(t+i\tau) = \Tr\!\bigl(\rho_\beta^{1/4} W \rho_\beta^{1/4} V(-t-i\tau)\rho_\beta^{1/4} W \rho_\beta^{1/4} V(-t-i\tau)\bigr).
\]
Repeating the same calculation with $(W,V,t+i\tau)$ replaced by $(V,W,-t-i\tau)$ yields $|F(t+i\tau)| \le \|W\|_\infty^2 C_{\beta,V}^W(2i\tau)$. This proves~\eqref{eq:main-otoc-bound}.

In finite volume, $F_N(z)$ is entire and satisfies~\eqref{eq:main-otoc-bound} on $\mathbb{C}$. On every compact subset $K\subset\{z\in\mathbb{C} : |\operatorname{Im}z| < \frac{\beta}{4}+\frac{\kappa}{2}\}$, the bound~\eqref{eq:exp-moment-bound}--\eqref{eq:W-imag-norm} yields a uniform-in-$N$ upper bound on $C_{\beta,W,N}^W(2i\tau)$, so by Montel's and Vitali's theorems any thermodynamic limit $F(t)$ of $F_N(t)$ along $\mathbb{R}$ extends holomorphically to $\{|\operatorname{Im}z| < \frac{\beta}{4}+\frac{\kappa}{2}\}$ and satisfies~\eqref{eq:main-otoc-bound} upon taking $\liminf_{N\to\infty}$ on the right-hand side (or directly with the limiting correlators whenever $C_{\beta,W,N}^W(2i\tau)$, $C_{\beta,V,N}^W(2i\tau)$, and $T_{\tau,N}(t)$ converge along the same subsequence). More generally, if $C_{\beta,W,N}^W(2i\tau)$ (respectively, $C_{\beta,V,N}^W(2i\tau)$) is uniformly bounded in $N$ on compact subsets of $|\tau| < r_W/2$ (respectively, $r_V/2$), the same normal-families argument extends holomorphy and~\eqref{eq:main-otoc-bound} to $\mathcal{S}_{\mathrm{OTOC}}$.
\end{proof}

\noindent\emph{Remark.}
In a clustering thermal state, at real times $t$ well past the dissipation timescale $t_d$, the time-ordered correlator $T_\tau(t) = \Tr(\rho_\beta^{1/4}W(t-i\tau)\rho_\beta^{1/2}W(t+i\tau)\rho_\beta^{1/4}V^2)$ factorizes asymptotically to $C_{\beta,W}^W(2i\tau)\langle V^2\rangle_\beta$, matching the time-ordered bound used in~\cite{MSS2016}.

\section{Quantum Glauber Bound and Entropic Background Decoupling}\label{sec:entropic}

We now address the second structural obstruction of Subsection~\ref{subsec:obstructions}: separating the dynamical growth of $\mathcal{O}(t)$ from the extensive $O(N)$ operator size of the thermal state $\rho_\beta^{1/2}$ inside $\mathcal{S}_\beta(\mathcal{O}(t))$, and connecting the summed commutators $C_W(A)$ and $\mathrm{I}(A)$ to the Wightman filtration and quantum Glauber superoperators. We begin with a general Wightman commutator estimate and its high-temperature Glauber bound, and then establish the self-contained trace cumulant and entropic background decoupling estimates.

\subsection{Wightman Commutator and Quantum Glauber Bounds}\label{subsec:1d-comm}

Because the Wightman inner product is unitarily equivalent to the Hilbert--Schmidt inner product via $\mathcal{J}_\beta(A) = \rho_\beta^{1/4}A\rho_\beta^{1/4}$ (Proposition~\ref{prop:isometry}), left- and right-multiplication by a local operator $V_x$ in $(\Hop,(\cdot|\cdot)_\beta^W)$ are controlled by the operator norm of the imaginary-time evolved operator $V_x(\pm i\beta/4)$.

\begin{lemma}[Single-Site Wightman Commutator Bound]\label{lem:wightman-comm}
Let $V_x$ be a Hermitian operator supported on site $x\in\Lambda$, and set $M_\beta(V_x) := \|V_x(i\beta/4)\|_\infty = \|V_x(-i\beta/4)\|_\infty$. Then for every $A\in\Hop$,
\begin{equation}\label{eq:single-site-comm}
  \bigl\|[V_x, A]\bigr\|_{\beta,W}
  \le 2 M_\beta(V_x)\,\bigl\|(\mathbf{1}-P_{x,0}^{(\beta)})A\bigr\|_{\beta,W}
  \le 2 M_\beta(V_x)\,\|\mathcal{S}_x(A)\|_{\beta,W}.
\end{equation}
\end{lemma}

\begin{proof}
Because $V_x$ is supported on site $x$, it commutes with every operator in the complementary subalgebra $\mathcal{V}_{(x^c)}$. Both $P_{x,0}^{(\beta)}A$ and $\mathcal{E}_x(A)$ lie in $\mathcal{V}_{(x^c)}$, so $[V_x, A] = [V_x, B]$ where $B := (\mathbf{1}-P_{x,0}^{(\beta)})A$. By Proposition~\ref{prop:isometry}(1) and~\eqref{eq:prelim-imag},
\begin{align*}
  \|V_x B\|_{\beta,W}
  &= \bigl\|\rho_\beta^{1/4} V_x B \rho_\beta^{1/4}\bigr\|_\HS
   = \bigl\|V_x(i\beta/4)\,\mathcal{J}_\beta(B)\bigr\|_\HS
   \le \|V_x(i\beta/4)\|_\infty\,\|B\|_{\beta,W}, \\
  \|B V_x\|_{\beta,W}
  &= \bigl\|\rho_\beta^{1/4} B V_x \rho_\beta^{1/4}\bigr\|_\HS
   = \bigl\|\mathcal{J}_\beta(B)\,V_x(-i\beta/4)\bigr\|_\HS
   \le \|V_x(-i\beta/4)\|_\infty\,\|B\|_{\beta,W}.
\end{align*}
Since $V_x = V_x^\dagger$,~\eqref{eq:prelim-adj} gives $\|V_x(-i\beta/4)\|_\infty = \|V_x(i\beta/4)^\dagger\|_\infty = M_\beta(V_x)$. The triangle inequality yields $\|[V_x,A]\|_{\beta,W} = \|[V_x,B]\|_{\beta,W} \le 2M_\beta(V_x)\|B\|_{\beta,W}$. Finally, since $P_{x,0}^{(\beta)}A$ is the $(\cdot|\cdot)_\beta^W$-best approximation to $A$ in $\mathcal{V}_{(x^c)}$ and $\mathcal{E}_x(A)\in\mathcal{V}_{(x^c)}$, we have $\|(\mathbf{1}-P_{x,0}^{(\beta)})A\|_{\beta,W} \le \|A-\mathcal{E}_x(A)\|_{\beta,W} = \|\mathcal{S}_x(A)\|_{\beta,W}$.
\end{proof}

In arbitrary spatial dimension $d\ge 1$ and high temperature $\beta\in[0,4\kappa)$, combining Lemma~\ref{lem:wightman-comm} with the nested-commutator estimate of Lemma~\ref{lem:nested}(2) immediately compares $C_W(A)$ to the quantum Glauber generator $\widetilde{\mathcal{Q}}^{(\beta)}$.

\begin{proposition}[High-Temperature Quantum Glauber Bound]\label{prop:glauber}
Let $H = \sum_{X\in\mathcal{E}} h_X$ be an $r$-local Hamiltonian ($r\ge 2$) on $N$ qubits with $J_{\mathrm{site}} \in (0,\infty)$ and $\kappa = (2(r-1)J_{\mathrm{site}})^{-1}$. Then for every $\beta\in[0,4\kappa)$ and every $A\in\Hop$,
\begin{equation}\label{eq:glauber-CW}
  C_W(A) \le \frac{3}{2}\left(1 - \frac{\beta}{4\kappa}\right)^{-2}\bigl(A\big|\widetilde{\mathcal{Q}}^{(\beta)}\big|A\bigr)_\beta^W.
\end{equation}
\end{proposition}

\begin{proof}
For $\beta\in[0,4\kappa)$, Lemma~\ref{lem:nested}(2) gives $M_\beta(\sigma_x^\alpha) = \|\sigma_x^\alpha(i\beta/4)\|_\infty \le (1-\beta/(4\kappa))^{-1}$ for every $x\in\Lambda$ and $\alpha\in\{1,2,3\}$. Substituting this bound into Lemma~\ref{lem:wightman-comm} and summing over $x\in\Lambda$ and $\alpha\in\{1,2,3\}$ with prefactor $1/8$ gives~\eqref{eq:glauber-CW}.
\end{proof}

\subsection{Entropic Background Decoupling and High-Temperature Estimates}\label{subsec:entropic}

We now turn to the MSS-regulated thermal operator size $\mathcal{S}_\beta(\mathcal{O}(t))$ in arbitrary spatial dimension $d\ge 1$. Commuting $\sigma_x^\alpha$ past $\mathcal{J}_\beta(A) = \rho_\beta^{1/4}A\rho_\beta^{1/4}$ separates cleanly into a dynamical commutator term (yielding $\mathrm{I}(A)$) and a thermal-background term whose expectation value in the time-evolved density matrix $\rho_{\mathcal{O}(t)} := \mathcal{J}_\beta(\mathcal{O}(t))^2$ is controlled uniformly in $t\in\mathbb{R}$ by unitary invariance of quantum relative entropy combined with a self-contained Koteck\'y--Preiss cumulant estimate on the normalized trace.

\begin{lemma}[High-Temperature Trace Cumulant Bound]\label{lem:kp-cumulant}
Let $H = \sum_{X\in\mathcal{E}} h_X$ be an $r$-local Hamiltonian ($r\ge 2$) on $N$ qubits with $J_{\mathrm{site}}\in(0,\infty)$, and suppose $\beta \le (8e^2 r J_{\mathrm{site}})^{-1}$. Let $X = \sum_{Z\in\mathcal{E}_X} X_Z$ be a Hermitian operator where each $X_Z$ is supported on $Z\subseteq\Lambda$ with $|Z| \le k_r$ ($k_r \ge r$) and per-site norm $g_X := \sup_{y\in\Lambda}\sum_{Z\ni y}\|X_Z\|_\infty < \infty$. Setting $\Lambda'_0 := \bigl(16\,k_r e^{1 + k_r/r} g_X\bigr)^{-1}$, there exists a numerical constant $C_0 > 0$ (independent of $N, k_r, g_X, \beta$) such that for all real $|\lambda| \le \Lambda'_0/2$,
\begin{equation}\label{eq:kp-cumulant}
  \ln\Tr\!\bigl(\rho_\beta\,e^{\lambda(X - \langle X\rangle_\beta)}\bigr) \le C_0\,N\,\frac{\lambda^2}{(\Lambda'_0)^2}.
\end{equation}
\end{lemma}

\begin{proof}
Let $\tau(A) := 2^{-N}\Tr(A)$ be the normalized trace, which satisfies $|\tau(A)|\le\|A\|_\infty$ and factorizes multiplicatively over operators with pairwise disjoint supports. Expanding $e^{-\beta H}$ and $e^{\zeta X}$ ($\zeta\in\mathbb{C}$) in powers of $-\beta h_Y$ ($Y\in\mathcal{E}$) and $\zeta X_Z$ ($Z\in\mathcal{E}_X$) and grouping into connected components $\gamma = (\gamma_H,\gamma_X)$ on the disjoint-union hypergraph $\mathcal{F} := \mathcal{E}\sqcup\mathcal{E}_X$ with support $U_\gamma := (\bigcup_{Y\in\gamma_H}Y)\cup(\bigcup_{Z\in\gamma_X}Z)$ gives
\[
  \tau\bigl(e^{-\beta H}e^{\zeta X}\bigr) = \sum_{\substack{\mathcal{G}\text{ pairwise}\\\text{support-disjoint}}} \prod_{\gamma\in\mathcal{G}} \tau\bigl(\widetilde{Q}_\gamma(\beta,\zeta)\bigr),
\]
where by the multinomial theorem $\|\widetilde{Q}_\gamma(\beta,\zeta)\|_\infty \le \prod_{Y\in\gamma_H}(e^{\beta\|h_Y\|_\infty}-1)\prod_{Z\in\gamma_X}(e^{|\zeta|\|X_Z\|_\infty}-1) =: z_\gamma(\beta,|\zeta|)$. Choose the Koteck\'y--Preiss weight $a_\gamma := \frac{|U_\gamma|}{2r}$ and define the base weights $\mu(Y) := e^{1/2}(e^{\beta\|h_Y\|_\infty}-1)$ ($Y\in\mathcal{E}$) and $\mu(Z) := e^{k_r/(2r)}(e^{|\zeta|\|X_Z\|_\infty}-1)$ ($Z\in\mathcal{E}_X$), so that $z_\gamma(\beta,|\zeta|)e^{a_\gamma} \le \prod_{E\in\gamma}\mu(E)$. Setting $c_0 := \frac{1}{2r}$, for $|\zeta|\le\Lambda'_0$ and $\beta\le(8e^2 r J_{\mathrm{site}})^{-1}$, we verify the exponential-pinning sum (using $e^u-1\le u e^u$ and $k_r\ge r$):
\begin{align*}
  \sup_{y\in\Lambda}\sum_{Y\in\mathcal{E}:\, y\in Y} \mu(Y)e^{|Y|/(2r)}
  &\le \sup_{y\in\Lambda}\sum_{Y\ni y} e\,\beta\|h_Y\|_\infty e^{\beta J_{\mathrm{site}}}
   \le e\cdot\frac{1}{8e^2 r}e^{1/(16e^2)} < \frac{1}{8r}, \\
  \sup_{y\in\Lambda}\sum_{Z\in\mathcal{E}_X:\, y\in Z} \mu(Z)e^{|Z|/(2r)}
  &\le \sup_{y\in\Lambda}\sum_{Z\ni y} e^{k_r/r}\Lambda'_0\|X_Z\|_\infty e^{\Lambda'_0 g_X}
   \le \frac{e^{\Lambda'_0 g_X}}{16e k_r} < \frac{1}{8r}.
\end{align*}
Thus $\Phi := \sup_{y\in\Lambda}\sum_{E\in\mathcal{F}:\, y\in E}\mu(E)e^{c_0|E|} < \frac{1}{4r} < \frac{1}{2r} = c_0$. Applying the multi-type unordered spanning-tree bound (Lemma~\ref{lem:app-multitree} in Appendix~\ref{app:trees}) yields, for every $y\in\Lambda$,
\[
  \sum_{\gamma:\, y\in U_\gamma} z_\gamma(\beta,|\zeta|) e^{a_\gamma}
  \le \Phi < \frac{1}{4r}, \qquad\text{so}\qquad
  \sum_{\gamma':\, U_{\gamma'}\cap U_\gamma\ne\emptyset} z_{\gamma'} e^{a_{\gamma'}} \le \frac{|U_\gamma|}{4r} = \frac{1}{2}a_\gamma < a_\gamma.
\]
By the Koteck\'y--Preiss theorem~\cite{Kotecky1986,Ueltschi2004}, the cluster expansion for $f(\zeta) := \ln\tau(e^{-\beta H}e^{\zeta X}) - \ln\tau(e^{-\beta H}) - \zeta\langle X\rangle_\beta$ converges absolutely and holomorphically on $|\zeta|\le\Lambda'_0$, with $|\ln\tau(e^{-\beta H}e^{\zeta X}) - \ln\tau(e^{-\beta H})| \le \sum_{y\in\Lambda}\frac{1}{4r}\le N$ and $|\zeta\langle X\rangle_\beta|\le N$ (since $\|X\|_\infty \le N g_X$ and $\Lambda'_0 N g_X \le N$). Thus $|f(\zeta)|\le 2N$ on $|\zeta|\le\Lambda'_0$. Since $f(0) = f'(0) = 0$, the Schwarz lemma (applied to $f(\zeta)/\zeta^2$) yields $|f(\lambda)| \le 8N\lambda^2/(\Lambda'_0)^2$ for all $|\lambda|\le\Lambda'_0/2$.
\end{proof}

\begin{proposition}[Entropic Background Decoupling]\label{prop:entropic}
Let $H = \sum_{X\in\mathcal{E}} h_X$ be an $r$-local Hamiltonian ($r\ge 2$) on $N$ qubits with $J_{\mathrm{site}}\in(0,\infty)$ and $\kappa = (2(r-1)J_{\mathrm{site}})^{-1}$. Let $\mathcal{O}$ be an $s_0$-local Hermitian operator normalized to $\|\mathcal{O}\|_{\beta,W} = 1$.
\begin{enumerate}
  \setlength{\itemsep}{3pt}\setlength{\parskip}{0pt}
  \item[\emph{(1)}] \emph{Three-term decomposition.} For every Hermitian $A\in\Hop$ with $\|A\|_{\beta,W}=1$,
  \begin{equation}\label{eq:three-term}
    \mathcal{S}_\beta(A) = \mathrm{I}(A) + \mathrm{II}(A) + \mathrm{III}(A),
  \end{equation}
  where $\mathrm{I}(A)$ is given by~\eqref{eq:intro-I}, and setting $D_{x,\alpha} := [\sigma_x^\alpha,\rho_\beta^{1/2}]$, $\rho_A := \mathcal{J}_\beta(A)^2$, and
  \begin{equation}\label{eq:Bbeta-def}
    \mathcal{B}_\beta := \frac{1}{8}\sum_{x=1}^N\sum_{\alpha=1}^3 \bigl(\sigma_x^\alpha(-i\beta/2)-\sigma_x^\alpha\bigr)\bigl(\sigma_x^\alpha(-i\beta/2)-\sigma_x^\alpha\bigr)^\dagger \succeq 0,
  \end{equation}
  we have
  \begin{align}
    \mathrm{II}(A) &= \frac{1}{4}\sum_{x=1}^N\sum_{\alpha=1}^3 \operatorname{Re}\bigl\langle\mathcal{J}_\beta([\sigma_x^\alpha(-i\beta/4),A]),\, A(i\beta/4)D_{x,\alpha}\bigr\rangle_\HS, \label{eq:II-def}\\
    \mathrm{III}(A) &= \Tr\!\bigl((\rho_A - \rho_\beta)\mathcal{B}_\beta\bigr). \label{eq:III-def}
  \end{align}
  \item[\emph{(2)}] \emph{High-temperature background decoupling.} For every $\beta\in[0,\beta_{\mathrm{wc}}]$ (and more generally whenever $\beta \le (8e^2 r J_{\mathrm{site}})^{-1}$ and $\beta < (erJ_{\mathrm{site}})^{-1}$) and every $s_0$-local Hermitian $\mathcal{O}$ with $\|\mathcal{O}\|_{\beta,W}=1$ and $\|\mathcal{O}\|_\infty = O(1)$, there exist $N$-independent constants $C_1'(\beta), C_2(\beta) > 0$ such that for all $t\in\mathbb{R}$,
  \begin{equation}\label{eq:entropic-decoupling}
    \bigl|\mathrm{III}(\mathcal{O}(t))\bigr| \le C_2(\beta)\sqrt{N}, \qquad
    \mathcal{S}_\beta(\mathcal{O}(t)) \le \mathrm{I}(\mathcal{O}(t)) + C_1'(\beta)\sqrt{N\,\mathrm{I}(\mathcal{O}(t))} + C_2(\beta)\sqrt{N}.
  \end{equation}
\end{enumerate}
\end{proposition}

\begin{proof}
\emph{Part~(1).} For each site $x\in\Lambda$ and $\alpha\in\{1,2,3\}$, using~\eqref{eq:prelim-imag} we rewrite the commutator $[\sigma_x^\alpha,\mathcal{J}_\beta(A)]$ as
\begin{align*}
  [\sigma_x^\alpha, \rho_\beta^{1/4} A \rho_\beta^{1/4}]
  &= \rho_\beta^{1/4}\sigma_x^\alpha(-i\beta/4)A\rho_\beta^{1/4} - \rho_\beta^{1/4} A \sigma_x^\alpha(i\beta/4)\rho_\beta^{1/4} \\
  &= \rho_\beta^{1/4}\bigl[\sigma_x^\alpha(-i\beta/4),A\bigr]\rho_\beta^{1/4}
     + \rho_\beta^{1/4} A \bigl(\sigma_x^\alpha(-i\beta/4) - \sigma_x^\alpha(i\beta/4)\bigr)\rho_\beta^{1/4}.
\end{align*}
Since $\rho_\beta^{1/4} A (\sigma_x^\alpha(-i\beta/4)-\sigma_x^\alpha(i\beta/4))\rho_\beta^{1/4} = \rho_\beta^{1/4}A\rho_\beta^{-1/4}(\sigma_x^\alpha\rho_\beta^{1/2} - \rho_\beta^{1/2}\sigma_x^\alpha) = A(i\beta/4)D_{x,\alpha}$, setting $X_{x,\alpha} := \mathcal{J}_\beta([\sigma_x^\alpha(-i\beta/4),A])$ and $Y_{x,\alpha} := A(i\beta/4)D_{x,\alpha}$ yields $[\sigma_x^\alpha,\mathcal{J}_\beta(A)] = X_{x,\alpha} + Y_{x,\alpha}$. Expanding $\|[ \sigma_x^\alpha,\mathcal{J}_\beta(A)]\|_\HS^2 = \|X_{x,\alpha}\|_\HS^2 + 2\operatorname{Re}\langle X_{x,\alpha}, Y_{x,\alpha}\rangle_\HS + \|Y_{x,\alpha}\|_\HS^2$, summing over $x$ and $\alpha$ with factor $1/8$, and subtracting $\langle\rho_\beta^{1/2},\mathcal{S}\rho_\beta^{1/2}\rangle_\HS = \frac{1}{8}\sum_{x,\alpha}\|D_{x,\alpha}\|_\HS^2$ gives~\eqref{eq:three-term} with $\mathrm{I}(A) = \frac{1}{8}\sum_{x,\alpha}\|X_{x,\alpha}\|_\HS^2$ (by Proposition~\ref{prop:isometry}(1)) and $\mathrm{II}(A)$ given by~\eqref{eq:II-def}.

To rewrite the third term, observe that $\rho_\beta^{-1/2}D_{x,\alpha} = \sigma_x^\alpha(-i\beta/2) - \sigma_x^\alpha$ and $D_{x,\alpha}^\dagger\rho_\beta^{-1/2} = (\sigma_x^\alpha(-i\beta/2)-\sigma_x^\alpha)^\dagger$. Because $Y_{x,\alpha} = \rho_\beta^{1/4}A\rho_\beta^{-1/4}D_{x,\alpha}$ and $A = A^\dagger$, cyclicity of the trace gives
\begin{align*}
  \|Y_{x,\alpha}\|_\HS^2
  &= \Tr\!\bigl(D_{x,\alpha}^\dagger \rho_\beta^{-1/4} A \rho_\beta^{1/2} A \rho_\beta^{-1/4} D_{x,\alpha}\bigr) \\
  &= \Tr\!\bigl(\rho_\beta^{1/4} A \rho_\beta^{1/2} A \rho_\beta^{1/4}\,\rho_\beta^{-1/2} D_{x,\alpha} D_{x,\alpha}^\dagger \rho_\beta^{-1/2}\bigr)
   = \Tr\!\bigl(\rho_A\,(\rho_\beta^{-1/2} D_{x,\alpha})(D_{x,\alpha}^\dagger \rho_\beta^{-1/2})\bigr),
\end{align*}
where $\rho_A = \mathcal{J}_\beta(A)^2 \succeq 0$ has unit trace $\Tr(\rho_A) = \|A\|_{\beta,W}^2 = 1$. Likewise,
\[
  \Tr\bigl(\rho_\beta(\rho_\beta^{-1/2}D_{x,\alpha})(D_{x,\alpha}^\dagger\rho_\beta^{-1/2})\bigr)
  = \Tr(D_{x,\alpha}D_{x,\alpha}^\dagger) = \|D_{x,\alpha}\|_\HS^2.
\]
Summing over $x,\alpha$ with factor $1/8$ yields~\eqref{eq:III-def}.

\emph{Part~(2).} Set $A = \mathcal{O}(t)$. Because $\mathcal{J}_\beta(\mathcal{O}(t)) = e^{iHt}\mathcal{J}_\beta(\mathcal{O})e^{-iHt}$, we have $\rho_{\mathcal{O}(t)} = e^{iHt}\rho_{\mathcal{O}}e^{-iHt}$. By unitary invariance of the quantum relative entropy $D(\rho\|\sigma) := \Tr(\rho(\ln\rho - \ln\sigma))$ and $[e^{iHt},\rho_\beta]=0$,
\[
  D\bigl(\rho_{\mathcal{O}(t)}\,\big\|\,\rho_\beta\bigr) = D\bigl(\rho_{\mathcal{O}}\,\big\|\,\rho_\beta\bigr) \qquad \forall t\in\mathbb{R}.
\]
To bound $D(\rho_{\mathcal{O}}\|\rho_\beta)$ independently of $N$, write $\mathcal{J}_\beta(\mathcal{O}) = \rho_\beta^{1/2}\mathcal{O}(-i\beta/4) = \mathcal{O}(i\beta/4)\rho_\beta^{1/2}$, so that
\[
  \rho_{\mathcal{O}} = \rho_\beta^{1/2}\,\mathcal{O}(-i\beta/4)\mathcal{O}(i\beta/4)\,\rho_\beta^{1/2}
  = \rho_\beta^{1/2}\,\mathcal{O}(i\beta/4)^\dagger\mathcal{O}(i\beta/4)\,\rho_\beta^{1/2}
  \preceq \|\mathcal{O}(i\beta/4)\|_\infty^2\,\rho_\beta.
\]
Because $\rho_{\mathcal{O}} \preceq \|\mathcal{O}(i\beta/4)\|_\infty^2\,\rho_\beta$, the max-relative entropy satisfies $D(\rho_{\mathcal{O}}\|\rho_\beta) \le D_{\max}(\rho_{\mathcal{O}}\|\rho_\beta) \le 2\ln\|\mathcal{O}(i\beta/4)\|_\infty$ (which holds even when $\rho_{\mathcal{O}}$ is singular), whence
\begin{equation}\label{eq:rel-ent-bound}
  D\bigl(\rho_{\mathcal{O}(t)}\,\big\|\,\rho_\beta\bigr) = D(\rho_{\mathcal{O}}\|\rho_\beta) \le 2\ln\|\mathcal{O}(i\beta/4)\|_\infty.
\end{equation}
Since $\beta < 2\kappa < 4\kappa$ and $\mathcal{O}$ is $s_0$-local with $\|\mathcal{O}\|_\infty = O(1)$, Lemma~\ref{lem:nested}(1) gives $\|\mathcal{O}(i\beta/4)\|_\infty \le \|\mathcal{O}\|_\infty (1 - \beta/(4\kappa))^{-s_0/(r-1)} = O(1)$, so $S_0 := D(\rho_{\mathcal{O}}\|\rho_\beta)$ is bounded uniformly in $N$ and $t$.

Now expand $\mathcal{B}_\beta = \sum_{k=2}^\infty \mathcal{B}_\beta^{(k)}$ into individual $(1+k(r-1))$-local Hermitian terms using the causal-sequence expansion of Lemma~\ref{lem:nested}(1):
\begin{align*}
  \sigma_x^\alpha(-i\beta/2) - \sigma_x^\alpha &= \sum_{m=1}^\infty \frac{(\beta/2)^m}{m!}\sum_{\Gamma\in\mathcal{C}_m(x)} T_\Gamma^{(x,\alpha)}, \\
  \mathcal{B}_\beta^{(k)} &:= \frac{1}{8}\sum_{x=1}^N\sum_{\alpha=1}^3 \sum_{\substack{m_1,m_2\ge 1\\ m_1+m_2=k}} \sum_{\substack{\Gamma_1\in\mathcal{C}_{m_1}(x)\\ \Gamma_2\in\mathcal{C}_{m_2}(x)}} \frac{(\beta/2)^k}{m_1!m_2!}\,B_{x,\alpha,\Gamma_1,\Gamma_2},
\end{align*}
where $B_{x,\alpha,\Gamma_1,\Gamma_2} := \frac{1}{2}\bigl(T_{\Gamma_1}^{(x,\alpha)}(T_{\Gamma_2}^{(x,\alpha)})^\dagger + T_{\Gamma_2}^{(x,\alpha)}(T_{\Gamma_1}^{(x,\alpha)})^\dagger\bigr)$ is Hermitian and supported on $U_{\Gamma_1}\cup U_{\Gamma_2}$ of cardinality at most $1 + k(r-1) =: k_r$. To apply Lemma~\ref{lem:kp-cumulant}, we bound the per-site norm $g_{\mathcal{B}_\beta^{(k)}}$ for any fixed site $y\in\Lambda$: if $y\in U_{\Gamma_1}\cup U_{\Gamma_2}$, then $y\in U_{\Gamma_1}$ or $y\in U_{\Gamma_2}$. Applying the Cayley spanning-tree root-reversal estimate (Lemma~\ref{lem:app-cayley} in Appendix~\ref{app:trees}, picking up $2^{m_1}$ from the $m_1$ nested commutators) to the leg containing $y$ and Lemma~\ref{lem:nested}(1) to the other leg rooted at $x$ yields
\[
  g_{\mathcal{B}_\beta^{(k)}} \le 2\cdot\frac{3}{8}\sum_{\substack{m_1,m_2\ge 1\\ m_1+m_2=k}} (\beta e r J_{\mathrm{site}})^{m_1}\left(\frac{\beta}{2\kappa}\right)^{m_2}
  \le \frac{3}{4}(k-1)(\beta e r J_{\mathrm{site}})^k.
\]
For each shell $k\ge 2$, the Gibbs--Donsker--Varadhan variational inequality (which follows from Peierls--Bogoliubov and the Golden--Thompson inequality $\Tr(e^{\ln\rho_\beta + \lambda X}) \le \Tr(\rho_\beta e^{\lambda X})$ for non-commuting $\rho_\beta$ and $X$) gives, for every $\lambda > 0$,
\begin{equation}\label{eq:gdv}
  \pm\Tr\!\bigl((\rho_{\mathcal{O}(t)}-\rho_\beta)\mathcal{B}_\beta^{(k)}\bigr)
  \le \frac{1}{\lambda}\left(S_0 + \ln\Tr\!\bigl(\rho_\beta\,e^{\pm\lambda(\mathcal{B}_\beta^{(k)}-\langle\mathcal{B}_\beta^{(k)}\rangle_\beta)}\bigr)\right).
\end{equation}
Applying Lemma~\ref{lem:kp-cumulant} to $X = \mathcal{B}_\beta^{(k)}$ with $\Lambda'_k := (16 k_r e^{1+k_r/r} g_{\mathcal{B}_\beta^{(k)}})^{-1}$ and choosing $\lambda = \min(\Lambda'_k/2, \sqrt{S_0/(C_0 N)}\Lambda'_k)$ yields, for all $N\ge 4S_0/C_0$,
\[
  \bigl|\Tr\!\bigl((\rho_{\mathcal{O}(t)}-\rho_\beta)\mathcal{B}_\beta^{(k)}\bigr)\bigr|
  \le \frac{2\sqrt{C_0 S_0}}{\Lambda'_k}\sqrt{N}
  \le C'(\beta)\,k_r e^{k_r/r}(k-1)(\beta e r J_{\mathrm{site}})^k\sqrt{N}.
\]
Because $\frac{k_r}{r} \le \frac{1}{r} + k$ and $\beta\le (8e^2 r J_{\mathrm{site}})^{-1}$ implies $\beta e^2 r J_{\mathrm{site}} \le \frac{1}{8} < 1$ (and for $\beta\le\beta_{\mathrm{wc}}$, $\beta_{\mathrm{wc}} e^2 r J_{\mathrm{site}} = \frac{1}{4\cdot 6^r} \le \frac{1}{144} < 1$), summing over $k\ge 2$ converges geometrically and yields $|\mathrm{III}(\mathcal{O}(t))| \le C_2(\beta)\sqrt{N}$ for all $N\ge 4S_0/C_0$. For $1\le N < 4S_0/C_0$, $|\mathrm{III}(\mathcal{O}(t))|\le 2\|\mathcal{B}_\beta\|_\infty \le c(\beta)N \le c'(\beta)\sqrt{N}$ is absorbed into $C_2(\beta)$, so the bound holds for all $N\ge 1$ and $t\in\mathbb{R}$.

Finally, by Cauchy--Schwarz applied to~\eqref{eq:II-def},
\[
  |\mathrm{II}(\mathcal{O}(t))|
  \le 2\sqrt{\mathrm{I}(\mathcal{O}(t))}\left(\frac{1}{8}\sum_{x,\alpha}\|Y_{x,\alpha}\|_\HS^2\right)^{1/2}
  = 2\sqrt{\mathrm{I}(\mathcal{O}(t))}\,\sqrt{\langle\mathcal{B}_\beta\rangle_\beta + \mathrm{III}(\mathcal{O}(t))}.
\]
Since $\|\mathcal{B}_\beta\|_\infty \le N\cdot\frac{3}{8}(\frac{q}{1-q})^2$ ($q=\beta/(2\kappa)<1$), we have $\langle\mathcal{B}_\beta\rangle_\beta + |\mathrm{III}(\mathcal{O}(t))| \le c_B(\beta) N$ for all $N\ge 1$, whence $|\mathrm{II}(\mathcal{O}(t))| \le C_1'(\beta)\sqrt{N\,\mathrm{I}(\mathcal{O}(t))}$, proving~\eqref{eq:entropic-decoupling}.
\end{proof}

\noindent\emph{Remark.}
For fixed-locality observables, the Gaussian concentration theorem of Kuwahara and Saito~\cite{Kuwahara2020a} (see also~\cite{Anshu2021}) extends high-temperature cumulant bounds up to the cluster-expansion threshold $\beta_c$. By proving Lemma~\ref{lem:kp-cumulant} directly on the normalized trace, Proposition~\ref{prop:entropic}(2) remains self-contained for all $\beta\in[0,\beta_{\mathrm{wc}}]$ without external $k_r\to\infty$ hypotheses.

\begin{corollary}[Krylov--Lieb--Robinson Bound]\label{cor:klr}
Suppose in addition to the hypotheses of Proposition~\ref{prop:entropic} that $H$ is defined on a $d$-dimensional lattice $\Lambda\subset\mathbb{Z}^d$ with finite-range interactions ($\diam(X)\le R$) and Lieb--Robinson velocity $v_{\mathrm{LR}}$~\cite{Lieb1972,Hastings2006,Nachtergaele2006}. Then for every $\beta\in[0,\beta_{\mathrm{wc}}]$ and every $t\ge 0$, there exist $N$-independent constants $C_{\mathrm{LR}}(\beta), C'_{\mathrm{LR}}(\beta) > 0$ such that
\begin{equation}\label{eq:klr-bound}
  \mathcal{S}_\beta(\mathcal{O}(t))
  \le C_{\mathrm{LR}}(\beta)(v_{\mathrm{LR}}t + 1)^d
    + C'_{\mathrm{LR}}(\beta)\sqrt{N(v_{\mathrm{LR}}t + 1)^d}
    + C_2(\beta)\sqrt{N}.
\end{equation}
\end{corollary}

\begin{proof}
For any site $x\in\Lambda$ at distance $\ell := \dist(x,\supp(\mathcal{O})) > 2v_{\mathrm{LR}}t + 2R$, truncate the causal-sequence expansion of $\sigma_x^\alpha(-i\beta/4)$ from Lemma~\ref{lem:nested}(1) at order $m_0 := \lfloor \ell/(2R)\rfloor$: the strictly local part $W_{x,\alpha}^{(\le m_0)}$ is supported within distance $\ell/2$ of $x$ (hence at distance at least $\ell/2 > v_{\mathrm{LR}}t$ from $\supp(\mathcal{O})$), while the remainder has norm at most $\sum_{m>m_0}(\beta/(4\kappa))^m = O(e^{-c\ell})$. By Cauchy--Schwarz, $\|B\|_{\beta,W}^2 = \Tr(\rho_\beta^{1/2}B^\dagger\rho_\beta^{1/2}B) \le \Tr(\rho_\beta B^\dagger B) \le \|B\|_\infty^2$ for every $B\in\Hop$. Invoking the Lieb--Robinson bound on $[W_{x,\alpha}^{(\le m_0)}, \mathcal{O}(t)]$ (absorbing the polynomial support-size prefactor $|\supp(W_{x,\alpha}^{(\le m_0)})| \le C\ell^d$ into the exponential decay $e^{-\mu(\ell/2 - v_{\mathrm{LR}}t)}$) shows that $\|[\sigma_x^\alpha(-i\beta/4),\mathcal{O}(t)]\|_{\beta,W}$ decays exponentially in $\ell - v_{\mathrm{LR}}t$, so $\mathrm{I}(\mathcal{O}(t)) \le C(v_{\mathrm{LR}}t+1)^d$. Substituting this into~\eqref{eq:entropic-decoupling} yields~\eqref{eq:klr-bound}.
\end{proof}

Proposition~\ref{prop:entropic}(2) reduces the finite-temperature MSS bound~\eqref{eq:intro-uogh} at high temperature $\beta\in[0,\beta_{\mathrm{wc}}]$ to a single global operator inequality: dominating $\mathrm{I}(A)$ by $(A|\mathcal{Q}^{(\beta)}|A)_\beta^W$. Together with Proposition~\ref{prop:glauber}, this motivates the main technical construction of the paper, which we now develop in Section~\ref{sec:highT}.

\section{Depolarized $2$-Design Cluster Expansion and Proof of Theorem~\ref{thm:main-highT}}\label{sec:highT}

In this section we prove Theorem~\ref{thm:main-highT}. The central challenge is to establish a global reverse-Poincar\'e comparison between the sitewise Hilbert--Schmidt projection sum $\sum_{x=1}^N \|\mathcal{S}_x(A)\|_{\beta,W}^2$ (and the imaginary-time shifted commutator sum $\mathrm{I}(A)$) and the Wightman filtration superoperator $(A|\mathcal{Q}^{(\beta)}|A)_\beta^W$ for arbitrary non-commuting $r$-local qubit Hamiltonians in any spatial dimension $d\ge 1$.

Our strategy adapts the high-temperature separability cluster expansion of Bakshi, Liu, Moitra, and Tang~\cite{Bakshi2024} and Kim, Yang, and Lee~\cite{Kim2026}, with one essential modification. In~\cite{Bakshi2024,Kim2026}, $\rho_\beta$ is expressed as a convex combination of tensor products of rank-$1$ single-qubit projectors. While rank-$1$ projectors suffice for separability, they are singular and cannot induce a non-degenerate inner product on $\Hop$. By depolarizing the six Pauli eigenprojectors to strictly positive operators with uniform condition number $3$, we obtain a convex decomposition of $\rho_{\beta/2}$ whose induced Wightman inner product is a positive average of well-conditioned product inner products. We organize the construction into four subsections: the single-site depolarized frame, the polymer cluster expansion and convergence lemma, the product-space reverse-Poincar\'e comparison, and the causal-sequence root-reversal estimate.

\subsection{Depolarized Single-Site Projective $2$-Design Frame}\label{subsec:2design-frame}

On a single qubit $\mathcal{H}_x = \mathbb{C}^2$, let $\{P_{s_x}\}_{s_x=1}^6$ denote the six rank-$1$ eigenprojectors of the Pauli operators $\sigma_x^1,\sigma_x^2,\sigma_x^3$:
\[
  P_{2\alpha-1} := \frac{1}{2}(I_x + \sigma_x^\alpha), \qquad
  P_{2\alpha} := \frac{1}{2}(I_x - \sigma_x^\alpha) \qquad (\alpha\in\{1,2,3\}).
\]
We define the \emph{depolarized single-site operators}
\begin{equation}\label{eq:Ptilde-def}
  \widetilde{P}_{s_x} := \frac{1}{2} P_{s_x} + \frac{1}{4} I_x
  = \frac{1}{2} I_x \pm \frac{1}{4}\sigma_x^\alpha \qquad (s_x\in\{1,\dots,6\}),
\end{equation}
together with the associated \emph{dual frame operators}
\begin{equation}\label{eq:Ftilde-def}
  \widetilde{F}_{s_x} := 24\,\widetilde{P}_{s_x} - 11\,I_x
  = I_x \pm 6\,\sigma_x^\alpha \qquad (s_x\in\{1,\dots,6\}).
\end{equation}
The basic algebraic and spectral properties of this frame are summarized in the following lemma.

\begin{lemma}[Depolarized Single-Site Frame Identity]\label{lem:frame}
For each site $x\in\Lambda$, the operators $\{\widetilde{P}_{s_x},\widetilde{F}_{s_x}\}_{s_x=1}^6$ satisfy:
\begin{enumerate}
  \setlength{\itemsep}{2pt}\setlength{\parskip}{0pt}
  \item[\emph{(1)}] \emph{Positivity and uniform conditioning:} Every $\widetilde{P}_{s_x}$ is a Hermitian density matrix ($\Tr(\widetilde{P}_{s_x}) = 1$) with eigenvalues $\frac{3}{4}$ and $\frac{1}{4}$; in particular, $\frac{1}{4}I_x \preceq \widetilde{P}_{s_x} \preceq \frac{3}{4}I_x$ and $\frac{1}{6}\sum_{s_x=1}^6 \widetilde{P}_{s_x} = \frac{1}{2}I_x$.
  \item[\emph{(2)}] \emph{Dual reconstruction:} Every operator $B\in M_2(\mathbb{C})$ admits the representation
  \begin{equation}\label{eq:single-frame}
    \frac{1}{2}\,B = \frac{1}{6}\sum_{s_x=1}^6 \left(\frac{1}{2}\Tr\bigl(B\,\widetilde{F}_{s_x}\bigr)\right)\widetilde{P}_{s_x},
  \end{equation}
  and if $B = B^\dagger$ then $\frac{1}{2}\Tr(B\,\widetilde{F}_{s_x})\in\mathbb{R}$ for every $s_x\in\{1,\dots,6\}$.
  \item[\emph{(3)}] \emph{Dual norm bound:} The dual operator $\frac{1}{2}\widetilde{F}_{s_x}$ has eigenvalues $\frac{7}{2}$ and $-\frac{5}{2}$, hence trace norm $\|\frac{1}{2}\widetilde{F}_{s_x}\|_1 = 6$. Consequently,
  \begin{equation}\label{eq:dual-bound}
    \left|\frac{1}{2}\Tr\bigl(B\,\widetilde{F}_{s_x}\bigr)\right| \le 6\,\|B\|_\infty \qquad \forall B\in M_2(\mathbb{C}),\; \forall s_x\in\{1,\dots,6\}.
  \end{equation}
\end{enumerate}
\end{lemma}

\begin{proof}
Part~(1) is immediate from~\eqref{eq:Ptilde-def}, since the eigenvalues of $\sigma_x^\alpha$ are $\pm 1$, giving eigenvalues $\frac{1}{2}\pm\frac{1}{4}\in\{\frac{3}{4},\frac{1}{4}\}$, trace $\frac{3}{4}+\frac{1}{4}=1$, and average $\frac{1}{6}\sum_{s_x=1}^6 \widetilde{P}_{s_x} = \frac{1}{2}I_x$.

For part~(2), it suffices by linearity to check~\eqref{eq:single-frame} on the Pauli basis $\{I_x,\sigma_x^1,\sigma_x^2,\sigma_x^3\}$. For $B = I_x$, $\frac{1}{2}\Tr(I_x\widetilde{F}_{s_x}) = \frac{1}{2}\Tr(I_x) = 1$, so the right-hand side is $\frac{1}{6}\sum_{s_x=1}^6 \widetilde{P}_{s_x} = \frac{1}{2}I_x$. For $B = \sigma_x^\beta$ ($\beta\in\{1,2,3\}$), we have $\frac{1}{2}\Tr(\sigma_x^\beta\widetilde{F}_{s_x}) = \pm 6\delta_{\alpha\beta}$, so only the two indices $s_x\in\{2\beta-1, 2\beta\}$ contribute:
\[
  \frac{1}{6}\left[6\left(\frac{1}{2}I_x + \frac{1}{4}\sigma_x^\beta\right) - 6\left(\frac{1}{2}I_x - \frac{1}{4}\sigma_x^\beta\right)\right]
  = \frac{1}{2}\sigma_x^\beta.
\]
Since $\widetilde{F}_{s_x} = \widetilde{F}_{s_x}^\dagger$, the trace $\frac{1}{2}\Tr(B\widetilde{F}_{s_x})$ is real whenever $B = B^\dagger$.

For part~(3), the eigenvalues of $\frac{1}{2}\widetilde{F}_{s_x} = \frac{1}{2}I_x \pm 3\sigma_x^\alpha$ are $\frac{1}{2}+3 = \frac{7}{2}$ and $\frac{1}{2}-3 = -\frac{5}{2}$, so $\|\frac{1}{2}\widetilde{F}_{s_x}\|_1 = \frac{7}{2}+\frac{5}{2} = 6$. H\"older's inequality for Schatten norms gives $|\frac{1}{2}\Tr(B\widetilde{F}_{s_x})| \le \|B\|_\infty\|\frac{1}{2}\widetilde{F}_{s_x}\|_1 = 6\|B\|_\infty$.
\end{proof}

Tensoring~\eqref{eq:single-frame} across all sites $x\in\Lambda$, set $\widetilde{P}_{\mathbf{s}} := \bigotimes_{x\in\Lambda}\widetilde{P}_{s_x}$ and $\widetilde{F}_{\mathbf{s}} := \bigotimes_{x\in\Lambda}\widetilde{F}_{s_x}$ for $\mathbf{s} = (s_1,\dots,s_N)\in\{1,\dots,6\}^N$, and define the real-linear functional
\begin{equation}\label{eq:omegatilde-def}
  \widetilde{\omega}_{\mathbf{s}}(B) := \frac{1}{2^N}\Tr\bigl(B\,\widetilde{F}_{\mathbf{s}}\bigr) \qquad (B\in\Hop).
\end{equation}
Then every $B\in\Hop$ admits the global representation
\begin{equation}\label{eq:global-frame}
  \frac{1}{2^N}\,B = \frac{1}{6^N}\sum_{\mathbf{s}\in\{1,\dots,6\}^N} \widetilde{\omega}_{\mathbf{s}}(B)\,\widetilde{P}_{\mathbf{s}}.
\end{equation}
Moreover, if $B$ is supported on a subset $U\subseteq\Lambda$ (meaning $B = B_U\otimes I_{U^c}$), then $\frac{1}{2}\Tr_x(I_x\widetilde{F}_{s_x}) = 1$ for every $x\notin U$, so $\widetilde{\omega}_{\mathbf{s}}(B) = 2^{-|U|}\Tr_U(B_U\widetilde{F}_{\mathbf{s}_U})$ depends only on $\mathbf{s}_U$ and satisfies $|\widetilde{\omega}_{\mathbf{s}}(B)| \le 6^{|U|}\|B\|_\infty$. In particular, if $B_1,\dots,B_k$ have pairwise disjoint supports $U_1,\dots,U_k$, then $\widetilde{\omega}_{\mathbf{s}}$ factorizes multiplicatively:
\begin{equation}\label{eq:omegatilde-factor}
  \widetilde{\omega}_{\mathbf{s}}\!\left(\prod_{j=1}^k B_j\right) = \prod_{j=1}^k \widetilde{\omega}_{\mathbf{s}}(B_j).
\end{equation}

\subsection{Polymer Expansion and Positive Convex Separability of $\rho_{\beta/2}$}\label{subsec:polymer}

Applying~\eqref{eq:global-frame} to $B = e^{-\beta' H}$ with $\beta' := \beta/2$, we seek conditions under which $\widetilde{\omega}_{\mathbf{s}}(e^{-\beta' H}) > 0$ for every configuration $\mathbf{s}\in\{1,\dots,6\}^N$. Because $\Tr(\widetilde{P}_{\mathbf{s}}) = 1$, positivity of $\widetilde{\omega}_{\mathbf{s}}(e^{-\beta' H})$ immediately expresses $\rho_{\beta/2} = e^{-\beta' H}/Z(\beta/2)$ as a convex mixture of the strictly positive product states $\widetilde{P}_{\mathbf{s}}$.

\begin{lemma}[Depolarized $2$-Design Convex Decomposition]\label{lem:2design-sep}
Let $H = \sum_{X\in\mathcal{E}} h_X$ be any $r$-local Hamiltonian ($r\ge 2$) on $N$ qubits with per-site interaction strength $J_{\mathrm{site}} < \infty$, and recall $\beta_{\mathrm{wc}} = (4e^2 r 6^r J_{\mathrm{site}})^{-1}$. Then for every $\beta\in[0,\beta_{\mathrm{wc}}]$ and every $\mathbf{s}\in\{1,\dots,6\}^N$, we have $\widetilde{\omega}_{\mathbf{s}}(e^{-\beta H/2}) > 0$. Consequently,
\begin{equation}\label{eq:rho-convex}
  \rho_{\beta/2} = \sum_{\mathbf{s}\in\{1,\dots,6\}^N} \widetilde{\pi}(\mathbf{s})\,\widetilde{P}_{\mathbf{s}}, \qquad
  \widetilde{\pi}(\mathbf{s}) := \frac{2^N\,\widetilde{\omega}_{\mathbf{s}}(e^{-\beta H/2})}{6^N\,Z(\beta/2)} > 0, \qquad
  \sum_{\mathbf{s}}\widetilde{\pi}(\mathbf{s}) = 1.
\end{equation}
\end{lemma}

\begin{proof}
Set $\beta' := \beta/2$ and expand $e^{-\beta' H}$ as a power series in $H = \sum_{X\in\mathcal{E}} h_X$:
\[
  e^{-\beta' H} = \sum_{k=0}^\infty \frac{(-\beta')^k}{k!}\sum_{Y_1,\dots,Y_k\in\mathcal{E}} h_{Y_1}\cdots h_{Y_k}.
\]
For each ordered tuple $(Y_1,\dots,Y_k)\in\mathcal{E}^k$, let $\mathcal{I}(Y_1,\dots,Y_k)$ be the set of distinct hyperedges appearing in $(Y_1,\dots,Y_k)$, viewed as a hypergraph on $\Lambda$ where two hyperedges $X,X'$ are adjacent if $X\cap X'\ne\emptyset$. Decompose $\mathcal{I}(Y_1,\dots,Y_k)$ into its connected components (which we call \emph{polymers}) $\mathcal{G} = \{\gamma_1,\dots,\gamma_p\}$. Each polymer $\gamma$ is a nonempty finite set of hyperedges in $\mathcal{E}$ whose intersection graph is connected, with support $U_\gamma := \bigcup_{X\in\gamma} X$ and size $|\gamma|$ (number of distinct hyperedges in $\gamma$). Two polymers $\gamma,\gamma'$ are \emph{compatible} ($\gamma\sim\gamma'$) if $U_\gamma\cap U_{\gamma'} = \emptyset$.

Because operators supported on disjoint subsets commute, grouping all tuples $(Y_1,\dots,Y_k)$ according to their set of connected components $\mathcal{G}$ and the multiplicities $n_X \ge 1$ of each $X\in\gamma$ gives the polymer representation
\begin{equation}\label{eq:poly-op}
  e^{-\beta' H} = \sum_{\substack{\mathcal{G}\text{ pairwise}\\\text{compatible}}} \prod_{\gamma\in\mathcal{G}} Q_\gamma,
\end{equation}
where for each polymer $\gamma$ we define
\[
  Q_\gamma := \sum_{\{n_X\ge 1\}_{X\in\gamma}} \frac{(-\beta')^{\sum_{X\in\gamma} n_X}}{\bigl(\sum_{X\in\gamma} n_X\bigr)!}
  \sum_{\substack{(Y_1,\dots,Y_k)\\ \#\{j:\,Y_j=X\}=n_X}} h_{Y_1}\cdots h_{Y_k}.
\]
We verify two properties of $Q_\gamma$:
\begin{enumerate}
  \setlength{\itemsep}{2pt}\setlength{\parskip}{0pt}
  \item[\emph{(i)}] \emph{Hermiticity.} Because each $h_X$ is Hermitian, $(h_{Y_1}\cdots h_{Y_k})^\dagger = h_{Y_k}\cdots h_{Y_1}$. Reversing the order of a tuple $(Y_1,\dots,Y_k)\mapsto(Y_k,\dots,Y_1)$ preserves the multiplicity $n_X$ of every $X\in\gamma$ and is a bijection on the inner sum. Since $(-\beta')^k/k!\in\mathbb{R}$, we conclude $Q_\gamma^\dagger = Q_\gamma$.
  \item[\emph{(ii)}] \emph{Operator-norm bound.} For fixed multiplicities $\{n_X\}_{X\in\gamma}$ with $k = \sum_{X\in\gamma} n_X$, the number of tuples $(Y_1,\dots,Y_k)$ having multiplicity $n_X$ for each $X\in\gamma$ is the multinomial coefficient $k!/\prod_{X\in\gamma} n_X!$. By the triangle inequality and submultiplicativity of $\|\cdot\|_\infty$,
  \[
    \|Q_\gamma\|_\infty \le \sum_{\{n_X\ge 1\}_{X\in\gamma}} \prod_{X\in\gamma}\frac{(\beta'\|h_X\|_\infty)^{n_X}}{n_X!}
    = \prod_{X\in\gamma}\bigl(e^{\beta'\|h_X\|_\infty} - 1\bigr).
  \]
\end{enumerate}

Applying $\widetilde{\omega}_{\mathbf{s}}$ to~\eqref{eq:poly-op} and using the multiplicative factorization~\eqref{eq:omegatilde-factor} across pairwise disjoint supports $U_\gamma$ yields the abstract polymer partition function
\[
  \widetilde{\omega}_{\mathbf{s}}\bigl(e^{-\beta' H}\bigr)
  = \sum_{\substack{\mathcal{G}\text{ pairwise}\\\text{compatible}}} \prod_{\gamma\in\mathcal{G}} \widetilde{\omega}_{\mathbf{s}}(Q_\gamma).
\]
By Hermiticity of $Q_\gamma$ and Lemma~\ref{lem:frame}(2), every polymer activity $\widetilde{\omega}_{\mathbf{s}}(Q_\gamma)$ is real-valued. Moreover, since $|U_\gamma| \le r|\gamma|$, Lemma~\ref{lem:frame}(3) gives
\[
  \bigl|\widetilde{\omega}_{\mathbf{s}}(Q_\gamma)\bigr|
  \le 6^{|U_\gamma|}\|Q_\gamma\|_\infty
  \le \prod_{X\in\gamma} 6^r\bigl(e^{\beta'\|h_X\|_\infty} - 1\bigr) =: z_\gamma.
\]

We now verify the Koteck\'y--Preiss convergence criterion~\cite{Kotecky1986,Ueltschi2004} with weight $a_\gamma := |\gamma|$. Set $v_X := e\cdot 6^r(e^{\beta'\|h_X\|_\infty}-1)$ for each $X\in\mathcal{E}$, so that $z_\gamma e^{a_\gamma} = \prod_{X\in\gamma} v_X$. Using $e^u - 1 \le u e^u$ for $u\ge 0$ and $\|h_X\|_\infty \le J_{\mathrm{site}}$, we bound the per-site activity sum by
\[
  \sup_{x\in\Lambda}\sum_{X\ni x} v_X
  \le \sup_{x\in\Lambda}\sum_{X\ni x} e\cdot 6^r \beta'\|h_X\|_\infty e^{\beta' J_{\mathrm{site}}}
  \le \frac{\beta}{2}\,J_{\mathrm{site}}\,6^r\,e^{1+\frac{\beta}{2}J_{\mathrm{site}}} =: \eta_0.
\]
For $\beta\le\beta_{\mathrm{wc}} = (4e^2 r 6^r J_{\mathrm{site}})^{-1}$, we have $\frac{\beta}{2}J_{\mathrm{site}} \le (8e^2 r 6^r)^{-1} < \frac{1}{100} < 1$, whence $e^{1+\frac{\beta}{2}J_{\mathrm{site}}} < e^2$ and
\[
  \eta_0 < \frac{1}{8e^2 r 6^r}\cdot 6^r e^2 = \frac{1}{8r}.
\]
By Lemma~\ref{lem:app-planetree} (proved in Appendix~\ref{app:trees}), summing $\prod_{X\in\gamma} v_X$ over all connected polymers $\gamma$ of cardinality $|\gamma|=m$ whose support $U_\gamma$ contains a fixed site $x\in\Lambda$ is bounded by $\eta_0(4r\eta_0)^{m-1}$. Summing over $m\ge 1$ yields
\[
  \sum_{\gamma:\, x\in U_\gamma} z_\gamma e^{a_\gamma}
  \le \sum_{m=1}^\infty \eta_0(4r\eta_0)^{m-1}
  = \frac{\eta_0}{1 - 4r\eta_0}
  < \frac{1/(8r)}{1 - 1/2} = \frac{1}{4r}.
\]
Therefore, for any fixed polymer $\gamma$, summing over all sites $x\in U_\gamma$ (using $|U_\gamma|\le r|\gamma| = r a_\gamma$) gives
\[
  \sum_{\gamma':\, U_{\gamma'}\cap U_\gamma\ne\emptyset} \bigl|\widetilde{\omega}_{\mathbf{s}}(Q_{\gamma'})\bigr|\,e^{a_{\gamma'}}
  \le \sum_{x\in U_\gamma}\sum_{\gamma':\, x\in U_{\gamma'}} z_{\gamma'}e^{a_{\gamma'}}
  < |U_\gamma|\cdot\frac{1}{4r} \le \frac{1}{4}\,a_\gamma < a_\gamma.
\]
By the Koteck\'y--Preiss theorem~\cite{Kotecky1986,Ueltschi2004}, the cluster expansion for $\ln\widetilde{\omega}_{\mathbf{s}}(e^{-\beta' H})$ converges absolutely to a finite real number (since all activities $\widetilde{\omega}_{\mathbf{s}}(Q_\gamma)$ are real). Exponentiating yields $\widetilde{\omega}_{\mathbf{s}}(e^{-\beta' H}) = \exp(\ln\widetilde{\omega}_{\mathbf{s}}(e^{-\beta' H})) > 0$ for every $\mathbf{s}\in\{1,\dots,6\}^N$. Dividing~\eqref{eq:global-frame} (at $B = e^{-\beta' H}$) by $Z(\beta/2) = \Tr(e^{-\beta' H})$ and using $\Tr(\widetilde{P}_{\mathbf{s}})=1$ gives~\eqref{eq:rho-convex}.
\end{proof}

\subsection{Product-Space Reverse-Poincar\'e Comparison and Proof of Theorem~\ref{thm:main-highT}(a)}\label{subsec:highT-a}

We now show how the convex decomposition~\eqref{eq:rho-convex} converts the Wightman norm into a positive average of uniformly well-conditioned product inner products, yielding the global quantum reverse-Poincar\'e inequality~\eqref{eq:main-rp}.

\begin{lemma}[Product-Space Conditioning and Projection Comparison]\label{lem:prod-comp}
Fix any two configurations $\mathbf{s},\mathbf{s}'\in\{1,\dots,6\}^N$, and define the product inner product on $\Hop$ by
\begin{equation}\label{eq:ss-inner}
  \langle A, B\rangle_{\mathbf{s},\mathbf{s}'} := \Tr\bigl(\widetilde{P}_{\mathbf{s}} A^\dagger \widetilde{P}_{\mathbf{s}'} B\bigr),
  \qquad \|A\|_{\mathbf{s},\mathbf{s}'} := \langle A, A\rangle_{\mathbf{s},\mathbf{s}'}^{1/2}.
\end{equation}
Let $P_{x,0}^{(\mathbf{s},\mathbf{s}')}:\Hop\to\mathcal{V}_{(x^c)}$ and $\Pi_k^{(\mathbf{s},\mathbf{s}')}:\Hop\to\mathcal{V}_k$ denote the orthogonal projectors with respect to $\langle\cdot,\cdot\rangle_{\mathbf{s},\mathbf{s}'}$, and for any subset $U\subseteq\Lambda$ let $P_{U,0}^{(\mathbf{s},\mathbf{s}')}:\Hop\to\mathcal{V}_{(U^c)}$ denote the $\langle\cdot,\cdot\rangle_{\mathbf{s},\mathbf{s}'}$-orthogonal projector onto $\mathcal{V}_{(U^c)}$. Then for every $A\in\Hop$:
\begin{enumerate}
  \setlength{\itemsep}{2pt}\setlength{\parskip}{0pt}
  \item[\emph{(1)}] For every subset $U\subseteq\Lambda$,
  \begin{equation}\label{eq:U-cond}
    \bigl\|(\mathbf{1}-\mathcal{E}_U)A\bigr\|_{\mathbf{s},\mathbf{s}'}^2
    \le 9^{|U|}\bigl\|(\mathbf{1}-P_{U,0}^{(\mathbf{s},\mathbf{s}')})A\bigr\|_{\mathbf{s},\mathbf{s}'}^2
    \le 9^{|U|}\sum_{x\in U}\bigl\|(\mathbf{1}-P_{x,0}^{(\mathbf{s},\mathbf{s}')})A\bigr\|_{\mathbf{s},\mathbf{s}'}^2.
  \end{equation}
  In particular, for a single site $U=\{x\}$, $\|\mathcal{S}_x(A)\|_{\mathbf{s},\mathbf{s}'}^2 \le 9\|(\mathbf{1}-P_{x,0}^{(\mathbf{s},\mathbf{s}')})A\|_{\mathbf{s},\mathbf{s}'}^2$.
  \item[\emph{(2)}] The sitewise projection sum coincides with the product filtration superoperator and is bounded by the Wightman filtration projections:
  \begin{equation}\label{eq:prod-filt}
    \sum_{x=1}^N \bigl\|(\mathbf{1}-P_{x,0}^{(\mathbf{s},\mathbf{s}')})A\bigr\|_{\mathbf{s},\mathbf{s}'}^2
    = \sum_{k=0}^{N-1}\bigl\|(\mathbf{1}-\Pi_k^{(\mathbf{s},\mathbf{s}')})A\bigr\|_{\mathbf{s},\mathbf{s}'}^2
    \le \sum_{k=0}^{N-1}\bigl\|(\mathbf{1}-\Pi_k^{(\beta)})A\bigr\|_{\mathbf{s},\mathbf{s}'}^2.
  \end{equation}
\end{enumerate}
\end{lemma}

\begin{proof}
Because $\widetilde{P}_{\mathbf{s}} = \bigotimes_{x=1}^N\widetilde{P}_{s_x}$ and $\widetilde{P}_{\mathbf{s}'} = \bigotimes_{x=1}^N\widetilde{P}_{s'_x}$, the inner product $\langle\cdot,\cdot\rangle_{\mathbf{s},\mathbf{s}'}$ factorizes as the tensor product of the single-site inner products $\langle a,b\rangle_{s_x,s'_x} := \Tr(\widetilde{P}_{s_x} a^\dagger \widetilde{P}_{s'_x} b)$ on $M_2(\mathbb{C})$.

\emph{Part~(1).} On a subset $U\subseteq\Lambda$, write $\widetilde{P}_{\mathbf{s}_U} := \bigotimes_{x\in U}\widetilde{P}_{s_x}$. By Lemma~\ref{lem:frame}(1), $\frac{1}{4^{|U|}}I_U \preceq \widetilde{P}_{\mathbf{s}_U}, \widetilde{P}_{\mathbf{s}'_U} \preceq \frac{3^{|U|}}{4^{|U|}}I_U$. Therefore, for any operator $Y$ on $\bigotimes_{x\in U}\mathcal{H}_x$,
\[
  \frac{1}{16^{|U|}}\|Y\|_\HS^2
  \le \Tr\bigl(\widetilde{P}_{\mathbf{s}_U} Y^\dagger \widetilde{P}_{\mathbf{s}'_U} Y\bigr)
  \le \frac{9^{|U|}}{16^{|U|}}\|Y\|_\HS^2.
\]
Apply the upper bound to $Y = (\mathbf{1}-\mathcal{E}_U)B$ and the lower bound to $Y = B - c I_U$ for any scalar $c\in\mathbb{C}$; since $(\mathbf{1}-\mathcal{E}_U)$ is the Hilbert--Schmidt orthogonal projector onto traceless operators on $U$, we have $\|(\mathbf{1}-\mathcal{E}_U)B\|_\HS^2 \le \|B - cI_U\|_\HS^2$, whence
\[
  \bigl\|(\mathbf{1}-\mathcal{E}_U)B\bigr\|_{\mathbf{s}_U,\mathbf{s}'_U}^2
  \le \frac{9^{|U|}}{16^{|U|}}\bigl\|(\mathbf{1}-\mathcal{E}_U)B\bigr\|_\HS^2
  \le \frac{9^{|U|}}{16^{|U|}}\|B - cI_U\|_\HS^2
  \le 9^{|U|}\|B - cI_U\|_{\mathbf{s}_U,\mathbf{s}'_U}^2.
\]
Taking the infimum over $c\in\mathbb{C}$ gives $\|(\mathbf{1}-\mathcal{E}_U)B\|_{\mathbf{s}_U,\mathbf{s}'_U}^2 \le 9^{|U|}\|(\mathbf{1}-P_{U,0}^{(\mathbf{s}_U,\mathbf{s}'_U)})B\|_{\mathbf{s}_U,\mathbf{s}'_U}^2$. Tensoring with the identity on $U^c$ (with respect to $\langle\cdot,\cdot\rangle_{\mathbf{s}_{U^c},\mathbf{s}'_{U^c}}$) yields the first inequality in~\eqref{eq:U-cond}.

For the second inequality in~\eqref{eq:U-cond} and part~(2), choose for each site $x\in\Lambda$ an orthonormal basis $\{e_{x,0}, e_{x,1}, e_{x,2}, e_{x,3}\}$ of $(M_2(\mathbb{C}), \langle\cdot,\cdot\rangle_{s_x,s'_x})$ with $e_{x,0} \propto I_x$. The tensor products $\mathbf{e}_{\mathbf{a}} := \bigotimes_{x=1}^N e_{x,a_x}$ ($\mathbf{a}\in\{0,1,2,3\}^N$) form a $\langle\cdot,\cdot\rangle_{\mathbf{s},\mathbf{s}'}$-orthonormal basis of $\Hop$ that simultaneously diagonalizes $P_{x,0}^{(\mathbf{s},\mathbf{s}')}$, $P_{U,0}^{(\mathbf{s},\mathbf{s}')}$, and $\Pi_k^{(\mathbf{s},\mathbf{s}')}$. Expanding $A = \sum_{\mathbf{a}} c_{\mathbf{a}}\mathbf{e}_{\mathbf{a}}$ and writing $|\mathbf{a}| := \#\{x\in\Lambda : a_x\ne 0\}$, we have
\begin{align*}
  \bigl\|(\mathbf{1}-P_{U,0}^{(\mathbf{s},\mathbf{s}')})A\bigr\|_{\mathbf{s},\mathbf{s}'}^2
  &= \sum_{\mathbf{a}} |c_{\mathbf{a}}|^2\,\mathbf{1}\bigl(\exists x\in U : a_x\ne 0\bigr) \\
  &\le \sum_{\mathbf{a}} |c_{\mathbf{a}}|^2\sum_{x\in U}\mathbf{1}(a_x\ne 0)
   = \sum_{x\in U}\bigl\|(\mathbf{1}-P_{x,0}^{(\mathbf{s},\mathbf{s}')})A\bigr\|_{\mathbf{s},\mathbf{s}'}^2,
\end{align*}
establishing~\eqref{eq:U-cond}. Likewise,
\[
  \sum_{x=1}^N \bigl\|(\mathbf{1}-P_{x,0}^{(\mathbf{s},\mathbf{s}')})A\bigr\|_{\mathbf{s},\mathbf{s}'}^2
  = \sum_{\mathbf{a}} |\mathbf{a}|\,|c_{\mathbf{a}}|^2
  = \sum_{k=0}^{N-1}\sum_{|\mathbf{a}|>k}|c_{\mathbf{a}}|^2
  = \sum_{k=0}^{N-1}\bigl\|(\mathbf{1}-\Pi_k^{(\mathbf{s},\mathbf{s}')})A\bigr\|_{\mathbf{s},\mathbf{s}'}^2.
\]
Finally, since $\Pi_k^{(\beta)}A\in\mathcal{V}_k$ and $\Pi_k^{(\mathbf{s},\mathbf{s}')}A$ is the $\langle\cdot,\cdot\rangle_{\mathbf{s},\mathbf{s}'}$-best approximation to $A$ in $\mathcal{V}_k$, we have $\|(\mathbf{1}-\Pi_k^{(\mathbf{s},\mathbf{s}')})A\|_{\mathbf{s},\mathbf{s}'}^2 \le \|(\mathbf{1}-\Pi_k^{(\beta)})A\|_{\mathbf{s},\mathbf{s}'}^2$ for every $k$, proving~\eqref{eq:prod-filt}.
\end{proof}

\begin{proof}[Proof of Theorem~\ref{thm:main-highT}(a)]
Fix $\beta\in[0,\beta_{\mathrm{wc}}]$. By Lemma~\ref{lem:2design-sep}, $\rho_\beta^{1/2} = c_\beta\sum_{\mathbf{s}}\widetilde{\pi}(\mathbf{s})\widetilde{P}_{\mathbf{s}}$ where $c_\beta := Z(\beta/2)/Z(\beta)^{1/2} > 0$ and $\widetilde{\pi}(\mathbf{s}) > 0$. Expanding the Wightman norm squared yields the positive-weight mixture identity
\begin{equation}\label{eq:norm-mixture}
  \|X\|_{\beta,W}^2 = \Tr\bigl(\rho_\beta^{1/2} X^\dagger \rho_\beta^{1/2} X\bigr)
  = c_\beta^2 \sum_{\mathbf{s},\mathbf{s}'\in\{1,\dots,6\}^N} \widetilde{\pi}(\mathbf{s})\widetilde{\pi}(\mathbf{s}')\,\|X\|_{\mathbf{s},\mathbf{s}'}^2 \qquad \forall X\in\Hop.
\end{equation}
For each pair $(\mathbf{s},\mathbf{s}')$, Lemma~\ref{lem:prod-comp}(1)--(2) gives
\[
  \sum_{x=1}^N \|\mathcal{S}_x(A)\|_{\mathbf{s},\mathbf{s}'}^2
  \le 9\sum_{x=1}^N \bigl\|(\mathbf{1}-P_{x,0}^{(\mathbf{s},\mathbf{s}')})A\bigr\|_{\mathbf{s},\mathbf{s}'}^2
  \le 9\sum_{k=0}^{N-1}\bigl\|(\mathbf{1}-\Pi_k^{(\beta)})A\bigr\|_{\mathbf{s},\mathbf{s}'}^2.
\]
Multiplying both sides by the positive weight $c_\beta^2\widetilde{\pi}(\mathbf{s})\widetilde{\pi}(\mathbf{s}')$, summing over all $\mathbf{s},\mathbf{s}'$, and applying~\eqref{eq:norm-mixture} gives
\[
  \sum_{x=1}^N \|\mathcal{S}_x(A)\|_{\beta,W}^2
  \le 9\sum_{k=0}^{N-1}\bigl\|(\mathbf{1}-\Pi_k^{(\beta)})A\bigr\|_{\beta,W}^2
  = 9\bigl(A\big|\mathcal{Q}^{(\beta)}\big|A\bigr)_\beta^W.
\]
On the other hand, since $\mathcal{E}_x(A)\in\mathcal{V}_{(x^c)}$ and $P_{x,0}^{(\beta)}A$ is the $(\cdot|\cdot)_\beta^W$-best approximation to $A$ in $\mathcal{V}_{(x^c)}$, we have $(A|\widetilde{\mathcal{Q}}^{(\beta)}|A)_\beta^W = \sum_{x=1}^N \|(\mathbf{1}-P_{x,0}^{(\beta)})A\|_{\beta,W}^2 \le \sum_{x=1}^N \|\mathcal{S}_x(A)\|_{\beta,W}^2$, establishing~\eqref{eq:main-rp}.
\end{proof}

\subsection{Causal-Sequence Expansion, Root Reversal, and Proof of Theorem~\ref{thm:main-highT}(b)--(c)}\label{subsec:highT-bc}

To prove the global MSS-commutator bound $\mathrm{I}(A) \le 216(A|\mathcal{Q}^{(\beta)}|A)_\beta^W$, we combine the causal-sequence expansion of $\sigma_i^\alpha(-i\beta/4)$ with the subset conditioning estimate~\eqref{eq:U-cond} and a root-reversal counting bound derived from Cayley's spanning-tree formula (Appendix~\ref{app:trees}).

\begin{proof}[Proof of Theorem~\ref{thm:main-highT}(b)--(c)]
\emph{Part~(b): Wightman commutator $C_W(A)$.} Since $\beta \le \beta_{\mathrm{wc}} = (4e^2 r 6^r J_{\mathrm{site}})^{-1} < (r-1)^{-1}J_{\mathrm{site}}^{-1} = 2\kappa$, we have $1 - \frac{\beta}{4\kappa} > \frac{1}{2}$, so $(1-\frac{\beta}{4\kappa})^{-2} < 4$. Combining Proposition~\ref{prop:glauber} with Theorem~\ref{thm:main-highT}(a) immediately yields
\[
  C_W(A) \le \frac{3}{2}\cdot 4\,\bigl(A\big|\widetilde{\mathcal{Q}}^{(\beta)}\big|A\bigr)_\beta^W
  = 6\bigl(A\big|\widetilde{\mathcal{Q}}^{(\beta)}\big|A\bigr)_\beta^W
  \le 54\bigl(A\big|\mathcal{Q}^{(\beta)}\big|A\bigr)_\beta^W.
\]

\emph{Part~(b): Imaginary-time shifted commutator $\mathrm{I}(A)$.} For each site $i\in\Lambda$ and $\alpha\in\{1,2,3\}$, expand $\sigma_i^\alpha(-i\beta/4)$ into ordered causal sequences $\Gamma = (X_1,\dots,X_k)$ rooted at $i$ (meaning $X_1\ni i$ and $X_j\cap(\{i\}\cup X_1\cup\dots\cup X_{j-1})\ne\emptyset$ for $2\le j\le k$, with the empty sequence $\Gamma=\emptyset$ at $k=0$):
\[
  \sigma_i^\alpha(-i\beta/4) = \sum_{k=0}^\infty \sum_{\Gamma\in\mathcal{C}_k(i)} T_{i,\alpha,\Gamma},
  \qquad T_{i,\alpha,\Gamma} := \frac{(\beta/4)^k}{k!}\,\ad_{h_{X_k}}\cdots\ad_{h_{X_1}}(\sigma_i^\alpha).
\]
Each term $T_{i,\alpha,\Gamma}$ is supported on $U_\Gamma := \{i\}\cup X_1\cup\dots\cup X_k$, which has cardinality $|U_\Gamma| \le 1 + (r-1)k$. Because $T_{i,\alpha,\Gamma}$ commutes with $\mathcal{E}_{U_\Gamma}(A)\in\mathcal{V}_{(U_\Gamma^c)}$, we have $[T_{i,\alpha,\Gamma}, A] = [T_{i,\alpha,\Gamma}, (\mathbf{1}-\mathcal{E}_{U_\Gamma})A]$. Applying the imaginary-time multiplication bound from the proof of Lemma~\ref{lem:wightman-comm} gives
\[
  \bigl\|[T_{i,\alpha,\Gamma}, A]\bigr\|_{\beta,W}
  \le \Bigl(\bigl\|T_{i,\alpha,\Gamma}(i\beta/4)\bigr\|_\infty + \bigl\|T_{i,\alpha,\Gamma}(-i\beta/4)\bigr\|_\infty\Bigr)\,\bigl\|(\mathbf{1}-\mathcal{E}_{U_\Gamma})A\bigr\|_{\beta,W}.
\]
Because $e^{\mp\frac{\beta}{4}H}$ acts as an algebra automorphism,
\[
  T_{i,\alpha,\Gamma}(\pm i\beta/4)
  = \frac{(\beta/4)^k}{k!}\,\ad_{h_{X_k}(\pm i\beta/4)}\cdots\ad_{h_{X_1}(\pm i\beta/4)}\bigl(\sigma_i^\alpha(\pm i\beta/4)\bigr).
\]
For $\beta\le\beta_{\mathrm{wc}} < 2\kappa$, Lemma~\ref{lem:nested}(2)--(3) gives $\|\sigma_i^\alpha(\pm i\beta/4)\|_\infty \le 2$ and $\|h_X(\pm i\beta/4)\|_\infty \le 4\|h_X\|_\infty$. Using $\|[B,C]\|_\infty \le 2\|B\|_\infty\|C\|_\infty$ at each of the $k$ nested commutators yields
\[
  \bigl\|T_{i,\alpha,\Gamma}(\pm i\beta/4)\bigr\|_\infty
  \le \frac{(\beta/4)^k}{k!}\cdot 2^k \left(\prod_{j=1}^k 4\|h_{X_j}\|_\infty\right)\cdot 2
  = 2\,\frac{(2\beta)^k}{k!}\prod_{j=1}^k \|h_{X_j}\|_\infty,
\]
and therefore
\begin{equation}\label{eq:Tgamma-comm}
  \bigl\|[T_{i,\alpha,\Gamma}, A]\bigr\|_{\beta,W}
  \le 4\,\frac{(2\beta)^k}{k!}\left(\prod_{j=1}^k \|h_{X_j}\|_\infty\right)\,\bigl\|(\mathbf{1}-\mathcal{E}_{U_\Gamma})A\bigr\|_{\beta,W}.
\end{equation}

To bound $\|(\mathbf{1}-\mathcal{E}_{U_\Gamma})A\|_{\beta,W}$, average~\eqref{eq:U-cond} over $(\mathbf{s},\mathbf{s}')$ via~\eqref{eq:norm-mixture}. Defining the sitewise averaged projection norm
\[
  D_x(A)^2 := c_\beta^2 \sum_{\mathbf{s},\mathbf{s}'}\widetilde{\pi}(\mathbf{s})\widetilde{\pi}(\mathbf{s}')\,\bigl\|(\mathbf{1}-P_{x,0}^{(\mathbf{s},\mathbf{s}')})A\bigr\|_{\mathbf{s},\mathbf{s}'}^2 \qquad (x\in\Lambda),
\]
Lemma~\ref{lem:prod-comp}(2) and~\eqref{eq:norm-mixture} ensure that
\begin{equation}\label{eq:Dx-sum}
  \sum_{x=1}^N D_x(A)^2 \le \bigl(A\big|\mathcal{Q}^{(\beta)}\big|A\bigr)_\beta^W,
\end{equation}
while~\eqref{eq:U-cond} and $|U_\Gamma|\le 1+(r-1)k$ yield
\[
  \bigl\|(\mathbf{1}-\mathcal{E}_{U_\Gamma})A\bigr\|_{\beta,W}
  \le 3^{1+(r-1)k}\left(\sum_{x\in U_\Gamma} D_x(A)^2\right)^{1/2}.
\]
Define the nonnegative causal-sequence weight
\[
  w_{i,\Gamma} := 4\,\frac{(2\beta)^k}{k!}\,3^{1+(r-1)k}\prod_{j=1}^k \|h_{X_j}\|_\infty
  = 12\,\frac{(2\beta\cdot 3^{r-1})^k}{k!}\prod_{j=1}^k \|h_{X_j}\|_\infty.
\]
By the triangle inequality in $(\Hop,(\cdot|\cdot)_\beta^W)$ and the Cauchy--Schwarz inequality,
\begin{equation}\label{eq:I-CS}
  \bigl\|[\sigma_i^\alpha(-i\beta/4), A]\bigr\|_{\beta,W}^2
  \le \left(\sum_{k=0}^\infty\sum_{\Gamma\in\mathcal{C}_k(i)} w_{i,\Gamma}\right)
      \left(\sum_{k=0}^\infty\sum_{\Gamma\in\mathcal{C}_k(i)} w_{i,\Gamma}\sum_{x\in U_\Gamma} D_x(A)^2\right).
\end{equation}

We now bound both sums in~\eqref{eq:I-CS}. First, by the inductive argument of Lemma~\ref{lem:nested}(1) (with $s=1$, where $\prod_{j=0}^{k-1}(\frac{1}{r-1}+j)\le k!$), we have $\sum_{\Gamma\in\mathcal{C}_k(i)} \frac{1}{k!}\prod_{j=1}^k \|h_{X_j}\|_\infty \le ((r-1)J_{\mathrm{site}})^k$, whence
\[
  \sum_{k=0}^\infty\sum_{\Gamma\in\mathcal{C}_k(i)} w_{i,\Gamma}
  \le 12\sum_{k=0}^\infty \bigl(2\beta\cdot 3^{r-1}(r-1)J_{\mathrm{site}}\bigr)^k \le 24,
\]
because $\beta\le\beta_{\mathrm{wc}} = (4e^2 r 6^r J_{\mathrm{site}})^{-1} < (4er 3^{r-1} J_{\mathrm{site}})^{-1}$ ensures that the geometric ratio is less than $\frac{1}{2}$. Second, for every fixed site $x\in\Lambda$, Lemma~\ref{lem:app-cayley} (proved in Appendix~\ref{app:trees} via Cayley's spanning-tree formula) shows that reversing the root from $i$ to $x\in U_\Gamma$ yields
\[
  \sum_{i\in\Lambda}\sum_{\substack{\Gamma\in\mathcal{C}_k(i) \\ x\in U_\Gamma}} \frac{1}{k!}\prod_{j=1}^k \|h_{X_j}\|_\infty
  \le (er J_{\mathrm{site}})^k \qquad (k\ge 1),
\]
while at $k=0$ the only term is $i=x$ with weight $w_{x,\emptyset}=12$. Summing over $k\ge 0$ gives
\[
  \sum_{i\in\Lambda}\sum_{k=0}^\infty\sum_{\substack{\Gamma\in\mathcal{C}_k(i) \\ x\in U_\Gamma}} w_{i,\Gamma}
  \le 12 + 12\sum_{k=1}^\infty \bigl(2\beta\cdot 3^{r-1} e r J_{\mathrm{site}}\bigr)^k \le 24,
\]
again because $2\beta\cdot 3^{r-1} e r J_{\mathrm{site}} \le \frac{1}{2}$ for $\beta\le\beta_{\mathrm{wc}}$. Substituting both bounds into~\eqref{eq:I-CS}, summing over $i\in\Lambda$ and $\alpha\in\{1,2,3\}$ with prefactor $\frac{1}{8}$, and applying~\eqref{eq:Dx-sum} yields
\[
  \mathrm{I}(A) = \frac{1}{8}\sum_{i=1}^N\sum_{\alpha=1}^3 \bigl\|[\sigma_i^\alpha(-i\beta/4), A]\bigr\|_{\beta,W}^2
  \le \frac{3}{8}\cdot 24^2 \sum_{x=1}^N D_x(A)^2
  \le 216\bigl(A\big|\mathcal{Q}^{(\beta)}\big|A\bigr)_\beta^W,
\]
establishing~\eqref{eq:main-CW-I}. Combining~\eqref{eq:main-CW-I} with Proposition~\ref{prop:q-complexity}(2) gives~\eqref{eq:main-dyn-bounds}.

\emph{Part~(c).} For the Wightman-commutator OTOC at $\beta\in[0,\beta_{\mathrm{wc}}]$, combining~\eqref{eq:main-dyn-bounds} with Theorem~\ref{thm:main-tails-reduction}(a) gives $C_W(\mathcal{O}_N(t)) \le C e^{2\alpha_T^{(W)} t}$ uniformly in $N$, so $\lambda_{L,T}^{(W)} \le 2\alpha_T^{(W)}$, and Theorem~\ref{thm:main-spectral}(b) yields $2\alpha_T^{(W)} \le 2\pi/(\beta+2\kappa) < 2\pi/\beta$.

For the MSS-regulated OTOC at $\beta\in[0,\beta_{\mathrm{wc}}]$, note that $\beta_{\mathrm{wc}} \le (8e^2 r J_{\mathrm{site}})^{-1}$ and $\beta_{\mathrm{wc}} < (erJ_{\mathrm{site}})^{-1}$, so Proposition~\ref{prop:entropic}(2) applies. Substituting $\mathrm{I}(\mathcal{O}_N(t)) \le 216((r-1)K_\beta^W(t)+s_0)$ into~\eqref{eq:entropic-decoupling} gives~\eqref{eq:main-Sbeta-dyn}. Suppose that $\mathcal{S}_\beta(\mathcal{O}_N(t_N(\epsilon))) \ge c N^{1-\epsilon}$ along $t_N(\epsilon) = \frac{1-\epsilon}{\lambda_L}\ln N$ for arbitrarily small $\epsilon > 0$. By Theorem~\ref{thm:main-tails-reduction}(a), $K_\beta^W(t_N(\epsilon)) = O(N^{2\alpha_T^{(W)}(1-\epsilon)/\lambda_L})$, so~\eqref{eq:main-Sbeta-dyn} forces
\[
  1 - \epsilon \le \max\!\left((1-\epsilon)\frac{2\alpha_T^{(W)}}{\lambda_L},\; \frac{1}{2} + \frac{1}{2}(1-\epsilon)\frac{2\alpha_T^{(W)}}{\lambda_L},\; \frac{1}{2}\right).
\]
For every $\epsilon\in(0,\frac{1}{2})$, we have $1-\epsilon > \frac{1}{2}$, so the maximum is attained by one of the first two arguments; in both cases, letting $\epsilon\downarrow 0$ yields $1 \le 2\alpha_T^{(W)}/\lambda_L$, i.e., $\lambda_{L,T}^{(\MSS)} \le 2\alpha_T^{(W)}$. Combining with Theorem~\ref{thm:main-spectral}(b) proves~\eqref{eq:main-mss-resolved}.
\end{proof}

\section{All-Temperature Regimes: One-Dimensional Chains and Classical Ising Models}\label{sec:ising}

While Theorem~\ref{thm:main-highT} establishes the chaos--operator growth bound at high temperature $\beta\in[0,\beta_{\mathrm{wc}}]$ in arbitrary spatial dimension, two structurally distinct settings admit all-temperature bounds: one-dimensional finite-range quantum spin chains (via Araki's imaginary-time analyticity theorem~\cite{Araki1969}) and arbitrary $2$-local classical Ising models (via a discrete Hubbard--Stratonovich decomposition).

\subsection{One-Dimensional Finite-Range Quantum Spin Chains}\label{subsec:1d-allT}

Summing Lemma~\ref{lem:wightman-comm} over all sites $x\in\Lambda$ and $\alpha\in\{1,2,3\}$ yields an all-temperature comparison between $C_W(A)$ and the quantum Glauber generator $\widetilde{\mathcal{Q}}^{(\beta)}$ in one spatial dimension.

\begin{proposition}[All-Temperature One-Dimensional Commutator Bound]\label{prop:1D}
Let $H = \sum_{X\in\mathcal{E}} h_X$ be a finite-range $r$-local Hamiltonian on a one-dimensional spin-$1/2$ chain $\Lambda = \{1,\dots,N\}$ with $\sup_x\sum_{X\ni x}\|h_X\|_\infty \le J_{\mathrm{site}}$ and $\max_{X\in\mathcal{E}}\diam(X) \le R$. Then for every inverse temperature $\beta\in[0,\infty)$, there exists a finite constant $M_{1\mathrm{D}}(\beta) = M_{1\mathrm{D}}(\beta, r, R, J_{\mathrm{site}}) < \infty$, independent of $N$, such that
\begin{equation}\label{eq:1D-comm-bound}
  C_W(A) \le \frac{3}{2} M_{1\mathrm{D}}(\beta)^2 \bigl(A\big|\widetilde{\mathcal{Q}}^{(\beta)}\big|A\bigr)_\beta^W
  \le \frac{3}{2} M_{1\mathrm{D}}(\beta)^2 \sum_{x=1}^N \|\mathcal{S}_x(A)\|_{\beta,W}^2 \qquad \forall A\in\Hop.
\end{equation}
\end{proposition}

\begin{proof}
By Araki's theorem on the analyticity of imaginary-time evolution for one-dimensional finite-range quantum spin systems~\cite{Araki1969} (see also~\cite[Sec.~6.2]{Bratteli1997} and~\cite{Bluhm2025}), for every $\beta\ge 0$, every site $x\in\Lambda$, and every $\alpha\in\{1,2,3\}$, the imaginary-time evolved operator $\sigma_x^\alpha(i\beta/4)$ is uniformly bounded in operator norm by a finite constant $M_{1\mathrm{D}}(\beta)$ depending only on $\beta$, $r$, $R$, and $J_{\mathrm{site}}$, and independent of $N$. Squaring~\eqref{eq:single-site-comm} and summing over $x\in\{1,\dots,N\}$ and $\alpha\in\{1,2,3\}$ with the prefactor $1/8$ gives
\[
  C_W(A) = \frac{1}{8}\sum_{x=1}^N\sum_{\alpha=1}^3 \bigl\|[\sigma_x^\alpha, A]\bigr\|_{\beta,W}^2
  \le \frac{3}{2} M_{1\mathrm{D}}(\beta)^2 \sum_{x=1}^N \bigl\|(\mathbf{1}-P_{x,0}^{(\beta)})A\bigr\|_{\beta,W}^2.
\]
Since $(\mathbf{1}-P_{x,0}^{(\beta)})$ is an orthogonal projector in $(\Hop,(\cdot|\cdot)_\beta^W)$, the sum on the right equals $(A|\widetilde{\mathcal{Q}}^{(\beta)}|A)_\beta^W$, and the second inequality in~\eqref{eq:1D-comm-bound} follows from Lemma~\ref{lem:wightman-comm}.
\end{proof}

\subsection{All-Temperature Classical Ising Reverse-Poincar\'e Bound}\label{subsec:classical-ising}

Theorem~\ref{thm:main-highT}(a) established the quantum reverse-Poincar\'e inequality $\widetilde{\mathcal{Q}}^{(\beta)}\preceq 9\mathcal{Q}^{(\beta)}$ for arbitrary non-commuting $r$-local Hamiltonians at high temperature $\beta\le\beta_{\mathrm{wc}}$. It is natural to ask whether an analogous reverse-Poincar\'e inequality persists at arbitrarily low temperature $\beta\to\infty$, including across thermal phase transitions, in commuting models where the Gibbs state is a classical spin measure. In this section we show that for any $2$-local classical Ising model on an arbitrary graph, a discrete Hubbard--Stratonovich decomposition yields an all-temperature reverse-Poincar\'e bound with a closed-form factor controlled by the maximum single-site field.

Let $\mu_\beta$ be a classical probability measure on $\{-1,+1\}^N$ of the form
\begin{equation}\label{eq:ising-gibbs}
  \mu_\beta(z) = \frac{1}{Z(\beta)}\exp\!\left(\beta\sum_{1\le i<j\le N} J_{ij} z_i z_j + \beta\sum_{i=1}^N h_i z_i\right),
\end{equation}
where $J_{ij}, h_i\in\mathbb{R}$. For a function $f:\{-1,+1\}^N\to\mathbb{R}$, define the discrete derivative $\partial_i f(z) := \frac{1}{2}(f(z_i=+1,z_{-i}) - f(z_i=-1,z_{-i}))$, the Glauber Dirichlet form
\[
  \mathcal{E}_{\mu_\beta}(f,f) := \sum_{i=1}^N \E_{\mu_\beta}\bigl[\Var_{\mu_\beta}(f\mid z_{-i})\bigr]
  = \sum_{i=1}^N \E_{\mu_\beta}\bigl[\Var_{\mu_\beta}(z_i\mid z_{-i})\,(\partial_i f)^2\bigr],
\]
and the degree filtration $\mathcal{V}_k^{\mathrm{cl}} := \operatorname{span}\{\prod_{i\in S} z_i : |S|\le k\}$ with $L^2(\mu_\beta)$-orthogonal projectors $\Pi_k^{(\mu_\beta)}$ and filtration quadratic form
\[
  \mathcal{Q}_{\mu_\beta}(f,f) := \sum_{k=0}^{N-1}\bigl\|(\mathbf{1}-\Pi_k^{(\mu_\beta)})f\bigr\|_{L^2(\mu_\beta)}^2.
\]
Note that for diagonal operators $A = f(Z_1,\dots,Z_N)$ in a classical Ising Hamiltonian $H_{\mathrm{cl}} = -\sum_{i<j} J_{ij} Z_i Z_j - \sum_i h_i Z_i$, we have $(A|\widetilde{\mathcal{Q}}^{(\beta)}|A)_\beta^W = \mathcal{E}_{\mu_\beta}(f,f)$ and $(A|\mathcal{Q}^{(\beta)}|A)_\beta^W = \mathcal{Q}_{\mu_\beta}(f,f)$.

\begin{lemma}[Discrete Hubbard--Stratonovich Mixture]\label{lem:hs-mixture}
Every $2$-local Ising Gibbs measure~\eqref{eq:ising-gibbs} on $\{-1,+1\}^N$ is a convex mixture of product Bernoulli measures on $\{-1,+1\}^N$:
\[
  \mu_\beta(z) = \int \nu_\xi(z)\,d\mathbb{P}(\xi),
\]
where each $\nu_\xi(z) = \prod_{i=1}^N \nu_{\xi,i}(z_i)$ is a product measure with single-site magnetizations $m_i(\xi) := \E_{\nu_\xi}[z_i]$ obeying
\begin{equation}\label{eq:hs-mag}
  \sup_{\xi,\,i} |m_i(\xi)| \le \tanh\!\left(\beta H_{\max} + \frac{\Delta}{2}\ln 2\right), \qquad
  H_{\max} := \max_{1\le i\le N}\left(|h_i| + \sum_{j\ne i}|J_{ij}|\right),
\end{equation}
where $\Delta := \max_i \#\{j : J_{ij}\ne 0\}$ is the maximum vertex degree of the interaction graph.
\end{lemma}

\begin{proof}
For $\beta = 0$, we have $\phi_{ij} = 0$ and $\mu_0$ is already the uniform product measure. For $\beta > 0$ and each edge $e = \{i,j\}$ with $J_{ij}\ne 0$, set $s_{ij} := \sgn(J_{ij})\in\{-1,+1\}$ and define $\phi_{ij} := \frac{1}{2}\operatorname{arccosh}(e^{2\beta|J_{ij}|}) > 0$. For $z_i,z_j\in\{-1,+1\}$, we have $|z_i + s_{ij}z_j| = 2$ when $s_{ij}z_i z_j = +1$ and $z_i + s_{ij}z_j = 0$ when $s_{ij}z_i z_j = -1$. Since $\cosh(2\phi_{ij}) = e^{2\beta|J_{ij}|}$, we obtain the discrete Hubbard--Stratonovich identity
\[
  e^{\beta J_{ij} z_i z_j}
  = e^{-\beta|J_{ij}|}\cdot \frac{1}{2}\sum_{\xi_{ij}\in\{-1,+1\}} e^{\phi_{ij}\xi_{ij}(z_i + s_{ij}z_j)},
\]
For every $u = \beta|J_{ij}| > 0$, we have $2\phi_{ij} = \ln(e^{2u}+\sqrt{e^{4u}-1}) \le \ln(2e^{2u}) = 2u + \ln 2$, whence $\phi_{ij} \le \beta|J_{ij}| + \frac{1}{2}\ln 2$. Multiplying over all edges $\{i,j\}$ with $J_{ij}\ne 0$ and normalizing gives $\mu_\beta(z) = \sum_\xi \mathbb{P}(\xi)\nu_\xi(z)$, where conditioned on $\xi = \{\xi_{ij}\}$, the spins $\{z_i\}$ are independent with effective local fields $h_i(\xi) = \beta h_i + \sum_{j\ne i} \phi_{ij}\xi_{ij}\epsilon_{ij}$ ($\epsilon_{ij}\in\{\pm 1\}$) obeying $|h_i(\xi)| \le \beta H_{\max} + \frac{\Delta}{2}\ln 2$, and magnetizations $m_i(\xi) = \tanh(h_i(\xi))$.
\end{proof}

\begin{proposition}[All-Temperature Classical Ising Reverse-Poincar\'e Bound]\label{prop:ising}
For any $2$-local classical Ising Gibbs measure~\eqref{eq:ising-gibbs} on $\{-1,+1\}^N$ with maximum degree $\Delta := \max_i \#\{j: J_{ij}\ne 0\}$ and maximum local field $H_{\max} := \max_i(|h_i|+\sum_{j\ne i}|J_{ij}|)$, set $\Phi_{\max}(\beta) := \beta H_{\max} + \frac{\Delta}{2}\ln 2$. Then for every $\beta\ge 0$ and every $f:\{-1,+1\}^N\to\mathbb{R}$,
\begin{equation}\label{eq:ising-rp}
  \mathcal{E}_{\mu_\beta}(f,f)
  \le \cosh^2\!\bigl(\Phi_{\max}(\beta)\bigr)\,\mathcal{Q}_{\mu_\beta}(f,f).
\end{equation}
\end{proposition}

\begin{proof}
For each auxiliary configuration $\xi$ in Lemma~\ref{lem:hs-mixture}, let $\nu_\xi = \bigotimes_{i=1}^N \nu_{\xi,i}$ be the corresponding product Bernoulli measure with magnetizations $m_i(\xi) = \tanh(h_i(\xi))$ obeying $|h_i(\xi)| \le \Phi_{\max}(\beta)$, so that
\[
  \Var_{\nu_{\xi,i}}(z_i) = 1 - m_i(\xi)^2 = \sech^2(h_i(\xi)) \ge \sech^2(\Phi_{\max}(\beta)).
\]
In the $p$-biased Fourier basis $\chi_S^{(\xi)}(z) := \prod_{i\in S} \frac{z_i - m_i(\xi)}{\sqrt{1-m_i(\xi)^2}}$ ($S\subseteq\{1,\dots,N\}$), which is orthonormal in $L^2(\nu_\xi)$ and respects the degree filtration $\mathcal{V}_k^{\mathrm{cl}} = \operatorname{span}\{\chi_S^{(\xi)} : |S|\le k\}$, write $f = \sum_S \widehat{f}_\xi(S)\chi_S^{(\xi)}$. Computing the discrete derivative gives $\partial_i\chi_S^{(\xi)} = \mathbf{1}(i\in S)(1-m_i(\xi)^2)^{-1/2}\chi_{S\setminus\{i\}}^{(\xi)}$, whence
\[
  \E_{\nu_\xi}\bigl[(\partial_i f)^2\bigr]
  = \frac{1}{1-m_i(\xi)^2}\sum_{S\ni i}\widehat{f}_\xi(S)^2
  \le \cosh^2\bigl(\Phi_{\max}(\beta)\bigr)\sum_{S\ni i}\widehat{f}_\xi(S)^2.
\]
Summing over $i=1,\dots,N$ yields
\begin{align*}
  \sum_{i=1}^N \E_{\nu_\xi}\bigl[(\partial_i f)^2\bigr]
  &\le \cosh^2\bigl(\Phi_{\max}(\beta)\bigr)\sum_{S\subseteq\Lambda} |S|\,\widehat{f}_\xi(S)^2 \\
  &= \cosh^2\bigl(\Phi_{\max}(\beta)\bigr)\sum_{k=0}^{N-1}\bigl\|(\mathbf{1}-\Pi_k^{(\nu_\xi)})f\bigr\|_{L^2(\nu_\xi)}^2 \\
  &\le \cosh^2\bigl(\Phi_{\max}(\beta)\bigr)\sum_{k=0}^{N-1}\bigl\|(\mathbf{1}-\Pi_k^{(\mu_\beta)})f\bigr\|_{L^2(\nu_\xi)}^2,
\end{align*}
where the last step uses the projection theorem in $L^2(\nu_\xi)$ since $\Pi_k^{(\mu_\beta)}f\in\mathcal{V}_k^{\mathrm{cl}}$. Averaging over $\xi\sim\mathbb{P}$ and using $\Var_{\mu_\beta}(z_i\mid z_{-i})\le 1$ gives
\[
  \mathcal{E}_{\mu_\beta}(f,f)
  \le \sum_{i=1}^N \E_{\mu_\beta}\bigl[(\partial_i f)^2\bigr]
  = \int \sum_{i=1}^N \E_{\nu_\xi}\bigl[(\partial_i f)^2\bigr]\,d\mathbb{P}(\xi)
  \le \cosh^2\bigl(\Phi_{\max}(\beta)\bigr)\,\mathcal{Q}_{\mu_\beta}(f,f). \qedhere
\]
\end{proof}

\noindent\emph{Remark.}
For the zero-field one-dimensional Ising ring $H_{\mathrm{cl}} = -J\sum_{i=1}^N Z_i Z_{i+1}$ ($J>0$), evaluating on the staggered mode $f(z) = \sum_{i=1}^N (-1)^i z_i$ shows that $\mathcal{Q}_{\mu_\beta}(f,f) = \|f\|_{L^2(\mu_\beta)}^2 \sim N e^{-2\beta J}$ while the sitewise derivative sum satisfies $\sum_i \E_{\mu_\beta}[(\partial_i f)^2] = N$, so the upper-bound ratio $\sum_i\E_{\mu_\beta}[(\partial_i f)^2]/\mathcal{Q}_{\mu_\beta}(f,f) \sim e^{2\beta J}$ grows exponentially with $\beta J$ (whereas the true Dirichlet-form ratio $\mathcal{E}_{\mu_\beta}(f,f)/\mathcal{Q}_{\mu_\beta}(f,f) \sim 2 = O(1)$ because $\E_{\mu_\beta}[\Var_{\mu_\beta}(z_i\mid z_{-i})] \sim 2e^{-2\beta J}$ cancels the denominator). Thus the exponential $\beta$-dependence of $\cosh^2(\Phi_{\max}(\beta))$ is inherent to any proof proceeding through the sitewise derivative sum $\sum_i\E_{\mu_\beta}[(\partial_i f)^2]$.

\section{Structural Necessity and Obstructions}\label{sec:nogo}

The positive bounds of Theorems~\ref{thm:main-tails-reduction}--\ref{thm:main-highT} operate within the standard physical setting of local thermal operator growth: microscopic $r$-locality of the Hamiltonian and initial probe $\mathcal{O}\in\mathcal{V}_{s_0}$, localization of the initial state at the origin of the Krylov chain, and thermal clustering of the equilibrium state (which holds automatically for $\beta\le\min(\beta_{\mathrm{wc}},\beta_c)$ by Lemma~\ref{lem:2design-sep} and Proposition~\ref{prop:entropic}(2), and enters Theorem~\ref{thm:main-tails-reduction}(b) via the static criterion~\eqref{eq:main-H2}). In this section we show that stepping outside any one of these three structural features leads to quantitative violations of the corresponding inequality, thereby clarifying the natural domain of validity of~\eqref{eq:intro-uogh}.

\begin{proposition}[Structural Necessity Results]\label{prop:nogo}
\begin{enumerate}
  \setlength{\itemsep}{3pt}\setlength{\parskip}{0pt}
  \item[\emph{(1)}] \emph{Necessity of microscopic locality.} Suppose a thermal system $(H,\rho_\beta)$ and local Hermitian operators $W_0, V$ exhibit exponential scrambling in the regulated four-point correlator,
  \begin{align*}
    F_{W_0,V}(t_1,t_2)
    &:= \Tr\!\bigl(\rho_\beta^{1/4}W_0(t_1)\rho_\beta^{1/4}V\rho_\beta^{1/4}W_0(t_2)\rho_\beta^{1/4}V\bigr) \\
    &\phantom{:}= F_d - \varepsilon\,e^{\lambda_L(t_1+t_2)/2} + o\!\left(\varepsilon\,e^{\lambda_L(t_1+t_2)/2}\right)
  \end{align*}
  with $\lambda_L > 0$ on a scrambling window $t_d \ll t_1, t_2 \ll t_*$ (where $t_* \sim \lambda_L^{-1}\ln(1/\varepsilon)$). For any $\sigma > 0$, define the Gaussian time-smeared operator
  \[
    W_\sigma := \int_{\mathbb{R}} \frac{e^{-s^2/(2\sigma^2)}}{\sqrt{2\pi}\,\sigma}\,W_0(s)\,ds.
  \]
  Then $d\mu_{\beta,W_\sigma}^W(\omega) = e^{-\sigma^2\omega^2}d\mu_{\beta,W_0}^W(\omega)$, so $C_{\beta,W_\sigma}^W(z)$ is an entire function of $z\in\mathbb{C}$, the Wightman Lanczos coefficients of $W_\sigma$ obey $(\prod_{k=1}^n b_{k,T}^{(W)}(W_\sigma))^{1/n} = O(\sqrt{n})$ and $\liminf_{n\to\infty} b_{n,T}^{(W)}(W_\sigma)/n = 0$, whereas the regulated four-point correlator $F_{W_\sigma,V}(t)$ retains the same Lyapunov exponent $\lambda_L > 0$ throughout $t_d + \sigma\sqrt{\ln(1/\varepsilon)} \ll t \ll t_* - \sigma\sqrt{\ln(1/\varepsilon)}$.
  \item[\emph{(2)}] \emph{Failure of Hypothesis~\emph{(H2)} without thermal clustering.} On $N\ge 2$ qubits, consider the non-clustering classical mixture
  \begin{equation}\label{eq:bimodal-state}
    \rho := \frac{1}{2}\,|0\rangle\langle 0|_1 \otimes \frac{I_{N-1}}{2^{N-1}}
      + \frac{1}{2}\,|1\rangle\langle 1|_1 \otimes |0\cdots 0\rangle\langle 0\cdots 0|_{2\dots N},
  \end{equation}
  and the normalized $1$-local Hermitian operator $A := \sqrt{2}\,\frac{I - Z_1}{2} = \sqrt{2}\,|1\rangle\langle 1|_1 \in \mathcal{V}_1$ (which satisfies $\|A\|_{\rho,W} = 1$). Then
  \begin{equation}\label{eq:bimodal-size}
    \mathcal{S}_\rho(A) = \frac{N-1}{4} + 2^{-(N+1)/2} \qquad \forall N\ge 2.
  \end{equation}
  In particular, $\Lambda_\rho^{\MSS}(1) = \Theta(N)$, violating~\eqref{eq:main-H2} for every $\delta\in[0,1)$.
  \item[\emph{(3)}] \emph{Obstruction to arbitrary-spread $q$-complexity bounds.} For every integer $D\ge 4$ and barrier height $E_B > 2$, there exist a $1$-banded symmetric Hermitian operator $\mathcal{L}$ on $\mathbb{C}^D$ in the eigenbasis $\{|0\rangle,\dots,|D-1\rangle\}$ of $\mathcal{Q} = \operatorname{diag}(0,1,\dots,D-1)$ and a normalized initial state $|\psi_0\rangle\in\mathbb{C}^D$ with variance $\sigma_0^2 := \Var_{\psi_0}(\mathcal{Q}) = O(E_B^{-2}) \to 0$ as $E_B\to\infty$, such that at time $t_*(E_B) = \pi/|E_- - E_+|$ the Krylov complexity satisfies $K(t_*) = 1$ while $\Delta Q(t_*) := \langle\psi(t_*)|\mathcal{Q}|\psi(t_*)\rangle - \langle\psi_0|\mathcal{Q}|\psi_0\rangle = (D-1) - O(E_B^{-2})$. Consequently,
  \begin{equation}\label{eq:doublewell-ratio}
    \lim_{E_B\to\infty} \frac{\Delta Q(t_*)}{K(t_*) + \sigma_0\sqrt{K(t_*)}} = D - 1,
  \end{equation}
  precluding any universal bound of the form $\Delta Q(t) \le C_1 K(t) + C_2 \sigma_0\sqrt{K(t)}$ independent of $D$.
\end{enumerate}
\end{proposition}

\begin{proof}
\emph{Part~(1).} By the convolution theorem for the Fourier transform, the Wightman autocorrelation function of $W_\sigma$ satisfies $C_{\beta,W_\sigma}^W(t) = \int_{\mathbb{R}} e^{i\omega t} e^{-\sigma^2\omega^2}\,d\mu_{\beta,W_0}^W(\omega)$, so $d\mu_{\beta,W_\sigma}^W(\omega) = e^{-\sigma^2\omega^2}\,d\mu_{\beta,W_0}^W(\omega)$ and $C_{\beta,W_\sigma}^W(z)$ extends to an entire function on $\mathbb{C}$. Estimating the $2n$-th normalized moment $\mu_{2n}(W_\sigma)$ gives
\[
  \mu_{2n}(W_\sigma)
  = \frac{1}{\|W_\sigma\|_{\beta,W}^2}\int_{\mathbb{R}} \omega^{2n} e^{-\sigma^2\omega^2}\,d\mu_{\beta,W_0}^W(\omega)
  \le C'\left(\frac{n}{e\sigma^2}\right)^n
\]
for a constant $C' := \|W_0\|_{\beta,W}^2/\|W_\sigma\|_{\beta,W}^2$ independent of $n$ (since $\sup_{\omega\in\mathbb{R}} \omega^{2n} e^{-\sigma^2\omega^2} = (n/(e\sigma^2))^n$ and $\int d\mu_{\beta,W_0}^W(\omega) = \|W_0\|_{\beta,W}^2 < \infty$). In the Dyck-path expansion $\mu_{2n}(W_\sigma) = \sum_{\gamma\in\mathcal{D}_{2n}}\prod_{k=1}^{2n} b_{\max(\gamma_{k-1},\gamma_k)}$, keeping only the single path that climbs from $0$ to $n$ and descends back to $0$ yields the moment--Lanczos inequality $\prod_{k=1}^n b_k^2 \le \mu_{2n}(W_\sigma)$. Taking the $(2n)$-th root gives $(\prod_{k=1}^n b_k)^{1/n} = O(\sqrt{n})$, which forces $\liminf_{n\to\infty} b_n/n = 0$.

On the other hand, expanding $F_{W_\sigma,V}(t) = \iint \frac{e^{-(s_1^2+s_2^2)/(2\sigma^2)}}{2\pi\sigma^2} F_{W_0,V}(t+s_1, t+s_2)\,ds_1\,ds_2$ and splitting the integral into $|s_1|,|s_2| \le \Delta t$ and its complement (choosing $\Delta t$ with $t - t_d, t_* - t \gg \Delta t \gg \sigma\sqrt{\ln(1/\varepsilon)}$), the tail integral is bounded by $O(\|W_0\|_\infty^2\|V\|_\infty^2 e^{-\Delta t^2/(2\sigma^2)}) = o(\varepsilon)$, while the core integral evaluates to
\[
  F_d - \varepsilon\,e^{\lambda_L t}\left(\int_{-\Delta t}^{\Delta t} \frac{e^{-s^2/(2\sigma^2) + \lambda_L s/2}}{\sqrt{2\pi}\,\sigma}\,ds\right)^2 + o(\varepsilon e^{\lambda_L t})
  = F_d - \varepsilon\,e^{\sigma^2\lambda_L^2/4}\,e^{\lambda_L t} + o(\varepsilon e^{\lambda_L t}).
\]
Thus $F_{W_\sigma,V}(t)$ retains the Lyapunov exponent $\lambda_L > 0$ despite $\liminf_{n\to\infty} b_n/n = 0$.

\emph{Part~(2).} Because the two summands in~\eqref{eq:bimodal-state} have orthogonal supports on qubit $1$, the square root $\rho^{1/2}$ is
\[
  \rho^{1/2} = \frac{1}{\sqrt{2}}\,|0\rangle\langle 0|_1 \otimes \frac{I_{N-1}}{2^{(N-1)/2}}
             + \frac{1}{\sqrt{2}}\,|1\rangle\langle 1|_1 \otimes |0\cdots 0\rangle\langle 0\cdots 0|_{2\dots N}.
\]
Let us compute $\langle\rho^{1/2},\mathcal{S}\rho^{1/2}\rangle_\HS$ directly in the Pauli basis. On sites $2,\dots,N$, we have $|0\cdots 0\rangle\langle 0\cdots 0| = 2^{-(N-1)}\sum_{S\subseteq\{2,\dots,N\}} Z_S$, where $Z_S := \prod_{x\in S} Z_x$. Writing $|0\rangle\langle 0|_1 = \frac{1}{2}(I_1+Z_1)$ and $|1\rangle\langle 1|_1 = \frac{1}{2}(I_1-Z_1)$ gives
\begin{align*}
  \rho^{1/2}
  &= \frac{1}{2\sqrt{2}}\left(2^{-(N-1)/2} + 2^{-(N-1)}\right) I
   + \frac{1}{2\sqrt{2}}\left(2^{-(N-1)/2} - 2^{-(N-1)}\right) Z_1 \\
  &\quad + \frac{1}{2\sqrt{2}}\,2^{-(N-1)}\sum_{\substack{S\subseteq\{2,\dots,N\} \\ S\ne\emptyset}} (I_1 - Z_1)\otimes Z_S.
\end{align*}
Since $\|P\|_\HS^2 = 2^N$ for every Pauli string $P$, the contribution of $Z_1$ (weight $1$) to $\langle\rho^{1/2},\mathcal{S}\rho^{1/2}\rangle_\HS$ is
\[
  2^N\cdot \frac{1}{8}\left(2^{-(N-1)/2} - 2^{-(N-1)}\right)^2
  = \frac{1}{4}\left(1 - 2^{-(N-1)/2}\right)^2
  = \frac{1}{4} - 2^{-(N+1)/2} + 2^{-(N+1)}.
\]
For each nonempty $S\subseteq\{2,\dots,N\}$ of cardinality $k = |S|\ge 1$, the string $I_1\otimes Z_S$ has weight $k$ and the string $-Z_1\otimes Z_S$ has weight $k+1$, each with squared coefficient $\frac{1}{8}2^{-2(N-1)} = 2^{-(2N+1)}$. Multiplying by $2^N$ and summing over all nonempty $S$ gives
\[
  2^{-(N+1)}\sum_{k=1}^{N-1}\binom{N-1}{k}(2k+1)
  = 2^{-(N+1)}\Bigl((N-1)2^{N-1} + 2^{N-1} - 1\Bigr)
  = \frac{N}{4} - 2^{-(N+1)}.
\]
Adding the $Z_1$ contribution cancels $2^{-(N+1)}$ and yields
\[
  \langle\rho^{1/2},\mathcal{S}\rho^{1/2}\rangle_\HS = \frac{N+1}{4} - 2^{-(N+1)/2}.
\]
On the other hand, for $A = \sqrt{2}\,|1\rangle\langle 1|_1 \in \mathcal{V}_1$, we have $\mathcal{J}_\rho(A) = \rho^{1/4}A\rho^{1/4} = |1\rangle\langle 1|_1 \otimes |0\cdots 0\rangle\langle 0\cdots 0|_{2\dots N}$, which is a pure computational-basis projector of norm $\|\mathcal{J}_\rho(A)\|_\HS = 1$ and Hilbert--Schmidt operator size $\langle\mathcal{J}_\rho(A),\mathcal{S}\mathcal{J}_\rho(A)\rangle_\HS = \frac{N}{2}$ (since each of the $N$ sites carries $|0\rangle\langle 0|$ or $|1\rangle\langle 1|$, contributing $\frac{1}{2}$ per site). Subtracting gives
\[
  \mathcal{S}_\rho(A) = \frac{N}{2} - \left(\frac{N+1}{4} - 2^{-(N+1)/2}\right) = \frac{N-1}{4} + 2^{-(N+1)/2},
\]
proving~\eqref{eq:bimodal-size}. Notice that although $\rho$ in~\eqref{eq:bimodal-state} is singular, perturbing it to the full-rank state $\rho_\epsilon := (1-\epsilon)\rho + \epsilon I/2^N$ ($\epsilon > 0$) preserves $\Lambda_{\rho_\epsilon}^{\MSS}(1) = \frac{N-1}{4} + O(\epsilon N)$.

\emph{Part~(3).} Fix an integer $D\ge 4$ and a barrier height $E_B > 2$. Define the $D\times D$ real symmetric tridiagonal ($1$-banded) matrix $\mathcal{L}$ in the basis $\{|0\rangle,\dots,|D-1\rangle\}$ by diagonal entries $a_0 = a_{D-1} = 0$, $a_1 = \dots = a_{D-2} = E_B$, and unit off-diagonal entries $b_0 = b_1 = \dots = b_{D-2} = 1$. Because every off-diagonal entry is nonzero, $\mathcal{L}$ is unreduced, so all $D$ eigenvalues are simple; moreover, $\mathcal{L}$ commutes with the spatial inversion unitary $P|j\rangle = |D-1-j\rangle$, so every eigenvector has definite parity $P|\phi\rangle = \pm|\phi\rangle$.

By Lemma~\ref{lem:app-doublewell} (proved in Appendix~\ref{app:doublewell}), for all $E_B > 2$ the two lowest eigenvalues $E_+, E_-$ of $\mathcal{L}$ (with $E_+$ even parity and $E_-$ odd parity) are distinct and lie in $[-(E_B-2)^{-1}, 0]$, and their normalized eigenstates $|\phi_\pm\rangle$ satisfy
\[
  \langle 0|\phi_\pm\rangle = \frac{1}{\sqrt{2}} - O(E_B^{-2}), \qquad
  \langle D-1|\phi_\pm\rangle = \pm\left(\frac{1}{\sqrt{2}} - O(E_B^{-2})\right),
\]
and $|\langle j|\phi_\pm\rangle| = O(E_B^{-\min(j,D-1-j)})$ for $1\le j\le D-2$. Define the normalized initial state
\[
  |\psi_0\rangle := \frac{1}{\sqrt{2}}\bigl(|\phi_+\rangle + |\phi_-\rangle\bigr).
\]
Then $\langle 0|\psi_0\rangle = 1 - O(E_B^{-2})$, $|\langle D-1|\psi_0\rangle| = \frac{1}{\sqrt{2}}|\langle 0|\phi_+\rangle - \langle 0|\phi_-\rangle| = O(E_B^{-2})$ (from~\eqref{eq:dw-boundary}), and $|\langle j|\psi_0\rangle| = O(E_B^{-\min(j,D-1-j)})$ for $1\le j\le D-2$. In particular,
\begin{align*}
  \langle\psi_0|\mathcal{Q}|\psi_0\rangle &= (D-1)|\langle D-1|\psi_0\rangle|^2 + \sum_{j=1}^{D-2} j|\langle j|\psi_0\rangle|^2 = O(E_B^{-2}), \\
  \sigma_0^2 &= \Var_{\psi_0}(\mathcal{Q}) \le \langle\psi_0|\mathcal{Q}^2|\psi_0\rangle = O(E_B^{-2}),
\end{align*}
so $\sigma_0 = O(E_B^{-1}) \to 0$ as $E_B\to\infty$.

Because $|\psi_0\rangle$ lies in the $2$-dimensional $\mathcal{L}$-invariant subspace $\operatorname{span}\{|\phi_+\rangle,|\phi_-\rangle\}$, the Krylov subspace generated by $(\mathcal{L},|\psi_0\rangle)$ has dimension $2$, with orthonormal Lanczos basis
\[
  |K_0\rangle = |\psi_0\rangle, \qquad
  |K_1\rangle = \frac{1}{\sqrt{2}}\bigl(|\phi_+\rangle - |\phi_-\rangle\bigr) = P|\psi_0\rangle.
\]
Under $|\psi(t)\rangle = e^{-i\mathcal{L}t}|\psi_0\rangle$, we obtain
\[
  |\psi(t)\rangle
  = e^{-i\frac{E_++E_-}{2}t}\left(\cos\!\left(\frac{E_--E_+}{2}\,t\right)|K_0\rangle
    + i\sin\!\left(\frac{E_--E_+}{2}\,t\right)|K_1\rangle\right).
\]
At the tunneling time $t_*(E_B) := \frac{\pi}{|E_--E_+|}$, we have $|\psi(t_*)\rangle = i e^{-i\frac{E_++E_-}{2}t_*}|K_1\rangle = i e^{-i\frac{E_++E_-}{2}t_*}P|\psi_0\rangle$. Consequently, the Krylov complexity is $K(t_*) = 1$, whereas
\[
  \langle\psi(t_*)|\mathcal{Q}|\psi(t_*)\rangle
  = \langle\psi_0|P\mathcal{Q}P|\psi_0\rangle
  = (D-1) - \langle\psi_0|\mathcal{Q}|\psi_0\rangle
  = (D-1) - O(E_B^{-2}),
\]
giving $\Delta Q(t_*) = (D-1) - O(E_B^{-2})$. Dividing by $K(t_*) + \sigma_0\sqrt{K(t_*)} = 1 + O(E_B^{-1})$ and letting $E_B\to\infty$ yields~\eqref{eq:doublewell-ratio}.
\end{proof}

\noindent\emph{Remark.}
Proposition~\ref{prop:nogo}(1) explains why effective field theory models with shift-symmetric non-microscopic correlators~\cite{Demulder2026} can exhibit $2\alpha_T^{(W)} = \lambda_L/2 < \lambda_L$: the microscopic $r$-locality condition $\mathcal{L}(\mathcal{V}_s)\subseteq\mathcal{V}_{s+r-1}$ is absent in the continuum hydrodynamic approximation. Proposition~\ref{prop:nogo}(3) similarly shows why the $q$-complexity bound of Proposition~\ref{prop:q-complexity}(2) cannot be extended to arbitrary initial states with small but non-vanishing tails in the filtration: resonant tunneling through an arbitrarily high barrier produces an $O(D)$ jump in $\mathcal{Q}$ while keeping the Krylov dimension equal to $2$.

\section{Finite-Size Numerical Benchmarks and Open Frontier}\label{sec:numerics}

To complement our analytical bounds and test their quantitative tightness on concrete spin and fermionic models, we performed numerical diagonalization checks across five benchmark systems:
\begin{enumerate}
  \setlength{\itemsep}{2pt}\setlength{\parskip}{0pt}
  \item the $\eta=1$ Meixner chain $b_n = n$ and the staggered linear sequence $b_n = n + \frac{1}{2}(-1)^n$ (truncated at $N_{\max}=40$, with $\theta_0 = 4$ and $t\in[0,1.5]$), testing Theorem~\ref{thm:main-tails-reduction}(a);
  \item a full-rank product thermal state on $L=4$ qubits with uniform polarization $m_x = 0.5$ at weight $s=1$, testing Proposition~\ref{prop:product};
  \item the one-dimensional mixed-field Ising chain $H = \sum_{i=1}^{L-1} Z_i Z_{i+1} + \sum_{i=1}^L (-1.05 X_i + 0.5 Z_i)$ at $L=4$ and $\beta=1.5, 2.0$ (with $W=Z_1$, $V=Z_4$, $\tau=0.1$, $t\in[0,2]$), testing Theorem~\ref{thm:main-spectral}(c) and Hypothesis~(H2) of Theorem~\ref{thm:main-tails-reduction}(b);
  \item the two-dimensional $2\times 2$ open-plaquette mixed-field Ising model $H = \sum_{\langle i,j\rangle} Z_i Z_j + \sum_{i=1}^4 (-1.05 X_i + 0.5 Z_i)$ ($4$ edges), testing Lemma~\ref{lem:2design-sep} over all $6^4 = 1296$ frame configurations $\mathbf{s}\in\{1,\dots,6\}^4$ at $\beta=0.015$ and the generalized eigenvalue ratios of $\sum_x \mathcal{S}_x^{\dagger W}\mathcal{S}_x$ and $\mathrm{I}_{\mathrm{op}}$ relative to $\mathcal{Q}^{(\beta)}$ on $\mathbf{1}^\perp$ at $\beta=0.2$ (Theorem~\ref{thm:main-highT}(a)--(b));
  \item a single disorder realization of the $N=8$ SYK-$4$ model at $\beta=6$ (with coupling variance $\E[J_{ijkl}^2] = 6\mathcal{J}^2/N^3$ and $\mathcal{J}=1$), the bimodal mixture~\eqref{eq:bimodal-state} at $L=5$, and the symmetric double-well Jacobi matrix of Proposition~\ref{prop:nogo}(3) at $D=12, E_B=10$.
\end{enumerate}

Table~\ref{tab:num} collects the observed numerical ratios alongside the corresponding analytical bounds. Notice that on the $2\times 2$ mixed-field Ising plaquette at $\beta = 0.2$, the largest generalized eigenvalues of $\sum_x \mathcal{S}_x^{\dagger W}\mathcal{S}_x$ and $\mathrm{I}_{\mathrm{op}}$ relative to $\mathcal{Q}^{(\beta)}$ on $\mathbf{1}^\perp$ are $1.5577$ and $1.6834$, well below the uniform analytical constants $9$ and $216$ of Theorem~\ref{thm:main-highT}. At low temperature ($\beta = 2.0$) on the same plaquette, however, the generalized eigenvalue ratio $\max(\mathrm{I}_{\mathrm{op}}/\mathcal{Q}^{(\beta)})$ increases to $1.32\times 10^6$, consistent with the exponential growth in $\beta$ observed in the classical Ising remark of Section~\ref{sec:ising}.

\begin{table}[!htbp]
\centering
\footnotesize
\begin{tabularx}{\linewidth}{@{}>{\raggedright\arraybackslash}p{4.4cm} >{\raggedright\arraybackslash}X c c@{}}
\toprule
\textbf{Result / Model} & \textbf{Quantity Tested} & \textbf{Observed} & \textbf{Bound} \\
\midrule
Thm.~\ref{thm:main-tails-reduction}(a), Meixner $\eta=1$ & $G(\theta(t),t)/e^{\theta_0/2}$ ($\theta_0=4, t\le 1.5$) & $0.4849$ & $\le 1.0000$ \\
Thm.~\ref{thm:main-tails-reduction}(a), staggered $b_n$ & $G(\theta(t),t)/e^{(\frac{1}{2}+\frac{\gamma}{\alpha})\theta_0}$ ($\theta_0=4, t\le 1.5$) & $0.0312$ & $\le 1.0000$ \\
Prop.~\ref{prop:product}, product ($L=4$) & $\Lambda_\rho^{\MSS}(1)/(1 + m_{\max}\sqrt{L})$ ($m_x=0.5$) & $0.5472$ & $\le 1.0000$ \\
Thm.~\ref{thm:main-spectral}(c), 1D Ising ($\beta=1.5$) & $|F(t+i\tau)|/\min(\|V\|_\infty^2 C_W,\|W\|_\infty^2 C_V,T_\tau)$ & $0.4942$ & $\le 1.0000$ \\
Lem.~\ref{lem:2design-sep}, 2D Ising ($\beta=0.015$) & $\min_{\mathbf{s}}\widetilde{\omega}_{\mathbf{s}}(e^{-\beta H/2})$ (1296 cfgs) & $+0.0355$ & $> 0$ \\
Thm.~\ref{thm:main-highT}(a), 2D Ising ($\beta=0.2$) & $\max(\sum_x \mathcal{S}_x^{\dagger W}\mathcal{S}_x / \mathcal{Q}^{(\beta)})$ on $\mathbf{1}^\perp$ & $1.5577$ & $\le 9.0000$ \\
Thm.~\ref{thm:main-highT}(b), 2D Ising ($\beta=0.2$) & $\max(\mathrm{I}_{\mathrm{op}}/\mathcal{Q}^{(\beta)})$ on $\mathbf{1}^\perp$ & $1.6834$ & $\le 216.0$ \\
Prop.~\ref{prop:nogo}(2), bimodal ($L=5$) & $\mathcal{S}_\rho(\sqrt{2}|1\rangle\langle 1|_1)$ & $1.1250$ & $= 1.1250$ \\
Prop.~\ref{prop:nogo}(3), double well ($D=12$) & $\Delta Q(t_*)/(K(t_*)+\sigma_0\sqrt{K(t_*)})$ & $9.9725$ & $\to 11$ \\
Thm.~\ref{thm:main-tails-reduction}(b) / (H2), SYK-$4$ ($N=8,\beta=6$) & $\Lambda_\beta^{\MSS}(s)$ ($s=1,2,3,4$) & $0.90, 2.55, 3.73, 5.73$ & $O(s)$ \\
Thm.~\ref{thm:main-tails-reduction}(b) / (H2), 1D Ising ($L=4,\beta=2$) & $\Lambda_\beta^{\MSS}(s)$ ($s=1,2,3,4$) & $0.90, 1.72, 2.17, 2.28$ & $O(s)$ \\
\bottomrule
\end{tabularx}
\caption{Finite-size numerical diagonalization benchmarks of the analytical bounds and constructions.}\label{tab:num}
\end{table}

We conclude by noting the remaining open frontier. Theorems~\ref{thm:main-tails-reduction}--\ref{thm:main-highT} establish $\lambda_{L,T}\le 2\alpha_T^{(W)}\le 2\pi/(\beta+2\kappa) < 2\pi/\beta$ without smoothness hypotheses for arbitrary non-commuting $r$-local qubit Hamiltonians in any spatial dimension $d\ge 1$ at high temperature $\beta\in[0,\beta_{\mathrm{wc}}]$, for both the Wightman-commutator and MSS-regulated OTOCs. Combined with our all-temperature estimates for product states (Proposition~\ref{prop:product}), one-dimensional local spin chains (Proposition~\ref{prop:1D}), and classical $2$-local Ising models (Proposition~\ref{prop:ising}), the open regime is localized to two settings: (i)~low-temperature entangled quantum phases in $d\ge 2$ beyond the cluster-expansion threshold $\beta > \beta_{\mathrm{wc}}$, where one expects~\eqref{eq:intro-uogh} to hold in any pure thermal phase exhibiting exponential decay of correlations~\cite{Kliesch2014,Frohlich2015,Kato2019,Kuwahara2020b,Alhambra2023,Rouze2026}; and (ii)~dense all-to-all $q$-local models such as finite-$q$ SYK~\cite{MS2016,Roberts2018,Qi2019,Yin2020,Garcia2025}, where $J_{\mathrm{site}}$ is finite only after $N$-dependent normalization. In both settings, Theorem~\ref{thm:main-tails-reduction}(b) reduces the left inequality of~\eqref{eq:intro-uogh} to establishing the static equilibrium criterion~\eqref{eq:main-H2}.

\appendix

\section{Rooted Plane-Tree and Cayley Spanning-Tree Polymer Estimates}\label{app:trees}

In this appendix we provide self-contained proofs of the two combinatorial tree-summation bounds used in Section~\ref{sec:highT}: the rooted plane-tree estimate for connected polymers (used in Lemma~\ref{lem:2design-sep}) and the Cayley spanning-tree root-reversal estimate for ordered causal sequences (used in the proof of Theorem~\ref{thm:main-highT}(b)).

\begin{lemma}[Rooted Plane-Tree Polymer Sum]\label{lem:app-planetree}
Let $\mathcal{E}$ be a collection of nonempty subsets $X\subseteq\Lambda$ of cardinality $|X|\le r$, and let $\{v_X\}_{X\in\mathcal{E}}$ be nonnegative weights satisfying $\sup_{x\in\Lambda}\sum_{X\ni x} v_X \le \eta_0$. For any fixed site $x_0\in\Lambda$ and any integer $m\ge 1$, the sum over all connected polymers $\gamma\subseteq\mathcal{E}$ of cardinality $|\gamma|=m$ with $x_0\in U_\gamma := \bigcup_{X\in\gamma} X$ satisfies
\begin{equation}\label{eq:app-planetree}
  \sum_{\substack{\gamma:\, |\gamma|=m \\ x_0\in U_\gamma}} \prod_{X\in\gamma} v_X
  \le \eta_0\,(4r\eta_0)^{m-1}.
\end{equation}
\end{lemma}

\begin{proof}
Fix a total order on $\mathcal{E}$ and on $\Lambda$. Given a connected polymer $\gamma = \{X_1,\dots,X_m\}$ with $x_0\in U_\gamma$, choose the first hyperedge $X_{\mathrm{root}}\in\gamma$ containing $x_0$, and construct a canonical rooted spanning tree $T(\gamma)$ of the intersection graph of $\gamma$ rooted at $X_{\mathrm{root}}$ (for example, by breadth-first search using the fixed order on $\mathcal{E}$). Ordering the children of each vertex according to the fixed order on $\mathcal{E}$ equips $T(\gamma)$ with the structure of a rooted plane tree on $m$ vertices.

A rooted plane tree on $m$ vertices is uniquely encoded by its depth-first contour walk of length $2(m-1)$ (a sequence of $m-1$ ``up'' steps $U$ and $m-1$ ``down'' steps $D$, i.e., a Dyck word of length $2(m-1)$). The number of rooted plane trees on $m$ vertices is the Catalan number $C_{m-1} = \frac{1}{m}\binom{2m-2}{m-1} \le 4^{m-1}$.

Now fix any rooted plane tree shape $\tau$ on $m$ vertices, and sum $\prod_{v\in\tau} v_{X_v}$ over all assignments $v\mapsto X_v\in\mathcal{E}$ such that $x_0\in X_{\mathrm{root}}$ and $X_v\cap X_{p(v)}\ne\emptyset$ whenever $p(v)$ is the parent of $v$ in $\tau$ (dropping the distinctness constraint on $\{X_v\}$ only increases the nonnegative sum). We evaluate this sum inductively from the leaves of $\tau$ to the root: for any leaf $v$ with parent $p(v)$, summing over all $X_v\in\mathcal{E}$ intersecting $X_{p(v)}$ gives
\[
  \sum_{X_v:\, X_v\cap X_{p(v)}\ne\emptyset} v_{X_v}
  \le \sum_{y\in X_{p(v)}}\sum_{X_v\ni y} v_{X_v}
  \le |X_{p(v)}|\,\eta_0 \le r\eta_0.
\]
Peeling off all $m-1$ non-root vertices inductively yields a factor $(r\eta_0)^{m-1}$, and the final sum over $X_{\mathrm{root}}\ni x_0$ contributes at most $\eta_0$. Summing over all $C_{m-1} \le 4^{m-1}$ plane-tree shapes $\tau$ yields~\eqref{eq:app-planetree}.
\end{proof}

\begin{lemma}[Cayley Spanning-Tree Root-Reversal Estimate]\label{lem:app-cayley}
Let $H = \sum_{X\in\mathcal{E}} h_X$ be an $r$-local Hamiltonian with per-site interaction strength $J_{\mathrm{site}} < \infty$. For any fixed site $x\in\Lambda$ and any integer $k\ge 1$, let $\mathcal{C}_k(i)$ denote the set of ordered causal sequences $\Gamma = (X_1,\dots,X_k)\in\mathcal{E}^k$ rooted at $i\in\Lambda$ ($X_1\ni i$, $X_j\cap(\{i\}\cup X_1\cup\dots\cup X_{j-1})\ne\emptyset$ for $2\le j\le k$), with support $U_\Gamma := \{i\}\cup X_1\cup\dots\cup X_k$. Then
\begin{equation}\label{eq:app-cayley}
  \sum_{i\in\Lambda}\sum_{\substack{\Gamma\in\mathcal{C}_k(i) \\ x\in U_\Gamma}} \frac{1}{k!}\prod_{j=1}^k \|h_{X_j}\|_\infty
  \le (er J_{\mathrm{site}})^k.
\end{equation}
\end{lemma}

\begin{proof}
Let $\Gamma = (X_1,\dots,X_k)\in\mathcal{C}_k(i)$ satisfy $x\in U_\Gamma = \{i\}\cup X_1\cup\dots\cup X_k$. Since $k\ge 1$ and $i\in X_1$, we have $U_\Gamma = X_1\cup\dots\cup X_k$, so $x\in X_a$ for at least one index $a\in\{1,\dots,k\}$. Moreover, the intersection graph on the $k$ labeled vertices $\{1,\dots,k\}$ (where $u$ and $v$ are adjacent if $X_u\cap X_v\ne\emptyset$) is connected, so it contains at least one spanning tree $\tau$ on $\{1,\dots,k\}$.

Root $\tau$ at the vertex $a\in\{1,\dots,k\}$ satisfying $x\in X_a$. By Cayley's theorem, there are $k^{k-2}$ labeled spanning trees on $\{1,\dots,k\}$, and $k$ choices for the root vertex $a\in\{1,\dots,k\}$, giving $k^{k-1}$ rooted labeled trees $(\tau,a)$ in total. For each fixed rooted labeled tree $(\tau,a)$ and each assignment $(X_1,\dots,X_k)$ satisfying $x\in X_a$ and $X_v\cap X_{p_\tau(v)}\ne\emptyset$ for all $v\ne a$, the original root site $i\in X_1$ has at most $|X_1|\le r$ possible values. Summing $\prod_{j=1}^k \|h_{X_j}\|_\infty$ from the leaves of $\tau$ up to the root $a$ just as in Lemma~\ref{lem:app-planetree} yields
\[
  \sum_{i\in X_1}\sum_{\substack{(X_1,\dots,X_k):\, x\in X_a \\ X_v\cap X_{p_\tau(v)}\ne\emptyset}} \prod_{j=1}^k \|h_{X_j}\|_\infty
  \le r\cdot (r J_{\mathrm{site}})^{k-1} J_{\mathrm{site}} = (r J_{\mathrm{site}})^k.
\]
Summing over all $k^{k-1}$ rooted labeled trees $(\tau,a)$ and dividing by $k!$ gives
\[
  \sum_{i\in\Lambda}\sum_{\substack{\Gamma\in\mathcal{C}_k(i) \\ x\in U_\Gamma}} \frac{1}{k!}\prod_{j=1}^k \|h_{X_j}\|_\infty
  \le \frac{k^{k-1}}{k!}\,(r J_{\mathrm{site}})^k
  \le (er J_{\mathrm{site}})^k,
\]
where the last step uses the elementary bound $k! \ge (k/e)^k$, i.e., $k^{k-1}/k! \le e^k/k \le e^k$.
\end{proof}

\begin{lemma}[Multi-Type Unordered Spanning-Tree Polymer Sum]\label{lem:app-multitree}
Let $\mathcal{F} := \bigsqcup_{j=1}^p \mathcal{E}_j$ be a finite collection of hyperedges on $\Lambda$ equipped with a fixed total order $<$ and nonnegative weights $\mu(E) \ge 0$ ($E\in\mathcal{F}$). Suppose there exists a constant $c_0 > 0$ such that
\begin{equation}\label{eq:app-pin}
  \Phi := \sup_{y\in\Lambda}\sum_{\substack{E\in\mathcal{F} \\ y\in E}} \mu(E)\,e^{c_0|E|} < c_0.
\end{equation}
Then for every site $y\in\Lambda$, the sum over all connected polymers $\gamma\subseteq\mathcal{F}$ with $y\in U_\gamma := \bigcup_{E\in\gamma} E$ satisfies
\begin{equation}\label{eq:app-multitree}
  \sum_{\substack{\gamma\subseteq\mathcal{F}\text{ connected} \\ y\in U_\gamma}} \prod_{E\in\gamma} \mu(E)
  \le \Phi.
\end{equation}
\end{lemma}

\begin{proof}
For each integer $h\ge 0$ and each hyperedge $E\in\mathcal{F}$, define $S_h(E)$ as the sum of $\prod_{u\in T}\mu(E_u)$ over all rooted trees $T$ of height at most $h$ with root $E_0 = E$, whose vertices carry distinct hyperedges in $\mathcal{F}$ such that at every node $u$, the children of $u$ are strictly increasing in the order $<$ and each child $E'$ satisfies $E'\cap E_u \ne \emptyset$. At $h = 0$, $S_0(E) = \mu(E) \le \mu(E)e^{c_0|E|}$. Inductively, for $h\ge 1$, grouping trees by the number $c\ge 0$ of children of the root $E$ (and dividing by $c!$ to account for unordered, strictly increasing child sets) gives
\begin{align*}
  S_h(E)
  &\le \mu(E)\sum_{c=0}^\infty \frac{1}{c!}\left(\sum_{\substack{E'\in\mathcal{F} \\ E'\cap E\ne\emptyset}} S_{h-1}(E')\right)^c
   \le \mu(E)\exp\!\left(\sum_{z\in E}\sum_{\substack{E'\in\mathcal{F} \\ z\in E'}} \mu(E')e^{c_0|E'|}\right) \\
  &\le \mu(E)e^{|E|\Phi} \le \mu(E)e^{c_0|E|},
\end{align*}
where the last inequality uses $\Phi < c_0$. Every connected polymer $\gamma\subseteq\mathcal{F}$ with $y\in U_\gamma$ admits a canonical spanning tree rooted at the $<$-minimal hyperedge $E_0\in\gamma$ containing $y$ (ordered by breadth-first search with strictly increasing children at each node). Passing to the limit $h\to\infty$ by monotone convergence therefore yields
\[
  \sum_{\substack{\gamma\subseteq\mathcal{F}\text{ connected} \\ y\in U_\gamma}} \prod_{E\in\gamma} \mu(E)
  \le \sum_{\substack{E_0\in\mathcal{F} \\ y\in E_0}} \lim_{h\to\infty} S_h(E_0)
  \le \sum_{\substack{E_0\in\mathcal{F} \\ y\in E_0}} \mu(E_0)e^{c_0|E_0|}
  \le \Phi. \qedhere
\]
\end{proof}

\section{Perturbation Estimates for the Symmetric Double-Well Jacobi Matrix}\label{app:doublewell}

In this appendix we prove the spectral and eigenvector localization estimates for the symmetric double-well tridiagonal matrix $\mathcal{L}$ used in Proposition~\ref{prop:nogo}(3).

\begin{lemma}[Double-Well Eigenvalue and Eigenvector Estimates]\label{lem:app-doublewell}
Fix an integer $D\ge 4$ and a barrier height $E_B > 2$. Let $\mathcal{L}$ be the $D\times D$ real symmetric tridiagonal matrix with diagonal entries $a_0 = a_{D-1} = 0$, $a_1 = \dots = a_{D-2} = E_B$, and unit off-diagonal entries $b_0 = \dots = b_{D-2} = 1$, and let $P|j\rangle = |D-1-j\rangle$ be spatial inversion. Then:
\begin{enumerate}
  \setlength{\itemsep}{2pt}\setlength{\parskip}{0pt}
  \item[\emph{(1)}] $\mathcal{L}$ has $D$ simple eigenvalues. Its two lowest eigenvalues $\{E_+, E_-\}$ (indexed by parity $P|\phi_\pm\rangle = \pm|\phi_\pm\rangle$, with sign convention $\langle 0|\phi_\pm\rangle > 0$) are distinct and satisfy
  \begin{equation}\label{eq:dw-evals}
    -\frac{1}{E_B - 2} \le E_\pm \le 0,
  \end{equation}
  while the remaining $D-2$ eigenvalues lie in $[E_B - 2, E_B + 2]$.
  \item[\emph{(2)}] For every $1\le j\le D-2$, setting $d(j) := \min(j, D-1-j)\ge 1$, the components of $|\phi_\pm\rangle$ satisfy
  \begin{equation}\label{eq:dw-decay}
    \bigl|\langle j|\phi_\pm\rangle\bigr| \le \frac{\sqrt{2}}{E_B - 2}\left(\frac{2}{E_B}\right)^{d(j)-1},
  \end{equation}
  and the boundary components satisfy
  \begin{equation}\label{eq:dw-boundary}
    \langle 0|\phi_\pm\rangle = \frac{1}{\sqrt{2}} - O(E_B^{-2}), \qquad
    \langle D-1|\phi_\pm\rangle = \pm\left(\frac{1}{\sqrt{2}} - O(E_B^{-2})\right)
  \end{equation}
  as $E_B\to\infty$ (with $D$ fixed).
\end{enumerate}
\end{lemma}

\begin{proof}
\emph{Part~(1).} Because $\mathcal{L}$ is real symmetric tridiagonal with all off-diagonal entries equal to $1\ne 0$, every eigenvalue of $\mathcal{L}$ is simple (Sturm's theorem). Write $\mathcal{L} = \mathcal{L}_0 + V$ where $\mathcal{L}_0 = \operatorname{diag}(0, E_B, \dots, E_B, 0)$ and $V$ is the adjacency matrix of the path graph on $D$ vertices ($\|V\|_\infty = 2\cos(\frac{\pi}{D+1}) < 2$). By Weyl's eigenvalue inequality, the two lowest eigenvalues $E_1 < E_2$ of $\mathcal{L}$ lie in $(-2, 2)$, whereas the remaining $D-2$ eigenvalues lie in $(E_B-2, E_B+2)$. Because $D\ge 4$, $\mathcal{L}$ vanishes identically on the $2$-dimensional subspace $\mathcal{H}_A := \operatorname{span}\{|0\rangle,|D-1\rangle\}$, so by the Courant--Fischer min-max theorem $E_1 < E_2 \le \max_{\phi\in\mathcal{H}_A,\|\phi\|=1}\langle\phi|\mathcal{L}|\phi\rangle = 0 < E_B - 2$ for every $E_B > 2$.

Decompose $\mathbb{C}^D = \mathcal{H}_A \oplus \mathcal{H}_B$ where $\mathcal{H}_B := \operatorname{span}\{|1\rangle,\dots,|D-2\rangle\}$. In block form,
\[
  \mathcal{L} = \begin{pmatrix} 0 & W^\dagger \\ W & B \end{pmatrix},
\]
where $B = E_B I_{D-2} + V_B$ is the $(D-2)\times(D-2)$ barrier block with $V_B$ the path-graph adjacency matrix on $\{1,\dots,D-2\}$ ($\|V_B\|_\infty < 2$), and $W|0\rangle = |1\rangle$, $W|D-1\rangle = |D-2\rangle$. For any $E < E_B - 2$, the matrix $B - E I_{D-2} \succeq (E_B - 2 - E)I_{D-2} \succ 0$ is strictly positive definite. Eliminating the barrier block $\phi_B = -(B - E I_{D-2})^{-1} W \phi_A$ reduces the eigenvalue equation $\mathcal{L}\phi = E\phi$ to the $2\times 2$ Schur-complement equation
\begin{equation}\label{eq:schur}
  -W^\dagger (B - E I_{D-2})^{-1} W\,\phi_A = E\,\phi_A.
\end{equation}
Because $(B - E I_{D-2})^{-1} \succ 0$, the matrix on the left-hand side of~\eqref{eq:schur} is negative semidefinite, so any eigenvalue $E < E_B - 2$ must satisfy $E \le 0$. Moreover, for $E\le 0$ we have $B - E I_{D-2} \succeq (E_B - 2)I_{D-2}$, and in fact the $2\times 2$ effective Hamiltonian $H_{\mathrm{eff}}(E) := -W^\dagger(B - E I_{D-2})^{-1}W$ has trace and off-diagonal entry given by
\begin{align*}
  \langle 0|H_{\mathrm{eff}}(E)|0\rangle &= \langle D-1|H_{\mathrm{eff}}(E)|D-1\rangle = -\langle 1|(B - E I_{D-2})^{-1}|1\rangle, \\
  \langle 0|H_{\mathrm{eff}}(E)|D-1\rangle &= -\langle 1|(B - E I_{D-2})^{-1}|D-2\rangle.
\end{align*}
Since $B - E I_{D-2} = (E_B - E)(I_{D-2} + \frac{1}{E_B-E}V_B)$ with $\|\frac{1}{E_B-E}V_B\|_\infty < \frac{2}{E_B} < 1$, expanding in a Neumann series gives
\[
  (B - E I_{D-2})^{-1} = \sum_{m=0}^\infty \frac{(-1)^m}{(E_B - E)^{m+1}} V_B^m.
\]
Notice that for $D\ge 4$, $\langle 1|V_B^m|D-2\rangle = 0$ for all $m < D-3$, while for $m = D-3$ there is a unique path of length $D-3$ from $1$ to $D-2$, giving $\langle 1|V_B^{D-3}|D-2\rangle = 1$. Since $P$ commutes with $\mathcal{L}$, the two eigenvectors $|\phi_\pm\rangle$ have definite parities $\pm 1$, with $\phi_{A,\pm} = c_\pm(|0\rangle \pm |D-1\rangle)$. Their eigenvalues solve
\[
  E_\pm = -\langle 1|(B - E_\pm I_{D-2})^{-1}|1\rangle \mp \langle 1|(B - E_\pm I_{D-2})^{-1}|D-2\rangle.
\]
By Sturm simplicity, $E_+ \ne E_-$. Since $E_\pm \le 0$ and $\|V_B\|_\infty < 2$, we have $(B - E_\pm I_{D-2})^{-1} \preceq (E_B - 2)^{-1} I_{D-2}$, whence
\[
  |E_\pm| \le \|H_{\mathrm{eff}}(E_\pm)\|_\infty \le \|(B - E_\pm I_{D-2})^{-1}\|_\infty \le \frac{1}{E_B - 2},
\]
establishing~\eqref{eq:dw-evals}.

\emph{Part~(2).} From $\phi_B = -(B - E_\pm I_{D-2})^{-1} W\phi_A$ and the Neumann series above, for any $1\le j\le D-2$, the shortest path in the barrier $\{1,\dots,D-2\}$ from $j$ to $\{1, D-2\}$ has length $d(j)-1$, where $d(j) = \min(j, D-1-j)$. Hence $\langle j|V_B^m|1\rangle = \langle j|V_B^m|D-2\rangle = 0$ for all $m < d(j)-1$. Using $|\langle 0|\phi_\pm\rangle| + |\langle D-1|\phi_\pm\rangle| \le \sqrt{2}$ (by Cauchy--Schwarz), $\|V_B\|_\infty < 2$, and $E_B - E_\pm \ge E_B$, we obtain
\[
  \bigl|\langle j|\phi_\pm\rangle\bigr|
  \le \sqrt{2}\sum_{m=d(j)-1}^\infty \frac{2^m}{E_B^{m+1}}
   = \frac{\sqrt{2}}{E_B - 2}\left(\frac{2}{E_B}\right)^{d(j)-1},
\]
which proves~\eqref{eq:dw-decay}. Finally, by normalization $1 = |\langle 0|\phi_\pm\rangle|^2 + |\langle D-1|\phi_\pm\rangle|^2 + \sum_{j=1}^{D-2}|\langle j|\phi_\pm\rangle|^2 = 2|\langle 0|\phi_\pm\rangle|^2 + O(E_B^{-2})$, taking square roots with $\langle 0|\phi_\pm\rangle > 0$ and $\langle D-1|\phi_\pm\rangle = \pm\langle 0|\phi_\pm\rangle$ yields~\eqref{eq:dw-boundary}.
\end{proof}

\section*{Use of Artificial Intelligence}

The author used Google's Gemini~4 during the preparation of this manuscript for proof exploration, technical and algebraic refinement, literature organization, numerical verification of finite-size benchmarks, and \LaTeX{} typesetting and stylistic editing.

\begin{thebibliography}{32}

\bibitem{Abanin2015}
D.~A.~Abanin, W.~De~Roeck, and F.~Huveneers,
\textit{Exponentially slow heating in periodically driven many-body systems},
Phys.\ Rev.\ Lett.\ \textbf{115} (2015), 256803.

\bibitem{Alhambra2023}
\'A.~M.~Alhambra,
\textit{Quantum many-body systems in thermal equilibrium},
PRX Quantum \textbf{4} (2023), 040201, arXiv:2204.08349.

\bibitem{Anshu2021}
A.~Anshu, S.~Arunachalam, T.~Kuwahara, and M.~Soleimanifar,
\textit{Sample-efficient learning of interacting quantum systems},
Nature Phys.\ \textbf{17} (2021), 931--935.

\bibitem{Arad2016}
I.~Arad, T.~Kuwahara, and Z.~Landau,
\textit{Connecting global and local energy distributions in quantum spin models on a lattice},
J.\ Stat.\ Mech.\ (2016), 033301.

\bibitem{Araki1969}
H.~Araki,
\textit{Gibbs states of a one dimensional quantum lattice},
Commun.\ Math.\ Phys.\ \textbf{14} (1969), 120--157.

\bibitem{Avdoshkin2020}
A.~Avdoshkin and A.~Dymarsky,
\textit{Euclidean operator growth and quantum chaos},
Phys.\ Rev.\ Research \textbf{2} (2020), 043234.

\bibitem{Baiguera2025}
E.~Rabinovici, A.~S\'anchez-Garrido, R.~Shir, and J.~Sonner,
\textit{Krylov complexity},
arXiv:2507.06286 (2025).

\bibitem{Bakshi2024}
A.~Bakshi, A.~Liu, A.~Moitra, and E.~Tang,
\textit{High-temperature Gibbs states are unentangled and efficiently preparable},
Proc.\ 65th IEEE Symp.\ Found.\ Comput.\ Sci.\ (FOCS) (2024), 1027--1036, arXiv:2403.16850.

\bibitem{Bluhm2025}
A.~Bluhm, \'A.~Capel, and A.~P\'erez-Hern\'andez,
\textit{Exponential decay of mutual information for Gibbs states of local Hamiltonians},
Quantum \textbf{6} (2022), 650.

\bibitem{Bratteli1997}
O.~Bratteli and D.~W.~Robinson,
\textit{Operator Algebras and Quantum Statistical Mechanics 2},
2nd ed., Springer, Berlin, 1997.

\bibitem{Chihara1978}
T.~S.~Chihara,
\textit{An Introduction to Orthogonal Polynomials},
Gordon and Breach, New York, 1978.

\bibitem{Demulder2026}
S.~Demulder, M.~Knysh, and A.~Rolph,
\textit{Krylov exponents and power spectra for maximal quantum chaos: an EFT approach},
J.\ High Energy Phys.\ (2026), arXiv:2508.05444.

\bibitem{Frohlich2015}
J.~Fr\"ohlich and D.~Ueltschi,
\textit{Some properties of correlations of quantum lattice systems in thermal equilibrium},
J.\ Math.\ Phys.\ \textbf{56} (2015), 053302.

\bibitem{Garcia2025}
A.~M.~Garc\'ia-Garc\'ia, C.~Liu, and J.~J.~M.~Verbaarschot,
\textit{Krylov complexity and chaos in deformed Sachdev--Ye--Kitaev models},
Phys.\ Rev.\ B \textbf{111} (2025), 035141, arXiv:2407.09604.

\bibitem{Hastings2006}
M.~B.~Hastings and T.~Koma,
\textit{Spectral gap and exponential decay of correlations},
Commun.\ Math.\ Phys.\ \textbf{265} (2006), 781--804.

\bibitem{Kato2019}
K.~Kato and F.~G.~S.~L.~Brand\~ao,
\textit{Quantum approximate Markov chains are thermal},
Commun.\ Math.\ Phys.\ \textbf{370} (2019), 117--149.

\bibitem{Kim2026}
A.~D.~Kim, T.~H.~Yang, and J.~Y.~Lee,
\textit{High-temperature Gibbs states are uniformly Markovian and trivial},
arXiv:2611.26626 (2026).

\bibitem{Kliesch2014}
M.~Kliesch, C.~Gogolin, M.~J.~Kastoryano, A.~Riera, and J.~Eisert,
\textit{Locality of temperature},
Phys.\ Rev.\ X \textbf{4} (2014), 031019.

\bibitem{Kotecky1986}
R.~Koteck\'y and D.~Preiss,
\textit{Cluster expansion for abstract polymer models},
Commun.\ Math.\ Phys.\ \textbf{103} (1986), 491--498.

\bibitem{Kuwahara2020a}
T.~Kuwahara and K.~Saito,
\textit{Gaussian concentration bound and ensemble equivalence in generic quantum many-body systems including long-range interactions},
Ann.\ Phys.\ \textbf{421} (2020), 168278, arXiv:1906.10872.

\bibitem{Kuwahara2020b}
T.~Kuwahara, K.~Kato, and F.~G.~S.~L.~Brand\~ao,
\textit{Clustering of conditional mutual information for quantum Gibbs states above a threshold temperature},
Phys.\ Rev.\ Lett.\ \textbf{124} (2020), 220601; Erratum, Phys.\ Rev.\ Lett.\ \textbf{134} (2025), 199901.

\bibitem{Lieb1972}
E.~H.~Lieb and D.~W.~Robinson,
\textit{The finite group velocity of quantum spin systems},
Commun.\ Math.\ Phys.\ \textbf{28} (1972), 251--257.

\bibitem{MSS2016}
J.~Maldacena, S.~H.~Shenker, and D.~Stanford,
\textit{A bound on chaos},
J.\ High Energy Phys.\ \textbf{08} (2016), 106, arXiv:1503.01409.

\bibitem{MS2016}
J.~Maldacena and D.~Stanford,
\textit{Remarks on the Sachdev--Ye--Kitaev model},
Phys.\ Rev.\ D \textbf{94} (2016), 106002, arXiv:1604.07818.

\bibitem{Murthy2019}
C.~Murthy and M.~Srednicki,
\textit{Bounds on chaos from the eigenstate thermalization hypothesis},
Phys.\ Rev.\ Lett.\ \textbf{123} (2019), 230606.

\bibitem{Nachtergaele2006}
B.~Nachtergaele and R.~Sims,
\textit{Lieb--Robinson bounds and the exponential clustering theorem},
Commun.\ Math.\ Phys.\ \textbf{265} (2006), 119--130.

\bibitem{Parker2019}
D.~E.~Parker, X.~Cao, A.~Avdoshkin, T.~Scaffidi, and E.~Altman,
\textit{A universal operator growth hypothesis},
Phys.\ Rev.\ X \textbf{9} (2019), 041017, arXiv:1812.08657.

\bibitem{Perugu2026}
R.~Perugu et al.,
\textit{Krylov winding and emergent coherence in operator growth dynamics},
Phys.\ Rev.\ Lett.\ \textbf{136}(15) (2026), arXiv:2509.25331.

\bibitem{Qi2019}
X.-L.~Qi and A.~Streicher,
\textit{Quantum epidemiology: operator growth, thermal effects, and SYK},
J.\ High Energy Phys.\ \textbf{08} (2019), 012.

\bibitem{Roberts2018}
D.~A.~Roberts, D.~Stanford, and A.~Streicher,
\textit{Operator growth in the SYK model},
J.\ High Energy Phys.\ \textbf{06} (2018), 122.

\bibitem{Rouze2026}
C.~Rouz\'e, D.~Stilck~Fran\c{c}a, and \'A.~M.~Alhambra,
\textit{Optimal quantum algorithm for Gibbs state preparation},
Phys.\ Rev.\ Lett.\ \textbf{136} (2026), 060601, arXiv:2411.04885.

\bibitem{Ueltschi2004}
D.~Ueltschi,
\textit{Cluster expansions and correlation functions},
Moscow Math.\ J.\ \textbf{4} (2004), 511--522.

\bibitem{Yin2020}
C.~Yin and A.~Lucas,
\textit{Bound on quantum scrambling with all-to-all interactions},
Phys.\ Rev.\ A \textbf{102} (2020), 022402.

\end{thebibliography}

\vspace{16pt}
\noindent\textsc{Prince Anand}\\
\emph{E-mail address:} \texttt{princeanand123pc@gmail.com}

\end{document}
